\section{Bench nopCommerce --- publiczny okaz}\label{sec:bench}

Repozytorium \repopath{letsgolegacy.secondkey-nopcommerce-bench} jest publicznym okazem
(specimen) frameworku letsgolegacy, w~praktyce nazywanym „bench nopCommerce”. Pokazuje
łańcuch Second Key na prawdziwym systemie legacy: nopCommerce 3.x, sklepie na klasycznym
ASP.NET MVC i~.NET Framework, który migruje agent migracyjny Microsoftu (modernize-dotnet),
a~Second Key sprawdza -- niezależnie i~deterministycznie -- względem spisanego kontraktu
tego, co sklep musi robić i~czego nie może robić nigdy. Second Key niczego nie migruje;
tylko sprawdza.

Rozdział opisuje migawkę repozytorium w~stanie \texttt{origin/main} (commit
\texttt{e6d4fe9}). W~tym stanie zadania P1--P5 są zrobione (build, host legacy, rozgrzewka
łańcucha, nagranie ruchu, kontrakt z~dowodem A/A), a~P6--P10 są zaplanowane. README
zaznacza, że repozytorium jest w~budowie i~że żaden evidence pack (pakiet dowodowy) nie
został jeszcze opublikowany; pakiety powstające w~każdym przebiegu CI są artefaktami
workflowów, a~nie publikacją. Kolejne podsekcje idą od decyzji (ADR 0001--0006) przez
nagranie ruchu, kontrakt, skrypty i~workflowy do wyników P3--P5, runbooka P6 i~planów
P7--P10.

\zrodla{\path{letsgolegacy.secondkey-nopcommerce-bench/README.md},
\path{letsgolegacy.secondkey-nopcommerce-bench/docs/WORKPLAN.md}}

\subsection{Cel okazu i~zawartość repozytorium}\label{sec:bench-cel}

Plan pracy (\repopath{docs/WORKPLAN.md}) nazywa okaz publicznym dowodem, że łańcuch
Second Key działa na prawdziwym kodzie legacy. Wybranym systemem jest starszy nopCommerce
(3.x, klasyczne ASP.NET MVC na .NET Framework): prawdziwy sklep z~rabatami, podatkami
i~wysyłką, w~którym reguły „never” coś znaczą. Migrację wykonuje agent modernize-dotnet;
Second Key tylko ją sprawdza.

Pokaz jest zbudowany wokół jednego momentu: agent migruje, build i~istniejące testy
przechodzą, a~Second Key mimo to znajduje prawdziwą różnicę w~zachowaniu. Najlepszym
kandydatem na taką różnicę jest globalizacja: przejście z~.NET Framework na nowoczesny
.NET przełącza ją z~NLS na ICU, więc porównywanie i~sortowanie napisów zależne od kultury
może się zmienić -- plan nazywa to znaną, udokumentowaną pułapką migracji. Miejsca
w~nopCommerce 3.90, w~których ta różnica może wyjść na wierzch, i~scenariusze, które je
sondują, opisuje sekcja~\ref{sec:bench-plany}.

\finding{Kod nopCommerce nie jest częścią repozytorium: CI pobiera przypięte wydanie
z~projektu upstream w~czasie buildu, a~repozytorium zawiera wyłącznie bench -- skrypty,
ruch, kontrakt i~wyniki.} Bench ma dzięki temu własne warunki (README: wszystkie prawa
zastrzeżone), a~nopCommerce zostaje na licencji swoich autorów.

Identyfikatory zadań (P1--P10) są zgodne z~backlogiem wspólnym dla wszystkich
repozytoriów frameworku. Jedno zadanie to jeden pull request, chyba że jest oznaczone jako praca człowieka.
Zadania należą do fazy 01 (demo), a~P9 tego repozytorium jest bramką fazy 01; status
zadań podaje sekcja~\ref{sec:bench-plany}. Tabela~\ref{tab:bench-repo} pokazuje, co
zawiera repozytorium.

\begin{table}[htbp]
\centering
\small
\caption{Zawartość repozytorium bench nopCommerce}\label{tab:bench-repo}
\begin{tabularx}{\linewidth}{@{}l X@{}}
\toprule
Ścieżka & Zawartość \\
\midrule
\path{pins.json} & jedyne źródło prawdy o~wszystkim, co bench pobiera lub buduje
  (sekcja~\ref{sec:bench-pins}) \\
\path{scripts/} & kroki 1--9 strony legacy, \path{collect-diagnostics.ps1}, biblioteki
  \path{lib/}, scenariusze ruchu \path{traffic/}, skrypty rozgrzewki \path{eshop/}
  (sekcja~\ref{sec:bench-skrypty}) \\
\path{contract/} & kontrakt \path{contract.yaml}, konfiguracja odtwarzania
  \path{secondkey.yaml}, README z~mapą reguł planu na klauzule
  (sekcja~\ref{sec:bench-kontrakt}) \\
\path{warmup/eshop/} & kontrakt i~konfiguracja rozgrzewki P3
  (sekcja~\ref{sec:bench-rozgrzewka}) \\
\path{docs/} & \path{WORKPLAN.md}, \path{TOPOLOGY.md}, \path{P4-TRAFFIC-PLAN.md},
  \path{P6-RUNBOOK.md}, \path{P10-ONBOARDING-HOURS.md}, \path{adr/} (ADR 0001--0006) \\
\path{results/p3/}, \path{p4/}, \path{p5/} & to, co zapisały przebiegi CI zamykające
  zadania P3--P5 (sekcja~\ref{sec:bench-wyniki}) \\
\path{.github/} & workflowy \path{legacy-build.yml}, \path{chain-tools.yml},
  \path{warmup-eshop.yml}, \path{secret-scan.yml}; \path{CODEOWNERS};
  \path{dependabot.yml} (sekcja~\ref{sec:bench-ci}) \\
\path{.githooks/}, \path{.gitleaks.toml} & skan sekretów przed commitem i~w~CI
  (sekcja~\ref{sec:bench-higiena}) \\
\path{PSScriptAnalyzerSettings.psd1} & ustawienia lintu skryptów \\
\path{.gitignore} & m.in.\ zakaz commitowania nagrań \path{*.skcap} i~\path{*.skrun} \\
\bottomrule
\end{tabularx}
\end{table}

\zrodla{\path{letsgolegacy.secondkey-nopcommerce-bench/README.md},
\path{letsgolegacy.secondkey-nopcommerce-bench/docs/WORKPLAN.md},
\path{letsgolegacy.secondkey-nopcommerce-bench/.gitignore}}

\subsection{Wybór i~przypięcie nopCommerce 3.90 (ADR 0001)}\label{sec:bench-pin}

\paragraph{Kontekst.} Bench potrzebuje prawdziwego sklepu legacy: klasycznego ASP.NET na
.NET Framework, SQL Server i~reguł biznesowych, które mają znaczenie (rabaty, podatki,
wysyłka), tak by reguły „must” i~„never” kontraktu P5 coś znaczyły. Plan pracy ustalił
produkt (nopCommerce) i~wersję główną (3.x); konkretne wydanie było otwarte. Upstream
taguje każde wydanie 3.x: \texttt{release-3.00} … \texttt{release-3.80},
\texttt{release-3.90} oraz dwie bety. Następna wersja główna, 4.00, przeszła na ASP.NET
Core, co przekreśla cel: obserwowana ma być migracja z~System.Web na nowoczesny .NET.

\paragraph{Decyzja.} Przypięte jest wydanie \texttt{release-3.90}:

\begin{center}
\small
\begin{tabularx}{\linewidth}{@{}l X@{}}
\toprule
Tag & \texttt{release-3.90} (obiekt tagu z~adnotacją
  \texttt{5846b2b59b102270e15d7eb8415298ae989a1abf}, otagowany 2017-03-15) \\
Commit & \texttt{12d1f0139149a9d6e84964d325d62bae6ac748f6} \\
Stos & ASP.NET MVC 5.2.3, Entity Framework 6.1.3, Autofac 4, .NET Framework 4.5.1,
  SQL Server (lub SQL CE) \\
Rozmiar & 5397 śledzonych plików; 94~MiB spakowane, 338~MB po checkoucie (upstream
  commituje swój katalog \texttt{packages/}) \\
\bottomrule
\end{tabularx}
\end{center}

Przypięcie jest w~jednym miejscu, \repopath{pins.json}. Krok~1
(\repopath{1-fetch-nopcommerce.ps1}) egzekwuje je przy każdym uruchomieniu dwa razy:
tag upstream musi nadal wskazywać przypięty commit (przesunięty tag przerywa przebieg,
zamiast po cichu zbudować coś innego), a~drzewo po checkoucie musi być dokładnie tym
commitem, bez zmodyfikowanego pliku śledzonego. Krok~2 powtarza drugi test przed buildem
i~po nim, więc opublikowana witryna jest dowodliwie buildem niezmienionego kodu upstream.

\paragraph{Dlaczego 3.90, a~nie wcześniejsze 3.x.}
\begin{itemize}
  \item To \textbf{ostatnie} wydanie na klasycznym ASP.NET MVC: zawiera wszystkie
    poprawki linii 3.x i~najpełniejszą powierzchnię biznesową -- reguły wymagań rabatów,
    dostawców podatku według kraju, stanu i~kodu pocztowego, kilku dostawców stawek
    wysyłki, wiele sklepów, dostawców (vendors), karty podarunkowe, punkty lojalnościowe,
    atrybuty zamówienia. Więcej reguł do zapisania w~kontrakcie i~więcej sposobów na
    regresję.
  \item Jego toolchain (projekty Visual Studio 2015, \texttt{packages.config}) jest
    najmłodszy w~linii 3.x, co daje najkrótszą drogę do buildu na obecnych obrazach
    runnerów.
  \item Żadne wcześniejsze 3.x nie wzmocniłoby pokazu; dodałoby tylko lat.
\end{itemize}

\paragraph{Dlaczego pobierany, a~nie dołączony do repozytorium.}
\begin{itemize}
  \item \textbf{Licencja.} nopCommerce 3.90 jest na nopCommerce Public License 3.0 (GPLv3
    z~wymogiem atrybucji „powered by nopCommerce”), a~repozytorium bench ma wszystkie prawa
    zastrzeżone. Trzymanie kodu upstream poza drzewem rozdziela oba zestawy warunków.
  \item \textbf{Rozmiar i~szum.} 338~MB i~około 5000 plików cudzego kodu przytłoczyłyby
    własną historię benchu i~każdy diff w~review.
\end{itemize}

\paragraph{Konsekwencje.}
\begin{itemize}
  \item CI potrzebuje w~czasie buildu dostępu do github.com; pobieranie ponawia błędy
    przejściowe i~na hostowanym runnerze trwa sekundy.
  \item Artefakt workflowu \repopath{legacy-site} jest buildem niezmienionego kodu upstream.
    Jego manifest (\path{state/build.json}) podaje repozytorium i~commit upstream, czyli
    miejsce publikacji źródeł. Stopka „powered by nopCommerce” pozostaje nietknięta.
    Artefakt jest przechowywany 7~dni -- wystarczająco dla jobów, które go konsumują.
  \item Zmiana przypięcia jest decyzją, a~nie rutynową czynnością: wymaga nowego ADR
    i~ponownego przebiegu smoke testu oraz zestawu ruchu (P4) na nowym wydaniu.
  \item Dane przykładowe 3.90 są datowane: rabat „20\% sumy zamówienia” wygasł
    2020-01-01. To zachowanie do nagrania i~zapisania w~kontrakcie (P4, P5), a~nie defekt
    do łatania.
\end{itemize}

\zrodla{\path{letsgolegacy.secondkey-nopcommerce-bench/docs/adr/0001-pin-nopcommerce-3.90.md},
\path{letsgolegacy.secondkey-nopcommerce-bench/README.md},
\path{letsgolegacy.secondkey-nopcommerce-bench/scripts/1-fetch-nopcommerce.ps1},
\path{letsgolegacy.secondkey-nopcommerce-bench/scripts/2-build-legacy.ps1}}

\subsubsection{Plik \texttt{pins.json}}\label{sec:bench-pins}

\repopath{pins.json} jest jedynym źródłem prawdy o~wszystkim, co bench pobiera lub buduje:
o~stronie legacy, o~narzędziach łańcucha Second Key i~o~aplikacji rozgrzewki. Każde
pobranie jest weryfikowane względem tych wartości, zanim zostanie użyte, a~zmiana wartości
jest decyzją, której powód zapisuje się w~\path{docs/adr/}. Plik ma cztery bloki:

\begin{itemize}
  \item \texttt{nopCommerce} -- repozytorium, tag, commit, ścieżki rozwiązania, projektu
    \texttt{Nop.Web} i~katalogu wtyczek oraz docelowa wersja frameworku; czytają go kroki~1
    i~2.
  \item \texttt{toolchain} -- \texttt{NuGet.CommandLine} 7.9.0 (\texttt{nuget.exe}, gdy
    żadnego nie ma na \texttt{PATH}), pakiet asemblacji referencyjnych .NET Framework 4.5.1
    (ADR 0002) oraz nośnik SQL Server 2022 Express z~oczekiwanym podpisującym,
    nazwą instancji i~kolacją (ADR 0003); czytają go kroki 2, 3 i~4. Każdy plik ma adres
    URL i~skrót SHA-256.
  \item \texttt{eShopModernizing} -- wyłącznie rozgrzewka P3 (ADR 0004): repozytorium,
    commit, projekt i~katalog źródeł eShopLegacyMVC oraz dwa pakiety asemblacji
    referencyjnych (4.6.1 i~4.7.2).
  \item \texttt{chain} -- commity \texttt{letsgolegacy.secondkey} i~\texttt{letsgolegacy.portcullis},
    z~których workflow \path{chain-tools.yml} buduje \repopath{sk} i~Portcullis za pomocą
    .NET 10 SDK.
\end{itemize}

Klucze pamięci podręcznej CI zawierają skrót \repopath{pins.json}, więc każda zmiana
pliku unieważnia zbuforowane pobrania. Przypięte wartości (adresy URL i~skróty SHA-256 pobrań są
w~pliku):

\begin{center}
\small
\begin{tabularx}{\linewidth}{@{}>{\raggedright\arraybackslash}p{4.1cm} >{\raggedright\arraybackslash}X@{}}
\toprule
Klucz & Wartość \\
\midrule
\texttt{nopCommerce} & \path{https://github.com/nopSolutions/nopCommerce.git}, tag
  \texttt{release-3.90}, commit \texttt{12d1f0139149a9d6e84964d325d62bae6ac748f6};
  rozwiązanie \path{src/NopCommerce.sln}, projekt
  \path{src/Presentation/Nop.Web/Nop.Web.csproj}, wtyczki \path{src/Plugins}, framework
  \texttt{v4.5.1} \\
\path{toolchain.nuget} & \texttt{NuGet.CommandLine} 7.9.0 (nuget.org) \\
\path{toolchain.referenceAssemblies} &
  \path{Microsoft.NETFramework.ReferenceAssemblies.net451} 1.0.3 (nuget.org) \\
\path{toolchain.sqlServerExpress} & „SQL Server 2022 Express (core media, x64,
  English)”, plik \path{SQLEXPR_x64_ENU.exe} z~download.microsoft.com, podpisujący
  \texttt{Microsoft Corporation}, instancja \texttt{SQLEXPRESS}, kolacja
  \texttt{SQL\_Latin1\_General\_CP1\_CI\_AS} \\
\texttt{eShopModernizing} & \path{https://github.com/dotnet-architecture/eShopModernizing.git},
  commit \texttt{63bc9ec4414d7e281dc9f9d7bdcf70030950457d}; projekt
  \path{eShopLegacyMVCSolution/src/eShopLegacyMVC/eShopLegacyMVC.csproj}, źródła
  \path{eShopLegacyMVCSolution/src/eShopLegacyMVC}; pakiety
  \path{Microsoft.NETFramework.ReferenceAssemblies.net461} i~\path{...net472} w~wersji
  1.0.3 (frameworki \texttt{v4.6.1}, \texttt{v4.7.2}) \\
\path{chain.secondKey} & \path{konradcinkusz/letsgolegacy.secondkey}, commit
  \texttt{967c49c0005c4412180914f43437d75b43f148dc}, projekt
  \path{src/SecondKey.Cli/SecondKey.Cli.csproj}, asemblacja \path{SecondKey.Cli.dll} \\
\path{chain.portcullis} & \path{konradcinkusz/letsgolegacy.portcullis}, commit
  \texttt{84e1925fe0d85cf23415f0d33c98590f86128722}, projekt
  \path{src/Portcullis.Cli/Portcullis.Cli.csproj}, asemblacja \path{Portcullis.Cli.dll} \\
\bottomrule
\end{tabularx}
\end{center}

\zrodla{\path{letsgolegacy.secondkey-nopcommerce-bench/pins.json},
\path{letsgolegacy.secondkey-nopcommerce-bench/.github/workflows/legacy-build.yml},
\path{letsgolegacy.secondkey-nopcommerce-bench/.github/workflows/chain-tools.yml}}

\subsection{Topologia demonstracyjna (decyzja D7)}\label{sec:bench-topologia}

\repopath{docs/TOPOLOGY.md} zapisuje decyzję D7 ze statusem „proposed”. Strona legacy jest
zbudowana i~działa w~CI (P1, P2); części oznaczone jako planowane opisują, jak miały się
podłączyć zadania P4--P7, które mogą je doprecyzować i~powinny wtedy zaktualizować tę
stronę.

\paragraph{Decyzja: jeden host Windows obsługuje stronę legacy, Linux -- resztę łańcucha.}
\begin{itemize}
  \item Sklep legacy potrzebuje Windows: nopCommerce 3.90 to ASP.NET MVC 5 na .NET
    Framework, serwowany przez IIS, z~SQL Server pod spodem. Jeden host Windows mieści
    wszystkie trzy.
  \item Łańcuch Second Key -- capture, replay, compare, evidence -- jest aplikacją .NET 10
    i~działa wszędzie; jego naturalnym miejscem jest Linux (szybszy start runnerów,
    kontenery).
  \item Komponenty wymieniają \textbf{pliki}, nigdy swoje wnętrza: \path{*.skcap},
    \path{contract.yaml}, \path{*.skrun}, \path{verdict.json}, \path{*.sarif}
    (sekcja~\ref{sec:formaty}). Granica między maszynami może więc przebiegać wszędzie
    tam, dokąd da się skopiować te pliki.
  \item Tylko dwa kroki muszą sięgać do działającego sklepu przez HTTP: \textbf{capture}
    (P4) i~\textbf{replay} (A/A w~P5, kandydat w~P7). Działają tam, skąd sklep jest
    osiągalny; wszystko dalej to „plik na wejściu, plik na wyjściu”.
\end{itemize}

Tabela~\ref{tab:bench-porty} zbiera nazwy, porty i~miejsca. \texttt{<work>} oznacza
katalog roboczy skryptów (sekcja~\ref{sec:bench-stan}).

{\footnotesize
\begin{xltabular}{\linewidth}{@{}>{\raggedright\arraybackslash}p{3.3cm} >{\raggedright\arraybackslash}X >{\raggedright\arraybackslash}p{3.6cm}@{}}
\caption{Nazwy, porty i~miejsca (TOPOLOGY)}\label{tab:bench-porty}\\
\toprule
Co & Wartość & Ustawia \\
\midrule
\endfirsthead
\toprule
Co & Wartość & Ustawia \\
\midrule
\endhead
\bottomrule
\endlastfoot
Sklep legacy & \url{http://localhost:8080/} & krok~4, \texttt{-Port} \\
Witryna i~pula aplikacji IIS & \texttt{nopcommerce-legacy}; .NET CLR v4.0, potok
  zintegrowany, bez limitu bezczynności, bez okresowego recyklingu & krok~4,
  \texttt{-SiteName} \\
Tożsamość witryny & \texttt{IIS APPPOOL\textbackslash{}nopcommerce-legacy}: Modify na
  katalogu witryny, \texttt{dbcreator} w~SQL Server, nic więcej & krok~4 \\
Katalog witryny & \path{C:\inetpub\nopcommerce-legacy} & krok~4, \texttt{-InstallDir} \\
SQL Server & SQL Server 2022 Express, instancja \texttt{.\textbackslash{}SQLEXPRESS},
  kolacja serwera \texttt{SQL\_Latin1\_General\_CP1\_CI\_AS}; połączenia lokalne przez
  pamięć współdzieloną & krok~3, \path{pins.json} \\
Baza danych & \texttt{nopcommerce\_legacy}, tworzona przez instalator nopCommerce
  z~przypiętą kolacją, danymi przykładowymi i~uwierzytelnianiem Windows (nigdzie hasła) &
  krok~4, \texttt{-DatabaseName} \\
Administrator sklepu & \texttt{admin@nopbench.invalid}, hasło generowane na przebieg,
  przechowywane w~\path{<work>\secrets\legacy-admin.json} & krok~4 \\
Proxy nagrywające (P4) & \url{http://127.0.0.1:8000/} $\rightarrow$
  \url{http://localhost:8080/} & krok~7, \texttt{-ProxyUrl} \\
Snapshot bazy (P4) & \texttt{nopcommerce\_legacy\_snapshot}, robiony po konfiguracji
  sklepu, przywracany przed każdym scenariuszem & krok~7 \\
Kandydat (planowany, P7) & \url{http://localhost:8090/}, własna kopia bazy & P7 \\
Rozgrzewka P3: eShopLegacyMVC & \url{http://localhost:8081/}, witryna IIS
  \texttt{eshop-legacy}, dane mock (bez bazy) & \path{scripts/eshop/2-deploy-eshop.ps1} \\
Rozgrzewka P3: proxy nagrywające & \url{http://127.0.0.1:8001/} $\rightarrow$
  \url{http://localhost:8081/} & \path{scripts/eshop/3-record-eshop.ps1} \\
\end{xltabular}
}

\paragraph{W~CI.} Runnery hostowane przez GitHub nie dzielą sieci: job na runnerze linuksowym nie dosięgnie
IIS na runnerze Windows. Dlatego capture i~replay działają \textbf{wewnątrz joba Windows},
obok IIS -- \repopath{sk} jest wieloplatformową aplikacją .NET 10, więc uruchomienie go na
Windows niczego nie komplikuje. To, co produkują (\path{*.skcap}, \path{*.skrun}),
opuszcza job jako artefakt, a~kroki „plik na wejściu, plik na wyjściu” (sprawdzenia
kontraktu, compare, evidence) działają w~jobie linuksowym. Rozgrzewka P3 była pierwszym
przebiegiem tego podziału (ADR 0004); rozkład jobów pokazuje
sekcja~\ref{sec:bench-ci}. Zmierzone na \texttt{windows-2022}: job buildu trwa
około 1,5~minuty; job hosta dodaje IIS, SQL Server Express i~instalator (czasy są
w~podsumowaniu każdego przebiegu oraz w~\path{state\platform.json}
i~\path{state\deploy.json}).

\paragraph{Na stacji roboczej.} Te same skrypty, w~tej samej kolejności, na jednym z:
\begin{itemize}
  \item \textbf{maszynie wirtualnej Windows Server 2022 pod Hyper-V} (do pokazu wystarczy
    nośnik ewaluacyjny; 4~vCPU i~8~GB są wygodne) albo
  \item \textbf{komputerze z~Windows 10/11 Pro lub Enterprise z~IIS} -- skrypty włączają
    IIS przez funkcje opcjonalne DISM, których nazwy są takie same w~Windows klienckim
    i~serwerowym.
\end{itemize}
Kroki 1--5 uruchamia się w~PowerShell z~podniesionymi uprawnieniami (kroki~1 i~2 ich nie
potrzebują). Żeby strona linuksowa (WSL2 albo maszyna Linux w~tym samym przełączniku
wirtualnym) dosięgła sklepu, trzeba otworzyć port na hoście Windows -- CI tego nie
potrzebuje, więc skrypty tego nie robią:

\begin{lstlisting}
New-NetFirewallRule -DisplayName 'nopbench legacy 8080' -Direction Inbound -Protocol TCP -LocalPort 8080 -Action Allow
\end{lstlisting}

Następnie z~Linuksa (smoke test działa także na PowerShell 7 w~Linuksie):

\begin{lstlisting}
pwsh scripts/5-smoke.ps1 -BaseUrl http://<windows-host>:8080/
\end{lstlisting}

\paragraph{Co host dokłada do zachowania -- i~dlatego jest ustalone.}

\begin{center}
\small
\begin{tabularx}{\linewidth}{@{}>{\raggedright\arraybackslash}p{2.6cm} >{\raggedright\arraybackslash}X >{\raggedright\arraybackslash}X@{}}
\toprule
Ustawienie & Wartość & Dlaczego ma znaczenie \\
\midrule
Kolacja SQL & \texttt{SQL\_Latin1\_General\_CP1\_CI\_AS}, przypięta dla serwera i~bazy &
  kolejność i~równość w~zapytaniach (sortowanie katalogu, wyszukiwanie) \\
Środowisko .NET Framework & 4.8 (hosta), uruchamiające aplikację zbudowaną dla 4.5.1 &
  środowisko zastępowane w~miejscu: tak działa każdy produkcyjny host tego sklepu \\
Kultura & \texttt{en-US}, ustawiana przez nopCommerce na każde żądanie z~języka sklepu &
  format cen i~dat; porównanie NLS $\rightarrow$ ICU w~P7 dotyczy środowiska, więc obie
  strony muszą widzieć tę samą kulturę \\
Strefa czasowa & UTC na runnerach hostowanych & ważność rabatów, daty zamówień, produkty
  „nowe” \\
Kompresja & niezainstalowana & nagrane treści zostają czytelne; kompresja należy do
  transportu, nie do kontraktu \\
Zegar & czas rzeczywisty & nie da się go kontrolować bez ruszania kodu legacy; zamiast
  tego komparator normalizuje czasy (zadanie rdzenia S12) \\
\bottomrule
\end{tabularx}
\end{center}

\zrodla{\path{letsgolegacy.secondkey-nopcommerce-bench/docs/TOPOLOGY.md},
\path{letsgolegacy.secondkey-nopcommerce-bench/docs/adr/0004-chain-tools-and-eshop-warmup.md}}

\subsection{Katalog roboczy, pliki przekazania i~jedyny sekret}\label{sec:bench-stan}

\paragraph{Katalog roboczy \texttt{<work>}.} Każdy skrypt rozwiązuje katalog roboczy
w~tej kolejności: parametr \texttt{-WorkRoot}, zmienna \texttt{NOPBENCH\_WORK}, a~potem
domyślnie \path{C:\nopbench} (dokładniej \texttt{\$env:SystemDrive\textbackslash{}nopbench})
w~Windows i~\path{$HOME/nopbench} w~innych systemach. W~CI workflowy ustawiają
\path{NOPBENCH_WORK=$env:RUNNER_TEMP\nopbench}, co na hostowanym
runnerze Windows daje \path{D:\a\_temp\nopbench}. Ścieżka ma być krótka: potok publikacji
webowej MSBuild głęboko zagnieżdża ścieżki, a~starsze zadania MSBuild nadal trafiają na
limit 260 znaków. Nic, co zapisuje którykolwiek krok, nie ląduje w~repozytorium.

Katalog roboczy zawiera: \path{legacy-src} (checkout, krok~1), \path{legacy-site}
(publikacja, krok~2), \path{downloads} (zweryfikowane pobrania: pakiety
\path{<id>.<wersja>.nupkg} i~\path{SQLEXPR_x64_ENU.exe}; w~CI buforowany),
\path{tools} (rozpakowane pakiety ze znacznikiem \path{.sha256}, nośnik
\path{sqlexpress-media}, narzędzia łańcucha w~\path{chain\secondkey}
i~\path{chain\portcullis}, asemblacje rozgrzewki w~\path{eshop-reference-assemblies}),
\path{logs} (binlogi MSBuild, logi proxy \path{traffic-capture.log} i~\path{.log.err}),
\path{state} (pliki przekazania), \path{secrets\legacy-admin.json},
\path{sqldata} (tylko gdy dysk systemowy zgłasza sektory atomowe ponad 4096~B),
\path{traffic\nopcommerce.skcap}, \path{replay} (\path{run.skrun}, \path{verdict.json},
\path{evidence}), \path{diagnostics} oraz katalogi rozgrzewki \path{eshop-src},
\path{eshop-site} i~\path{eshop}. Poza nim są tylko katalogi serwowane przez IIS
(\path{C:\inetpub\nopcommerce-legacy}, \path{C:\inetpub\eshop-legacy}) i~pliki snapshotu
bazy (\path{<snapshot>_<plik>.ss}) obok plików danych bazy.

\paragraph{Pliki przekazania.} Każdy krok zapisuje, co wyprodukował, w~pliku
\path{<work>\state\<krok>.json}, a~następny krok czyta go, gdy nie dostanie parametru:
najpierw jawny parametr, potem plik przekazania, na końcu wartość domyślna. Funkcja
\texttt{Write-BenchState} zapisuje JSON w~UTF-8 bez BOM i~z~zasady nie zapisuje sekretów.

{\footnotesize
\begin{xltabular}{\linewidth}{@{}>{\raggedright\arraybackslash}p{3.4cm} l >{\raggedright\arraybackslash}X@{}}
\caption{Pliki przekazania w~\texttt{<work>\textbackslash{}state}}\label{tab:bench-state}\\
\toprule
Plik & Zapisuje & Zawiera \\
\midrule
\endfirsthead
\toprule
Plik & Zapisuje & Zawiera \\
\midrule
\endhead
\bottomrule
\endlastfoot
\path{fetch.json} & krok~1 & tag, commit, ścieżkę checkoutu, liczbę plików śledzonych \\
\path{build.json} & krok~2 & zbudowany commit, wersje narzędzi i~pakietu asemblacji,
  ścieżkę witryny, liczbę plików i~wtyczek, skrót manifestu, obraz hosta, czasy faz \\
\path{legacy-site.sha256} & krok~2 & SHA-256 każdego opublikowanego pliku, zgodny
  z~\texttt{sha256sum -c}, sortowany porządkowo \\
\path{platform.json} & krok~3 & stan IIS, instancję, wersję i~kolację SQL Server \\
\path{deploy.json} & krok~4 & URL, witrynę, tożsamość, bazę i~jej kolację, connection
  string (bez sekretu), liczbę tabel i~produktów, ścieżkę pliku z~danymi administratora,
  czasy \\
\path{smoke.json} & krok~5 & każdą kontrolę ze statusem, czasem i~wynikiem \\
\path{store-config.json} & krok~6 & co skonfigurowano, z~nadanymi identyfikatorami
  (kategoria podatku, kraje, produkty, rabaty, rola klienta) \\
\path{traffic.json} & krok~7 & ścieżkę i~skrót nagrania, scenariusze z~identyfikatorami
  sesji i~liczbą żądań, snapshot, odpowiedzi na pierwsze żądanie po każdym przywróceniu,
  skrót manifestu witryny, znacznik \texttt{partial} \\
\path{replay.json} & krok~8 & ścieżkę i~skrót przebiegu, tryb (A/A albo legacy vs
  kandydat), oba URL-e, liczby scenariuszy, wyników i~resetów, statusy \\
\path{verdict.json} & krok~9 & wynik, podsumowanie, klauzule niespełnione, maski, pliki
  pakietu \\
\path{eshop-*.json} & rozgrzewka & odpowiedniki dla eShopLegacyMVC (\texttt{build},
  \texttt{deploy}, \texttt{traffic}, \texttt{replay}, \texttt{scan} -- SARIF, liczba
  znalezisk i~kod wyjścia Portcullis --, \texttt{verdict}) \\
\end{xltabular}
}

\paragraph{Jedyny sekret.} Krok~4 tworzy administratora sklepu z~hasłem generowanym na
dany przebieg: 32~znaki losowane generatorem \texttt{RandomNumberGenerator}
z~57-znakowego alfabetu (bez liter I, O, l i~cyfr 0, 1), chyba że ustawiono
\texttt{NOPBENCH\_ADMIN\_PASSWORD}. Hasło jest rejestrowane w~maskowaniu logów GitHub
Actions (\texttt{::add-mask::}), nigdy nie jest drukowane i~trafia wyłącznie do
\path{<work>\secrets\legacy-admin.json} (pola \texttt{url}, \texttt{email},
\texttt{password}, \texttt{generated}). ACL pliku jest ustawiana przez \texttt{icacls}
z~wyłączonym dziedziczeniem: pełny dostęp mają bieżący użytkownik, grupa Administratorzy
i~SYSTEM. Krok~6 czyta plik i~ponownie maskuje hasło; diagnostyka nigdy nie zbiera
katalogu \path{secrets}. Baza danych nie ma hasła w~ogóle: witryna łączy się tożsamością
swojej puli aplikacji (uwierzytelnianie Windows), a~adres \texttt{admin@nopbench.invalid}
należy do zarezerwowanej domeny, więc nic nie jest nigdy dostarczane.

\zrodla{\path{letsgolegacy.secondkey-nopcommerce-bench/scripts/README.md},
\path{letsgolegacy.secondkey-nopcommerce-bench/scripts/lib/NopBench.Common.ps1},
\path{letsgolegacy.secondkey-nopcommerce-bench/scripts/4-deploy-and-install.ps1},
\path{letsgolegacy.secondkey-nopcommerce-bench/docs/TOPOLOGY.md},
\path{letsgolegacy.secondkey-nopcommerce-bench/docs/adr/0003-legacy-host-iis-sqlexpress-unattended-install.md}}

\subsection{Build na \texttt{windows-2022} (ADR 0002)}\label{sec:bench-build}

\paragraph{Kontekst.} Wszystkie 31~projektów nopCommerce 3.90 celują w~.NET Framework
\textbf{4.5.1}. MSBuild kompiluje względem \emph{asemblacji referencyjnych} frameworku
docelowego, a~nie względem środowiska uruchomieniowego, a~pakiet docelowy (targeting pack)
4.5.1 nie jest już zainstalowany nigdzie, dokąd job CI może sięgnąć:

\begin{center}
\small
\begin{tabularx}{\linewidth}{@{}>{\raggedright\arraybackslash}p{4.4cm} >{\raggedright\arraybackslash}p{3.4cm} X@{}}
\toprule
Obraz hostowany (wrzesień 2026) & Visual Studio & Pakiety docelowe .NET Framework \\
\midrule
\texttt{windows-2022} & Enterprise 2022, 17.14 & 4.5.2, 4.6, 4.6.2, 4.7, 4.7.1, 4.7.2,
  4.8, 4.8.1 \\
\texttt{windows-2025} = \texttt{windows-latest} & Enterprise 2026, 18.x & 4.6.2, 4.7,
  4.7.1, 4.7.2, 4.8, 4.8.1 \\
\bottomrule
\end{tabularx}
\end{center}

Dane pochodzą z~plików README repozytorium \path{actions/runner-images}. Bez pakietu
build zatrzymuje się na błędzie:

\begin{lstlisting}
MSB3644: The reference assemblies for .NETFramework,Version=v4.5.1 were not found
\end{lstlisting}

Ograniczenie z~zadania: rozwiązać to bez modyfikowania kodu nopCommerce.

\paragraph{Decyzje.}
\begin{enumerate}
  \item \textbf{Asemblacje referencyjne z~NuGet, przekazane jako właściwość.} Krok~2
    pobiera z~nuget.org pakiet \texttt{Microsoft.NETFramework.ReferenceAssemblies.net451}
    1.0.3 (publikowany przez Microsoft pakiet, którego SDK używa w~tym samym celu),
    sprawdza jego przypięty SHA-256, rozpakowuje go w~katalogu roboczym i~przekazuje
    MSBuild \texttt{/p:TargetFrameworkRootPath=<pakiet>\textbackslash{}build}. MSBuild
    rozwiązuje wtedy \path{.NETFramework\v4.5.1} (asemblacje referencyjne, fasady,
    \path{RedistList\FrameworkList.xml}) z~pakietu zamiast z~\texttt{Program Files}. Nic nie
    jest instalowane na maszynie i~żaden plik upstream się nie zmienia; krok~2 sprawdza
    po buildzie, że checkout jest nadal czysty.
  \item \textbf{Runner \texttt{windows-2022}}, przypięty zamiast \texttt{windows-latest}:
    \begin{itemize}
      \item jego MSBuild 17.14 i~cele aplikacji webowych są najbliższym potomkiem toolchainu
        Visual Studio 2015, dla którego pisano projekty;
      \item \texttt{windows-latest} już raz się przesunął (na Windows Server 2025
        z~Visual Studio 2026); przypięta etykieta chroni build legacy przed zmianą pod
        spodem;
      \item skrypty nie zależą od niczego specyficznego dla obrazu (MSBuild jest szukany na
        \texttt{PATH} albo przez \texttt{vswhere}, asemblacje referencyjne pochodzą
        z~\path{pins.json}), więc przejście na \texttt{windows-2025} to zmiana jednej
        linii. Ręczny trigger workflowu ma wejście \texttt{runner}, żeby przenośność dało
        się sprawdzić na żądanie, bez macierzy przy każdym pushu.
    \end{itemize}
  \item \textbf{Te same skrypty lokalnie i~w~CI, w~Windows PowerShell 5.1.} CI wywołuje
    \path{scripts/1-fetch-nopcommerce.ps1} i~\path{scripts/2-build-legacy.ps1} bez
    dodatkowych argumentów. Akcje \texttt{microsoft/setup-msbuild} i~\texttt{NuGet/setup-nuget}
    (nuget.exe 7.9.0 -- ta sama wersja, której skrypty używają awaryjnie) tylko kładą
    narzędzia na \texttt{PATH}.
  \item \textbf{Publikacja tak, jak opisuje ją upstream.} \texttt{Deploying.Readme.txt}
    nopCommerce mówi: przebuduj rozwiązanie, potem opublikuj \texttt{Nop.Web} w~Release.
    Krok~2 robi dokładnie to z~wiersza poleceń (publikacja webowa do systemu plików
    z~\texttt{DeployDefaultTarget=WebPublish}). \texttt{Nop.Web.csproj} ustawia
    \texttt{FilesToIncludeForPublish=AllFilesInProjectFolder} i~własną listę wykluczeń,
    więc opublikowany katalog zawiera już zbudowane wtyczki (\path{Plugins\*}) i~obszar
    Administration, bez plików źródłowych, \path{*.csproj} i~plików stanu instalacji.
  \item \textbf{Build hermetyczny.} \path{ImportDirectoryBuildProps=false}
    i~\path{ImportDirectoryBuildTargets=false} nie pozwalają, by jakikolwiek plik
    \texttt{Directory.Build.*} nad checkoutem wpłynął na projekty upstream.
\end{enumerate}

\paragraph{Rozważone alternatywy.} Przestawienie projektów na 4.5.2 albo 4.8 modyfikuje
kod upstream i~zmienia to, co jest migrowane. Instalacja Developer Pack 4.5.1: framework
stracił wsparcie w~styczniu 2016, pakietu nie ma na żadnym obecnym obrazie, a~instalacja
w~każdym przebiegu jest wolna i~zmienia maszynę dewelopera. Skopiowanie asemblacji do
\texttt{Program Files (x86)\textbackslash{}Reference Assemblies\textbackslash{}…\textbackslash{}v4.5.1}
działa, ale pisze do miejsca ogólnomaszynowego, które przeżywa build, i~wymaga uprawnień
administratora -- zostaje jako udokumentowana ścieżka awaryjna dla narzędzi ignorujących
\texttt{TargetFrameworkRootPath}. \texttt{PackageReference} do pakietu asemblacji wymaga
edycji każdego projektu, bo projekty używają \texttt{packages.config}.
\texttt{windows-2019} wycofano z~runnerów hostowanych w~2025. Kontener Windows
(\path{mcr.microsoft.com/dotnet/framework/sdk}) dodaje warstwę do buildu oraz do IIS
i~SQL Server bez zysku: host ma MSBuild, a~problem asemblacji referencyjnych nie znika.

\paragraph{Konsekwencje.} Asemblacje referencyjne są częścią przypiętych wejść:
\path{pins.json} zapisuje wersję pakietu, URL i~SHA-256, a~\path{state\build.json} --
których użyto. Środowisko uruchomieniowe się nie zmienia: witryna działa na .NET Framework
4.8 maszyny, zastępującym 4.5.1 w~miejscu, jak na każdym produkcyjnym hoście tej
aplikacji. Jeśli GitHub wycofa \texttt{windows-2022}, domyślnym runnerem zostaje
\texttt{windows-2025} po zielonym przebiegu ręcznym z~\texttt{runner: windows-2025}.

\zrodla{\path{letsgolegacy.secondkey-nopcommerce-bench/docs/adr/0002-build-on-windows-2022-with-packaged-reference-assemblies.md},
\path{letsgolegacy.secondkey-nopcommerce-bench/scripts/2-build-legacy.ps1},
\path{letsgolegacy.secondkey-nopcommerce-bench/.github/workflows/legacy-build.yml}}

\subsection{Host legacy: IIS, SQL Server Express i~instalacja nienadzorowana (ADR 0003)}\label{sec:bench-host}

\paragraph{Kontekst.} P2 jest zrobione, gdy sklep odpowiada pod adresem URL: witryna
zbudowana w~P1 musi działać pod IIS z~SQL Server i~danymi przykładowymi nopCommerce,
bez nadzoru, na hostowanym runnerze Windows i~-- tymi samymi skryptami -- na stacji
roboczej. Nagranie i~odtworzenia, które następują potem (P4, P7), wymagają hosta
odtwarzalnego, startującego ze znanego stanu i~bez żadnego sekretu, który mógłby wyciec
do publicznego repozytorium albo logu.

\paragraph{1. SQL Server 2022 Express instalowany przez skrypt z~nośnika Microsoftu.}
Wybrana jest instancja nazwana \texttt{SQLEXPRESS} z~pliku \texttt{SQLEXPR\_x64\_ENU.exe}:
prawdziwa usługa Windows jak na produkcji, osiągalna dla tożsamości puli aplikacji IIS,
obsługująca snapshoty baz (zadanie rdzenia S11), z~małym nośnikiem (jeden plik), który
skrypt przypina przez SHA-256 i~sprawdza pod kątem ważnego podpisu Authenticode
Microsoftu; CI buforuje zweryfikowany plik. Odrzucone: LocalDB (instancje per użytkownik;
pula IIS wymagałaby ładowania profilu i~obejść współdzielenia, a~narzędzia odtwarzania
jako inny użytkownik nie widziałyby tej samej instancji -- nie na tym działa sklep legacy);
wydanie Developer (niepotrzebne funkcje, ponad 1~GB); kontener (Microsoft nie publikuje już
obrazów Windows dla SQL Server, a~kontenery Linux nie działają na runnerze Windows);
Chocolatey \texttt{sql-server-express} (ten sam nośnik i~przełączniki plus repozytorium
stron trzecich w~łańcuchu zaufania; jego suma kontrolna posłużyła tylko do zasiania pinu);
akcja GitHub w~rodzaju \texttt{mssqlsuite} (działałaby tylko w~CI, a~stacja robocza
i~tak potrzebuje skryptu).

Instalator działa z~\path{/SQLSYSADMINACCOUNTS=BUILTIN\Administrators}
(operator z~podniesionymi uprawnieniami jest administratorem),
\path{/SQLCOLLATION=SQL_Latin1_General_CP1_CI_AS} (kolejność i~równość w~bazie nie
mogą zależeć od ustawień regionalnych hosta) i~\path{/SKIPRULES=RebootRequiredCheck}
(włączenie IIS minutę wcześniej może zostawić oczekujący, niezwiązany restart). Pełną
listę przełączników podaje sekcja~\ref{sec:bench-krok3}.

\paragraph{2. Uwierzytelnianie Windows dla witryny, bez hasła w~żadnym connection
stringu.} Tożsamość puli aplikacji \texttt{IIS APPPOOL\textbackslash{}nopcommerce-legacy}
dostaje login Windows w~SQL Server wyłącznie z~rolą \textbf{\texttt{dbcreator}};
instalator nopCommerce tworzy bazę i~przez to jest jej właścicielem. Connection string
zapisywany przez nopCommerce w~\path{App_Data\Settings.txt} ma
\texttt{Integrated Security=True}: nie ma hasła SQL do generowania, przechowywania,
rotowania ani wycieku. Inne tożsamości, które później potrzebują bazy (reset stanu
odtwarzania, kandydat w~swojej puli), dostają własny login w~ten sam sposób.

\paragraph{3. IIS przez funkcje opcjonalne DISM, konfigurowany przez \texttt{appcmd}.}
\texttt{dism.exe /Online /Enable-Feature /All} z~nazwami funkcji IIS i~ASP.NET 4.x działa
tak samo w~Windows Server i~Windows 10/11, więc CI i~stacja robocza dzielą jedną ścieżkę
kodu. \texttt{appcmd.exe} jest dostarczany z~IIS i~zachowuje się tak samo w~Windows
PowerShell 5.1 i~PowerShell 7, w~przeciwieństwie do modułu \texttt{WebAdministration}.
Kompresja odpowiedzi nie jest instalowana, więc nagrane treści pozostają czytelne. Pula
aplikacji nie ma limitu bezczynności ani okresowego recyklingu: nagranie ani odtworzenie
nie może obejmować restartu obserwowanej aplikacji.

\paragraph{4. Dane przez instalator aplikacji, po HTTP.} Krok~4 wysyła formularz
instalacyjny nopCommerce 3.90 -- pola \texttt{InstallModel}
i~\path{Views/Install/Index.cshtml} -- dokładnie tak, jak wysyła go człowiek na
\texttt{/install}, z~włączonym tworzeniem bazy i~danymi przykładowymi oraz przypiętą
kolacją. Odrzucone alternatywy: uruchomienie SQL instalatora
(\path{App_Data\Install\Fast\*.sql}) wprost zasiałoby dane obok systemu zamiast przez
niego i~pominęłoby instalację wtyczek, którą wykonuje formularz; odtworzenie
przygotowanego \texttt{.bak} wprowadziłoby do łańcucha binarną bazę nieznanego
pochodzenia. Sukces nie jest przyjmowany na wiarę: formularz instalacji musiał zostać
najpierw podany (więc witryna nie była już zainstalowana), odpowiedzią musi być
przekierowanie instalatora na stronę główną, \path{App_Data\Settings.txt} musi istnieć,
strona główna musi odpowiedzieć treścią nopCommerce, a~baza musi zawierać produkty
przykładowe.

\paragraph{5. Czysta instalacja przy każdym przebiegu.} Krok~4 usuwa witrynę, pulę
aplikacji i~bazę przed wdrożeniem. Powtórzenie zbiega więc do tego samego stanu, zamiast
go kumulować -- na tej własności polegają nagranie (P4) i~odtworzenia (P7). Witryna jest
sprawdzana względem manifestu SHA-256 zapisanego przez build, więc to, co serwuje IIS,
jest dowodliwie tym, co zbudowano.

\paragraph{6. Jednorazowe dane administratora.} Instalator wymaga administratora sklepu;
jego hasło jest generowane na przebieg i~nigdy nie opuszcza
\path{<work>\secrets\legacy-admin.json} (sekcja~\ref{sec:bench-stan}).

\paragraph{7. Port 8080.} Port 80 zostaje przy Default Web Site; 8080 jest wolny na
hostowanym obrazie; proponowane porty proxy nagrywającego (8000) i~kandydata (8090) są
w~TOPOLOGY. Krok~4 odmawia startu, jeśli port jest zajęty.

\paragraph{Konsekwencje.}
\begin{itemize}
  \item Job hosta przy każdym przebiegu instaluje SQL Server Express (pobranie nośnika jest
    buforowane); podsumowanie joba podaje czas każdej fazy.
  \item Jeśli Microsoft podmieni plik pod przypiętym URL-em, krok~3 zatrzyma się na
    sprawdzeniu SHA-256 z~oboma skrótami w~komunikacie. Aktualizacja przypięcia jest
    wtedy zmianą przeglądaną w~review, nie automatyczną.
  \item \textbf{Zaobserwowane raz: instalator nie uruchamia silnika.} Pierwszy przebieg CI
    kroku~3 (przebieg 36229287543) zakończył się po 11~minutach błędem
    „Could not find the Database Engine startup handle” (\texttt{0x851A0019}); następny,
    z~identycznymi wejściami, zainstalował serwer w~2,5~minuty. Najlepiej znaną przyczynę
    tego błędu -- dysk systemowy zgłaszający sektory atomowe większe niż 4~KB, których SQL
    Server nie obsługuje -- wykluczono dla tego runnera: jego C: zgłasza 4096 bajtów.
    Krok~3 dlatego (a)~loguje rozmiary sektorów i~tylko wtedy, gdy dysk systemowy zgłasza
    więcej niż 4~KB, przenosi pliki danych na dysk katalogu roboczego; (b)~przy porażce
    drukuje \texttt{ERRORLOG} silnika i~zdarzenia SQL Server, żeby przyczyna była w~logu
    joba; (c)~usuwa na wpół zainstalowaną instancję i~uruchamia instalator jeszcze raz.
    Druga porażka kończy job błędem.
  \item Dane przykładowe to dane samego nopCommerce z~2017 roku. Zachowanie zależne od
    czasu (np. przykładowy rabat, który wygasł 2020-01-01) jest nagrywane takie, jakie
    jest.
\end{itemize}

\zrodla{\path{letsgolegacy.secondkey-nopcommerce-bench/docs/adr/0003-legacy-host-iis-sqlexpress-unattended-install.md},
\path{letsgolegacy.secondkey-nopcommerce-bench/scripts/3-install-iis-sql.ps1},
\path{letsgolegacy.secondkey-nopcommerce-bench/scripts/4-deploy-and-install.ps1}}

\subsection{Narzędzia łańcucha i~rozgrzewka na eShopLegacyMVC (ADR 0004, P3)}\label{sec:bench-rozgrzewka}

\paragraph{Kontekst.} P3 uruchamia cały łańcuch Second Key -- capture, replay, compare,
gate, evidence -- raz na przykładowej aplikacji legacy Microsoftu, zanim łańcuch zostanie
skierowany na nopCommerce, żeby problemy łańcucha wyszły na czymś małym. Zadanie jest
zrobione, gdy powstaje pakiet dowodowy. Trzeba było zdecydować o~trzech rzeczach, z~których
dwie przeżywają rozgrzewkę: skąd bench bierze narzędzia łańcucha i~jak nagranie jest
cięte na scenariusze.

\paragraph{1. Narzędzia łańcucha budowane z~przypiętych commitów, raz na przebieg
workflowu.} \repopath{sk} (Second Key) i~Portcullis (analyzery bramki) nie są jeszcze
publikowanymi pakietami. \path{pins.json} przypina każdy do commita
(\texttt{chain.secondKey}, \texttt{chain.portcullis}). Workflow wielokrotnego użytku
\path{chain-tools.yml} pobiera oba przez \texttt{actions/checkout} na tych commitach,
sprawdza pobrany commit, buduje oba projekty wiersza poleceń za pomocą .NET 10 SDK
wskazanego przez \path{global.json} repozytorium \repopath{sk} i~przekazuje je dalej jako
artefakt \texttt{chain-tools}. Build jest zależny od frameworku i~przenośny, więc te same
pliki działają na hoście Windows (capture, replay) i~w~Linuksie (compare, gate, evidence):
\texttt{dotnet SecondKey.Cli.dll}. Skrypty znajdują narzędzie przez
\texttt{NOPBENCH\_SK} / \texttt{NOPBENCH\_PORTCULLIS}, a~w~przeciwnym razie w~katalogu
\path{<work>\tools\chain\<tool>}. Przesunięcie przypięcia jest przeglądaną zmianą
\path{pins.json}, jak każde inne.

\paragraph{2. Jeden scenariusz to jedna sesja, nazwana nagłówkiem żądania.}
\texttt{sk replay} odtwarza nagranie scenariusz po scenariuszu, resetując każdą stronę
przed każdym z~nich i~dając każdemu świeży słoik ciasteczek (cookie jar); \texttt{sk
capture} rozstrzyga, do której sesji należy wymiana (exchange), na podstawie
skonfigurowanego ciasteczka albo nagłówka. Skrypty ruchu ustawiają
\texttt{x-bench-scenario: <id>} na każdym żądaniu scenariusza, a~capture dostaje ten
nagłówek jako klucz sesji. Konsekwencje:
\begin{itemize}
  \item lista scenariuszy planu ruchu \emph{jest} listą scenariuszy odtworzenia, jeden do
    jednego, a~werdykt da się odczytać wstecz: identyfikator sesji to \texttt{s-} i~pierwsze
    12 cyfr szesnastkowych SHA-256(\texttt{x-bench-scenario} + znak nowej linii + id),
    które krok nagrania drukuje dla każdego scenariusza;
  \item scenariusz przeżywa aplikację, która w~trakcie zmienia swoje ciasteczko sesji --
    nopCommerce wydaje nowe ciasteczko klienta przy logowaniu -- a~klucz oparty na
    ciasteczku rozciąłby go na dwa scenariusze z~resetem pomiędzy;
  \item jeden scenariusz ma jeden słoik ciasteczek, więc historia z~dwiema osobami to
    dwa scenariusze.
\end{itemize}
Nagłówek jest nagrywany i~odtwarzany jak każdy inny; aplikacje go ignorują.

\paragraph{3. Proxy nagrywające zatrzymuje się, gdy zawiera każdą wysłaną wymianę.}
Skrypty ruchu liczą każde żądanie wysłane przez proxy. \texttt{sk capture} zapisuje każdą
wymianę do pliku od razu, więc po ostatnim scenariuszu krok nagrania czeka, aż plik
będzie zawierał dokładnie tyle wymian, i~wtedy zatrzymuje proces. Niezgodna liczba
przerywa krok: wymiana, na którą aplikacja nie odpowiedziała, nie jest nagrana, a~taka,
której skrypty nie wysłały, nie jest ruchem ze skryptu. Opcja \texttt{-{}-exit-after}
wymaga liczby przed ruchem, a~czystego Ctrl+C nie da się dostarczyć procesowi w~tle
w~Windows bez pomocnika współdzielącego konsolę.

\paragraph{4. Aplikacja rozgrzewki: eShopLegacyMVC na własnych danych mock.}

\begin{center}
\small
\begin{tabularx}{\linewidth}{@{}l X@{}}
\toprule
Źródło & \texttt{dotnet-architecture/eShopModernizing} na
  \texttt{63bc9ec4414d7e281dc9f9d7bdcf70030950457d} (czubek \texttt{main} z~2023-10-25;
  repozytorium nie ma tagów), pobierane po identyfikatorze commita i~sprawdzane jak
  przypięcie nopCommerce \\
Aplikacja & \path{eShopLegacyMVCSolution/src/eShopLegacyMVC}: ASP.NET MVC 5.2.7 i~Web API
  na .NET Framework 4.7.2, Autofac, menedżer katalogu (lista, szczegóły, tworzenie,
  edycja, usuwanie) i~Web API marek \\
Dane & własny przełącznik aplikacji \texttt{UseMockData}, ustawiony na \texttt{true}
  w~\textbf{opublikowanym} \path{Web.config} (źródło się nie zmienia): katalog 12
  pozycji w~pamięci, bez bazy \\
Build & restore MSBuild (projekt używa \texttt{PackageReference}) i~publikacja webowa do
  systemu plików w~Release. Biblioteka narzędziowa celuje w~.NET Framework 4.6.1, którego
  pakietu docelowego obraz już nie ma; jak w~ADR 0002 asemblacje referencyjne 4.6.1
  i~4.7.2 pochodzą z~przypiętych pakietów NuGet, ułożone pod jednym
  \texttt{TargetFrameworkRootPath} \\
Host & IIS na runnerze \texttt{windows-2022}, witryna \texttt{eshop-legacy} na porcie
  8081, pula bez limitu bezczynności i~okresowego recyklingu; proxy nagrywające na 8001 \\
\bottomrule
\end{tabularx}
\end{center}

Katalog mock żyje na jednej liście w~pamięci przez cały czas życia procesu roboczego
(\texttt{CatalogServiceMock} zarejestrowany jako pojedyncza instancja), więc jedynym
resetem jest nowy proces: \path{scripts/eshop/Reset-EShop.ps1} zatrzymuje i~uruchamia pulę
aplikacji, po czym czeka na pełny katalog. Krok nagrania wywołuje go przed każdym
scenariuszem, a~\texttt{sk replay} -- przed każdym scenariuszem po każdej stronie (reset
\texttt{command} w~\path{warmup/eshop/secondkey.yaml}). Tworzenie, edycja i~usuwanie
sprawdzają token anti-forgery wydany z~formularzem; reguła korelacji odtwarzania przenosi
świeży token do formularza, który wysyła.

\paragraph{5. A/A i~bramka na całym drzewie.} Rozgrzewka nie ma kandydata. Obie strony
odtworzenia to ta sama witryna, więc werdykt pokazuje, co dokłada sam łańcuch: musi być
\texttt{pass}, z~każdą zaakceptowaną klauzulą małego kontraktu rozgrzewki spełnioną
i~żadną niewykonaną (unexercised). Portcullis skanuje kod legacy jako całe drzewo -- nie
ma pull requesta, do którego zmienionych linii można by go zawęzić. Jego znaleziska to
idiomy .NET Framework, które migracja musiałaby usunąć; \texttt{sk gate} raportuje, że
jego błędy zablokowałyby zmianę, a~krok zapisuje ten wynik, zamiast na nim padać. P8
bramkuje pull request kandydata tak, jak bramka ma być używana
(sekcja~\ref{sec:portcullis}).

\paragraph{6. Gdzie działa każdy krok.} Zgodnie z~TOPOLOGY: build, IIS, capture, replay
i~skan źródeł na runnerze Windows, obok aplikacji; compare, gate i~evidence w~Linuksie,
z~plików przesłanych przez job Windows.

\paragraph{Konsekwencje.} Każdy workflow, który uruchamia łańcuch, wywołuje
\path{chain-tools.yml}; aktualizacja narzędzia to jedna zmiana w~\path{pins.json}. Skrypty
ruchu muszą ustawiać nagłówek scenariusza na każdym żądaniu i~wysyłać każde żądanie przez
proxy; krok nagrania egzekwuje liczbę. Pakiet rozgrzewki jest artefaktem CI, a~nie
opublikowanym wynikiem; jego liczby są w~\path{results/p3/} (sekcja~\ref{sec:bench-wyniki}).

\zrodla{\path{letsgolegacy.secondkey-nopcommerce-bench/docs/adr/0004-chain-tools-and-eshop-warmup.md},
\path{letsgolegacy.secondkey-nopcommerce-bench/scripts/lib/NopBench.Chain.ps1},
\path{letsgolegacy.secondkey-nopcommerce-bench/scripts/eshop/Reset-EShop.ps1}}

\subsubsection{Ruch, kontrakt i~konfiguracja rozgrzewki}\label{sec:bench-rozgrzewka-ruch}

Ruch rozgrzewki to siedem scenariuszy z~\path{scripts/eshop/EShop.Scenarios.ps1} --
czynności osoby w~menedżerze katalogu eShopLegacyMVC. Każdy startuje z~tego samego
katalogu 12 pozycji (dane przykładowe \path{PreconfiguredData.cs}: pozycja~1 „.NET Bot
Black Hoodie” za \$19.50, pozycja~2 „.NET Black \& White Mug” za \$8.50, pozycja~3
„Prism White T-Shirt” za \$12.00; marki 1--5, z~których 2 to „.NET”; typy 1--4,
z~których 1 to „Mug”).

Scenariusze: \textbf{E01} -- przeglądanie (\texttt{/}, strony listy
\texttt{pageSize}/\texttt{pageIndex}, szczegóły pozycji 1 i~9, obraz
\texttt{/items/1/pic} jako \texttt{image/png}); \textbf{E02} -- pozycje, których nie ma
(szczegóły, edycja i~usuwanie pozycji 999 -- 404; \texttt{/Catalog/Details} bez
identyfikatora -- 400); \textbf{E03} -- utworzenie pozycji „Bench Mug” (\$9.99), która
pojawia się na liście („Showing 10 of 13 products”) i~pod \texttt{/Catalog/Details/13};
\textbf{E04} -- tworzenie z~błędnymi wartościami odrzucone z~trzema komunikatami walidacji;
\textbf{E05} -- zmiana ceny pozycji~2 na 10.50; \textbf{E06} -- usunięcie pozycji~3
(szczegóły dają potem 404, a~lista -- w~brzmieniu aplikacji „Showing 12 of 11 products”
-- jej nie zawiera); \textbf{E07} -- Web API marek (\texttt{/api/brands},
\texttt{/api/brands/2}, \texttt{/api/brands/42} -- 404).

Kontrakt rozgrzewki (\path{warmup/eshop/contract.yaml}, \texttt{metadata.name:
eshop-legacy-mvc-warmup}) jest celowo mały: ma ćwiczyć łańcuch. Porównuje nagłówki
\texttt{content-type} i~\texttt{location}, HTML -- przez ekstrakty, ma 6~ekstraktorów
(\texttt{catalogNames}, \texttt{catalogPrices}, \texttt{catalogTotal}, \texttt{itemName},
\texttt{itemPrice}, \texttt{validationMessages}) i~13~klauzul, z~czego 3~rodzaju \texttt{never} (23\%):
\path{ESHOP-NEVER-SERVER-ERROR}, \path{ESHOP-ITEM-PAGE-NEVER-WITHOUT-PRICE}
i~\path{ESHOP-REFUSED-FORM-NEVER-SILENT} (żadnej odpowiedzi 5xx, żadnej strony pozycji bez
ceny, żadnego odrzuconego formularza bez komunikatu walidacji) oraz 10 rodzaju
\texttt{must}: lista z~nazwami i~sumą, dodatnie ceny, stronicowanie 5 i~10 pozycji, 404
dla pozycji 999, 400 dla szczegółów bez identyfikatora, przekierowania po edycji
i~usunięciu, obrazek PNG, Web API z~pięcioma markami (w~tym „.NET”) i~404 dla marki 42.
Klauzule używają operatorów \texttt{gte}, \texttt{equals}, \texttt{absent},
\texttt{countEquals}, \texttt{countAtLeast}, \texttt{countAtMost}, \texttt{gt}
i~\texttt{matches}, a~klauzula marek -- kwantyfikatora \texttt{any} nad
\texttt{response.body.json[*].Brand}. Konfiguracja odtwarzania
\path{warmup/eshop/secondkey.yaml} ma \texttt{scenarioMode: session},
\texttt{timeoutSeconds: 120} (pierwsze żądanie po restarcie uruchamia aplikację
i~kompiluje widoki), po obu stronach reset \texttt{command} uruchamiający
\texttt{powershell -NoProfile -ExecutionPolicy Bypass -File}
\path{scripts\eshop\Reset-EShop.ps1} oraz regułę korelacji \texttt{csrf}, która wyciąga wyrażeniem regularnym wartość ukrytego
pola \texttt{\_\_RequestVerificationToken} ze strony i~wstawia ją do pola formularza
o~tej samej nazwie (tworzenie, edycja i~usuwanie sprawdzają token). Ścieżka polecenia resetu jest względna, więc \repopath{sk} uruchamia się z~korzenia
repozytorium. Sekcje \texttt{compare} i~\texttt{evidence} wskazują \path{contract.yaml}.
Kolejność kroków, skrypty i~ich produkty opisuje sekcja~\ref{sec:bench-skrypty-eshop};
liczby pierwszego pakietu -- sekcja~\ref{sec:bench-wyniki-p3}.

\zrodla{\path{letsgolegacy.secondkey-nopcommerce-bench/warmup/eshop/README.md},
\path{letsgolegacy.secondkey-nopcommerce-bench/warmup/eshop/contract.yaml},
\path{letsgolegacy.secondkey-nopcommerce-bench/warmup/eshop/secondkey.yaml},
\path{letsgolegacy.secondkey-nopcommerce-bench/scripts/eshop/EShop.Scenarios.ps1}}

\subsection{Nagranie ruchu nopCommerce (ADR 0005, P4)}\label{sec:bench-ruch}

P4 nagrywa zestaw ruchu z~\repopath{docs/P4-TRAFFIC-PLAN.md} -- 30~zaplanowanych
scenariuszy obejmujących katalog, koszyk, rabaty, podatki, wysyłkę i~checkout -- przez
\texttt{sk capture} przed sklepem z~P2. Zadanie jest zrobione, gdy liczba scenariuszy jest
zapisana. Nagranie jest wejściem wszystkiego, co następuje: P5 pisze na jego podstawie
kontrakt i~odtwarza je A/A, P7 odtworzy je na kandydacie. Plan opiera każdą regułę na
kodzie nopCommerce 3.90 i~cytuje go (\texttt{release-3.90}, commit \texttt{12d1f01};
ścieżki względem \path{src/}, numery linii tego commita): opisuje, jak zachowuje się
\emph{ten} sklep, a~nie jak sklep powinien się zachowywać w~ogóle. Identyfikatory
scenariuszy (\texttt{T}) i~reguł (\texttt{R}) są stabilne, żeby oba dokumenty mogły je
cytować.

\zrodla{\path{letsgolegacy.secondkey-nopcommerce-bench/docs/P4-TRAFFIC-PLAN.md},
\path{letsgolegacy.secondkey-nopcommerce-bench/docs/adr/0005-nopcommerce-traffic-recording.md}}

\subsubsection{Zasady nagrania}\label{sec:bench-ruch-zasady}

\begin{itemize}
  \item \textbf{Stan początkowy.} \path{scripts/4-deploy-and-install.ps1} (świeża witryna,
    świeża baza, dane przykładowe instalatora), potem konfiguracja z~sekcji~2 planu
    (sekcja~\ref{sec:bench-ruch-konfiguracja}), potem snapshot bazy SQL Server
    \texttt{nopcommerce\_legacy\_snapshot} robiony przez krok~7. Każdy scenariusz startuje
    z~tego snapshotu -- przy nagrywaniu oraz po stronie legacy i~kandydata w~każdym
    odtworzeniu -- więc scenariusze są niezależne od siebie i~od kolejności.
  \item \textbf{Wyłącznie przez proxy nagrywające} (\url{http://127.0.0.1:8000/}
    $\rightarrow$ \texttt{:8080}). Żądanie, które omija proxy, nie jest dowodem.
    \texttt{sk capture} to proxy nagrywające oparte na YARP, tylko dla ruchu
    przychodzącego, zapisujące \path{*.skcap}. System legacy pozostaje nietknięty: bez
    zmian w~kodzie, bez rekompilacji, bez modułu IIS.
  \item \textbf{Jeden scenariusz, jedna sesja.} Każdy scenariusz to jedna sesja
    przeglądarki z~własnym słoikiem ciasteczek, a~każde żądanie niesie nagłówek
    \texttt{x-bench-scenario: <id>} -- klucz sesji capture. Scenariusz jest więc jedną
    sesją w~\path{*.skcap} i~jednym scenariuszem w~odtworzeniu. Historia z~dwiema osobami
    to dwa scenariusze.
  \item \textbf{Jak przeglądarka.} Własne linki i~formularze sklepu; każdy formularz
    wysyłany ze wszystkimi polami, które daje strona; \texttt{\_\_RequestVerificationToken}
    czytany ze strony tam, gdzie jest sprawdzany. Sklep przykładowy włącza anti-forgery dla
    sklepu publicznego (\texttt{EnableXsrfProtectionForPublicStore = true}), co dotyczy
    akcji oznaczonych \texttt{[PublicAntiForgery]} -- w~tym zestawie szacowania wysyłki
    (\path{ShoppingCartController.GetEstimateShipping}) i~rejestracji
    (\path{CustomerController.Register}). Dodawanie do koszyka, zmiany w~koszyku,
    kupony, logowanie i~checkout jednostronicowy go nie sprawdzają.
  \item \textbf{Host.} \texttt{sk capture} przekazuje dalej host celu (domyślne
    zachowanie YARP), więc bezwzględne adresy, które nopCommerce buduje z~hosta żądania
    (adresy skryptów checkoutu jednostronicowego), wskazują \texttt{:8080}. Przeglądarka,
    która by za nimi poszła, opuściłaby proxy; sesje ze skryptu żądają wyłącznie ścieżek
    względnych, a~przekierowania sprowadzają do ścieżki i~zapytania
    (\texttt{Get-BenchLocationPath}), więc nic nie omija nagrania.
  \item \textbf{Sondy (probes).} Język kontraktu zawęża klauzulę wyłącznie metodą,
    ścieżką i~zapytaniem, więc żądanie, na którego odpowiedzi P5 stwierdza fakt właściwy
    jednemu scenariuszowi, niesie w~zapytaniu \texttt{probe=<nazwa>}
    (\path{/cart?probe=T11}, \path{/cart/estimateshipping?probe=T26-1024}). nopCommerce
    ignoruje nieznane parametry zapytania na tych trasach; parametr jedynie nazywa żądanie.
  \item \textbf{Deterministyczne wejścia.} Stałe produkty, ilości, adresy, kody kuponów
    i~klienci -- R1 rejestruje się zawsze jako \texttt{r1@nopbench.invalid}; każdy
    scenariusz startuje ze snapshotu, więc adres jest zawsze wolny.
  \item \textbf{Brak wywołań wychodzących.} Używane są tylko metody płatności offline
    (Check / Money Order, Manual) i~dostawcy stawek offline; PayPal i~wtyczki
    przewoźników pozostają bezczynne, więc zadanie rdzenia S16 (zaślepki wywołań
    wychodzących) nie jest temu zestawowi potrzebne.
\end{itemize}

\zrodla{\path{letsgolegacy.secondkey-nopcommerce-bench/docs/P4-TRAFFIC-PLAN.md},
\path{letsgolegacy.secondkey-nopcommerce-bench/docs/TOPOLOGY.md},
\path{letsgolegacy.secondkey-nopcommerce-bench/scripts/lib/NopBench.Chain.ps1}}

\subsubsection{Konfiguracja sklepu przed nagraniem}\label{sec:bench-ruch-konfiguracja}

Dane przykładowe wyłączają reguły, o~które chodzi w~zestawie:
\begin{itemize}
  \item \textbf{Podatek jest wszędzie zerowy.} \texttt{Tax.FixedOrByCountryStateZip}
    instaluje się w~trybie stałej stawki (\texttt{CountryStateZipEnabled = false}) i~nie
    zapisuje żadnej stawki, więc każda kategoria podatku daje 0.
  \item \textbf{Wysyłka nic nie kosztuje.} \texttt{Shipping.FixedOrByWeight} instaluje się
    w~trybie stałej stawki bez zapisanej stawki, a~brak stawki daje 0.
  \item \textbf{Żaden kupon przykładowy nie zmienia ceny.} \texttt{123} to rabat \$10
    „przypisany do SKU”, do którego nie jest przypisany żaden produkt; \texttt{456} (20\%
    sumy zamówienia) wygasł 2020-01-01.
\end{itemize}

\paragraph{Decyzja (ADR 0005, 1): sklep konfigurowany przez własny panel
administracyjny, przed nagraniem, z~odczytem zwrotnym.}
\path{scripts/6-configure-store.ps1} loguje się jako administrator utworzony w~kroku~4
i~używa stron administracyjnych oraz punktów AJAX, które klika człowiek -- nigdy SQL --
więc każdy wiersz, który zapisuje nopCommerce, jest wierszem, który zapisałby sam, razem
z~pamięcią podręczną i~zdarzeniami. Każda zmiana jest odczytywana z~powrotem, a~formularz
ustawień lub produktu wysłany z~powrotem jest ponownie czytany i~porównywany pole po polu,
więc nic innego nie zmienia się niezauważenie. Żądania administratora idą wprost do
witryny, nigdy przez proxy nagrywające: konfiguracja nie jest ruchem, a~hasło
administratora nie może trafić do nagrania.

{\footnotesize
\begin{xltabular}{\linewidth}{@{}>{\raggedright\arraybackslash}p{2.4cm} >{\raggedright\arraybackslash}X >{\raggedright\arraybackslash}p{2.1cm}@{}}
\caption{Konfiguracja sklepu przed nagraniem (plan, sekcja~2)}\label{tab:bench-konfig}\\
\toprule
Obszar & Konfiguracja & Używa \\
\midrule
\endfirsthead
\toprule
Obszar & Konfiguracja & Używa \\
\midrule
\endhead
\bottomrule
\endlastfoot
Dostawca podatku & przełączenie na „by country / state / zip”. Kategoria podatku
  \emph{Electronics \& Software}: US / New York / zip \texttt{10001} $\rightarrow$
  9,000\%; US / New York / dowolny zip $\rightarrow$ 8,875\%; US / California
  $\rightarrow$ 7,250\%; US / dowolny stan $\rightarrow$ 5,000\%; Canada / Ontario / zip
  \texttt{K1A 0B1} $\rightarrow$ 13,000\%. Nic dla Germany. „Apple MacBook Pro 13-inch”
  zwolniony z~podatku & T21--T23 \\
Wysyłka & stała stawka na metodę: Ground \$10.00, 2nd Day Air \$25.00, Next Day Air
  \$40.00. Darmowa wysyłka powyżej \$1,000.00, bez podatku & T24--T26 \\
Rabaty (każdy wymaga swojego kuponu, więc żaden nie przecieka do innych scenariuszy) &
  \texttt{SAVE10} 10\% sumy częściowej zamówienia · \texttt{TAKE50} \$50.00 od sumy
  zamówienia · \texttt{PCT20MAX100} 20\% sumy częściowej, najwyżej \$100.00 ·
  \texttt{SHIPFREE} 100\% wysyłki · \texttt{ONCE5} \$5.00 od sumy zamówienia, 1~raz na
  klienta · \texttt{FUTURE10} 10\% sumy częściowej od 2099-01-01 · \texttt{MEMBERS15}
  15\% sumy częściowej z~wymaganiem „customer role is Registered”
  (\texttt{DiscountRules.CustomerRoles}) & T11--T20 \\
Zadania w~tle & wyłączenie zadań harmonogramu włączonych przez instalację przykładową --
  „Send emails” (co 60~s), „Keep alive” (300~s), „Delete guests” (600~s), „Update
  currency exchange rates” (3600~s, wywołanie wychodzące) -- i~restart aplikacji. Wersja
  3.90 uruchamia je na timerach wewnątrz aplikacji webowej, a~\texttt{Task.Execute} czyta
  swój wiersz przed blokiem \texttt{try}: przy przywracaniu bazy ten odczyt rzuca wyjątek
  na wątku timera i~kładzie proces roboczy, a~sklep odpowiada 500, aż się zrestartuje.
  Każdy scenariusz zaczyna się od przywrócenia & każdy scenariusz \\
Katalog & wycofanie publikacji „HP Envy 6-1180ca 15.6-Inch Sleekbook” i~wyłączenie
  ustawienia katalogu „Allow viewing of unpublished product details page” -- instalacja
  przykładowa je włącza (\texttt{AllowViewUnpublishedProductPage = true}), przez co strona
  nieopublikowanego produktu jest podawana każdemu z~oznaczeniem „discontinued”. Stan
  „Lenovo IdeaCentre 600 All-in-One PC” ustawiony na 3 (produkt już zarządza stanem bez
  zamówień oczekujących) & T02, T05--T07, T28 \\
\end{xltabular}
}

Waluta (USD, format wyświetlania \texttt{en-US}), wyświetlanie podatku (bez podatku),
zaokrąglanie cen w~trakcie obliczeń, anonimowy checkout i~checkout jednostronicowy
zostają takie, jak po instalacji. Krok~6 tworzy każdy rabat formularzem
\path{/Admin/Discount/Create} z~\texttt{RequiresCouponCode = true} i~nazwą
„Bench <kod>”; typ to \texttt{DiscountTypeId} 20 (suma częściowa zamówienia) dla
\texttt{SAVE10}, \texttt{PCT20MAX100} (z~\texttt{MaximumDiscountAmount = 100}),
\texttt{FUTURE10} (z~\texttt{StartDateUtc = 2099-01-01 00:00}) i~\texttt{MEMBERS15},
1 (suma zamówienia) dla \texttt{TAKE50} i~\texttt{ONCE5} (z~\texttt{DiscountLimitationId
= 25}, „N razy na klienta”, i~\texttt{LimitationTimes = 1}) oraz 10 (wysyłka) dla
\texttt{SHIPFREE}. Wymaganie roli Registered dla \texttt{MEMBERS15} dołącza wtyczka
\texttt{DiscountRulesCustomerRoles} do domyślnej grupy wymagań rabatu.

Zadania harmonogramu są uruchamiane z~włączonych wierszy przy starcie aplikacji -- stąd
restart. Krok~6 wyłącza wszystkie, restartuje aplikację przez „Restart application”
w~panelu i~czeka na nowy wpis „Application started” w~logu sklepu. Żadne z~czterech zadań
(wysyłanie e-maili z~kolejki, pingowanie witryny, usuwanie gości, aktualizacja kursów
walut) nie jest częścią tego, co ćwiczy zestaw ruchu.

\zrodla{\path{letsgolegacy.secondkey-nopcommerce-bench/docs/P4-TRAFFIC-PLAN.md},
\path{letsgolegacy.secondkey-nopcommerce-bench/docs/adr/0005-nopcommerce-traffic-recording.md},
\path{letsgolegacy.secondkey-nopcommerce-bench/scripts/6-configure-store.ps1}}

\subsubsection{Persony i~scenariusze T01--T30}\label{sec:bench-ruch-scenariusze}

\begin{center}
\small
\begin{tabularx}{\linewidth}{@{}l l X@{}}
\toprule
Persona & Kto & Uwagi \\
\midrule
G1 & gość robiący zakupy & własny słoik ciasteczek \\
G2 & drugi gość & próbuje dostać się do cudzego zamówienia \\
R1 & klient, który rejestruje się w~T28 & \texttt{r1@nopbench.invalid}, z~hasłem podanym
  przy rejestracji \\
S1 & klient przykładowy \texttt{steve\_gates@nopCommerce.com} & zarejestrowany; właściciel
  przykładowego zamówienia~1; hasło nadawane przez instalator klientom przykładowym --
  publiczne dane przykładowe, nie sekret \\
\bottomrule
\end{tabularx}
\end{center}

Przykładowe zamówienie~2 należy do \texttt{arthur\_holmes@nopCommerce.com}, który nigdy
nie działa; jego zamówienie jest „cudzym zamówieniem” w~T30. Administrator konfiguruje
sklep i~nie występuje w~żadnym scenariuszu. Tabela~\ref{tab:bench-scenariusze} to plan
scenariuszy (plan podaje dla każdego także miejsce w~kodzie 3.90, które go uzasadnia).

{\footnotesize
\begin{xltabular}{\linewidth}{@{}l >{\raggedright\arraybackslash}p{1.35cm} l >{\raggedright\arraybackslash}X >{\raggedright\arraybackslash}X@{}}
\caption{Scenariusze zestawu ruchu (plan, sekcja~4)}\label{tab:bench-scenariusze}\\
\toprule
Id & Obszar & Persona & Kroki & Ćwiczona reguła biznesowa \\
\midrule
\endfirsthead
\toprule
Id & Obszar & Persona & Kroki & Ćwiczona reguła biznesowa \\
\midrule
\endhead
\bottomrule
\endlastfoot
T01 & Katalog & G1 & \texttt{GET /} & strona główna wymienia produkty oznaczone „show on
  home page”; atrybucja w~stopce \\
T02 & Katalog & G1 & \texttt{GET /notebooks}, potem \texttt{?orderby=10},
  \texttt{?orderby=5}, \texttt{?pagesize=3\&pagenumber=2} & lista kategorii: tylko
  produkty opublikowane, sortowane po cenie albo nazwie, stronicowane \\
T03 & Katalog & G1 & \texttt{GET /computers} & kategoria nadrzędna wymienia swoje
  podkategorie \\
T04 & Katalog & G1 & \texttt{GET} \path{/lenovo-ideacentre-600-all-in-one-pc},
  \path{/apple-macbook-pro-13-inch}, \path{/build-your-own-computer} & strona produktu:
  nazwa, SKU, cena, atrybuty z~dopłatami \\
T05 & Katalog & G1 & \texttt{GET} \path{/hp-envy-6-1180ca-156-inch-sleekbook} (nieopublikowany
  w~konfiguracji) & produkt nieopublikowany nie jest pokazywany \\
T06 & Wyszuki\-wanie & G1 & \texttt{GET /search?q=apple}, \texttt{?q=APPLE},
  \texttt{?q=AP\_MBP\_13}, \texttt{?q=ap}, \texttt{?q=\%20apple\%20} & wyszukiwanie po
  nazwach (kolacja bazy), dokładne dopasowanie SKU, przycięte wejście, minimalna długość~3 \\
T07 & Koszyk & G1 & Lenovo z~\texttt{/desktops} (ilość~1), potem ze strony produktu
  z~ilością~5 & dodanie do koszyka; limit stanu bez zamówień oczekujących \\
T08 & Koszyk & G1 & MacBook z~ilością~1, potem 2 & minimalna ilość zamówienia
  (w~danych przykładowych 2) \\
T09 & Koszyk & G1 & „Build your own computer” bez wymaganych atrybutów, potem
  z~wybranymi Processor, RAM i~HDD & atrybuty wymagane; dopłaty za atrybuty w~cenie
  jednostkowej \\
T10 & Koszyk & G1 & dwie pozycje; \texttt{POST /cart} \texttt{updatecart}
  z~\texttt{itemquantity\{id\}=3}; potem zaznaczenie \texttt{removefromcart} na jednej
  pozycji; potem \texttt{itemquantity\{id\}=0} na drugiej & zmiana ilości; usunięcie
  polem wyboru; ilość 0 usuwa \\
T11 & Rabaty & G1 & Lenovo × 2 (\$1,000.00); \texttt{POST /cart}
  \texttt{applydiscountcouponcode} z~\texttt{SAVE10} & rabat procentowy od sumy
  częściowej zamówienia \\
T12 & Rabaty & G1 & kod „\texttt{\textvisiblespace{}save10\textvisiblespace{}}” (spacje, małe litery); na świeżym koszyku
  \texttt{SAVE10} z~U+180E na końcu & kod kuponu jest przycinany i~porównywany bez względu
  na wielkość liter \\
T13 & Rabaty & G1 & kody \texttt{456} (przykładowy, wygasły), \texttt{FUTURE10},
  \texttt{NOSUCH} & kupony wygasłe, jeszcze nieaktywne i~nieznane są odrzucane, każdy
  z~własnym komunikatem \\
T14 & Rabaty & G1 & „Pride and Prejudice” (\$24.00, wysyłany, bez stawki podatku dla
  Books); checkout do potwierdzenia z~Ground; \texttt{TAKE50} & stały rabat od sumy
  zamówienia większy niż całe zamówienie \\
T15 & Rabaty & G1 & MacBook × 2 (\$3,600.00); \texttt{PCT20MAX100} & maksymalna kwota
  rabatu procentowego \\
T16 & Rabaty & G1 & Lenovo × 1; checkout do kroku potwierdzenia z~Ground;
  \texttt{SHIPFREE} & rabat na wysyłkę \\
T17 & Rabaty & G1 & „\$25 Virtual Gift Card” i~„Pride and Prejudice” w~koszyku;
  \texttt{SAVE10} & rabaty od sumy częściowej i~całkowitej są odrzucane, gdy w~koszyku
  jest karta podarunkowa \\
T18 & Rabaty & G1, S1 & \texttt{MEMBERS15} jako G1; logowanie jako S1
  i~\texttt{MEMBERS15} & wymaganie rabatu: rola klienta \\
T19 & Rabaty & S1 & zamówienie z~\texttt{ONCE5}; nowy koszyk i~ponownie \texttt{ONCE5} &
  ograniczenie „N razy na klienta”, śledzone dla klientów zarejestrowanych \\
T20 & Rabaty & G1 & przykładowy kupon \texttt{123} & rabat na SKU nieprzypisany do żadnego
  produktu jest przyjmowany i~niczego nie zmienia \\
T21 & Podatki & G1 & Lenovo × 1; checkout z~adresem rozliczeniowym US / New York /
  \texttt{10001}; stop na potwierdzeniu & stawka według kraju, stanu i~dokładnego kodu
  pocztowego \\
T22 & Podatki & G1 & jak T21 z~US / New York / \texttt{10002}, potem US / California,
  US / Texas, Germany; potem MacBook × 2 z~US / New York / \texttt{10001} & stawki
  zastępcze (dowolny kod, dowolny stan, brak stawki) i~produkt zwolniony z~podatku \\
T23 & Podatki & G1 & jak T21 z~Canada / Ontario / \texttt{k1a 0b1} (małe litery) &
  dopasowanie kodu pocztowego ignoruje wielkość liter \\
T24 & Wysyłka & G1 & Lenovo × 1: checkout do kroku metody wysyłki; osobno szacowanie
  wysyłki na \texttt{/cart} dla US / New York / \texttt{10001} & metody i~stałe stawki;
  szacunek zgodny z~checkoutem \\
T25 & Wysyłka & G1 & tylko MacBook × 2 (produkt z~darmową wysyłką); szacowanie & darmowa
  wysyłka, gdy każda pozycja ma darmową wysyłkę \\
T26 & Wysyłka & G1 & Lenovo × 2 (dokładnie \$1,000.00), szacowanie; dodanie „Pride and
  Prejudice”, ponowne szacowanie & „darmowa wysyłka powyżej X” znaczy ściśle powyżej \\
T27 & Checkout & G1 & checkout jednostronicowy gościa: Lenovo × 1, adres rozliczeniowy
  US / New York / \texttt{10001}, wysyłka na ten sam adres, Ground, Check / Money Order,
  potwierdzenie; \texttt{GET /checkout/completed/\{id\}}, \texttt{GET /orderdetails/\{id\}}
  & checkout gościa od początku do końca; sumy przeliczane na serwerze \\
T28 & Checkout & R1 & rejestracja na \texttt{/register}; Lenovo × 2 z~\texttt{SAVE10};
  Next Day Air; płatność Manual, najpierw kartą \texttt{4111 1111 1111 1112}, potem
  \texttt{4111 1111 1111 1111}; potwierdzenie; strona Lenovo & checkout klienta
  zarejestrowanego; cyfra kontrolna karty; stan maleje (3 $\rightarrow$ 1); koszyk
  opróżniony \\
T29 & Zamówie\-nia & S1 & logowanie; \texttt{GET /order/history},
  \texttt{/orderdetails/1}, \texttt{/orderdetails/pdf/1} & historia zamówień i~szczegóły
  własnego zamówienia \\
T30 & Zamówie\-nia & G1, G2, S1 & G1 składa zamówienie gościa (jak T27); G2 żąda
  \texttt{/orderdetails/\{id\}}, \texttt{/orderdetails/print/\{id\}},
  \texttt{/orderdetails/pdf/\{id\}}, \texttt{/reorder/\{id\}} i~\texttt{/order/history};
  S1 żąda \texttt{/orderdetails/2} & nikt nie widzi ani nie ponawia cudzego zamówienia;
  goście nie mają historii zamówień \\
\end{xltabular}
}

\zrodla{\path{letsgolegacy.secondkey-nopcommerce-bench/docs/P4-TRAFFIC-PLAN.md}}

\subsubsection{Nagranie: sesje, nagłówek scenariusza, reset i~sondy}\label{sec:bench-ruch-nagranie}

Nagranie realizuje plan z~doprecyzowaniami, z~których każde wymusiła zasada „jeden
scenariusz = jedna sesja startująca ze snapshotu” albo to, jak zachował się sklep.
\finding{Nagrano 38~sesji scenariuszy obejmujących 30~zaplanowanych scenariuszy
i~290~wymian.}

\begin{center}
\small
\begin{tabularx}{\linewidth}{@{}>{\raggedright\arraybackslash}p{2.3cm} >{\raggedright\arraybackslash}X >{\raggedright\arraybackslash}X@{}}
\toprule
Plan & Nagrane jako & Dlaczego \\
\midrule
T12 & T12 (\texttt{\textvisiblespace{}save10\textvisiblespace{}}), T12B (\texttt{SAVE10}+U+180E) & „na świeżym koszyku”
  oznacza świeżą sesję \\
T18 & T18A (G1), T18B (S1) & dwie osoby, dwie sesje \\
T22 & T22A (US/NY/\texttt{10002}), T22B (US/CA), T22C (US/TX), T22D (Germany), T22E
  (MacBook × 2, US/NY/\texttt{10001}) & jeden adres na checkout, każdy ze snapshotu \\
-- & T22F: rozliczenie na US/NY/\texttt{10001}, wysyłka na US/TX; podatek nowojorski &
  R30 („podatek idzie za adresem rozliczeniowym”) wymaga adresu rozliczeniowego innego niż
  adres wysyłki; w~T21--T23 są takie same \\
T27 & \texttt{GET /checkout/completed/} bez identyfikatora & tego żąda skrypt checkoutu
  (\texttt{ConfirmOrder.init(…, '…checkout/completed/')}); strona pokazuje ostatnie
  zamówienie klienta \\
T28 „strona Lenovo” & dodanie 2 kolejnych Lenovo jest odrzucane: „The maximum quantity
  that can be added is 1.” & pokazuje spadek stanu 3 $\rightarrow$ 1 w~odpowiedzi, a~nie
  tylko na stronie \\
T30 & T30A (G2: szczegóły, wydruk, PDF i~ponowienie zamówienia~2 oraz
  \texttt{/order/history}), T30B (S1: \texttt{/orderdetails/2}) & zamówienie G1 nie
  istniałoby w~sesji G2 -- każdy scenariusz startuje ze snapshotu -- więc „cudzym
  zamówieniem” jest przykładowe zamówienie~2 Arthura Holmesa \\
\bottomrule
\end{tabularx}
\end{center}

\paragraph{Reset przed scenariuszem (ADR 0005, 2).} Po kroku~6
\path{scripts/7-record-traffic.ps1} robi snapshot bazy i~przywraca go przed każdym
scenariuszem -- dokładnie tak, jak reset \texttt{sqlServerSnapshot} w~\texttt{sk replay}
robi to przed każdym scenariuszem po każdej stronie. W~efekcie nagranie i~każde
odtworzenie zaczynają każdy scenariusz od tych samych wierszy, łącznie z~wartościami
identity -- zamówienie złożone w~scenariuszu dostaje za każdym razem ten sam numer;
scenariusze są niezależne od siebie i~od kolejności, a~każdy da się nagrać osobno
(\texttt{-Only}); historia z~dwiema osobami nie może zakładać, że zamówienie pierwszej
istnieje w~sesji drugiej (stąd T30 na przykładowym zamówieniu). Przywrócenie zamyka każde
połączenie z~bazą, w~tym połączenia z~puli sklepu. Krok~7 wysyła po każdym przywróceniu,
przed scenariuszem, jedno żądanie wprost do witryny i~zapisuje odpowiedź: to dowód, czy
odtworzenie może przywrócić bazę i~od razu ruszyć. To samo wymaga, by nic nie używało
bazy według własnego harmonogramu -- stąd wyłączone zadania w~tle
(sekcja~\ref{sec:bench-ruch-konfiguracja}).

\paragraph{Sondy w~zapytaniu (ADR 0005, 3).} Wiele reguł planu to fakty o~odpowiedzi
jednego scenariusza pod adresem, którego używa każdy scenariusz -- koszyk po
\texttt{SAVE10} (\texttt{POST /cart}), szacunek dla koszyka za \$1,000.00
(\texttt{POST /cart/estimateshipping}). Takie żądania niosą dodatkowy parametr
\texttt{probe=<scenariusz lub krok>}, który nopCommerce ignoruje na tych trasach,
a~kontrakt może po nim wybierać. Alternatywa -- jedna klauzula na URL, stwierdzająca to,
co wspólne wszystkim scenariuszom -- sprowadziłaby większość reguł do „odpowiada 200”.
Zmiana żądania jest dopuszczalna, bo sonda jest częścią nagrania: legacy i~kandydat
dostają te same bajty. Sondę niosą: T07, T08, T09 (odrzucone dodanie i~koszyk);
T10-quantity, T10-remove, T10-zero; T11, T12, T12B, T13-expired, T13-future,
T13-unknown, T14, T15, T16, T17, T18A, T18B, T19, T20 (odpowiedź na kupon albo koszyk);
T21, T22A--T22F, T23 (koszyk po checkoucie); T24 (szacunek i~odpowiedź kroku
rozliczenia, która wymienia metody wysyłki); T25, T26-1000, T26-1024 (szacunki); T27
(strona zakończenia i~szczegóły zamówienia); T28 (odrzucona karta, strona zakończenia,
szczegóły zamówienia, pusty koszyk \texttt{T28-after}, odrzucone dodanie); T29 (historia
zamówień); T30A i~T30B (każde żądanie: historia zamówień gościa jest odrzucana pod
adresem, którego używa też Steve, więc każde potrzebuje nazwy).

\paragraph{Scenariusz, który zawiedzie, nie zatrzymuje nagrania (ADR 0005, 4).} Każdy
scenariusz sprawdza po drodze to, czego potrzebuje (\texttt{Invoke-BenchStep},
\texttt{Assert-NopTotal}): scenariusz, który nie dotarł do tego, co ćwiczy, nie jest
dowodem. Krok~7 uruchamia każdy scenariusz także po porażce innego, drukuje własny log
nopCommerce dla nieudanego, zanim następne przywrócenie go wymaże, i~na końcu kończy się
błędem z~listą wszystkich nieudanych scenariuszy. Krok~6 robi to samo dla obszarów
(podatki, katalog, wysyłka, rabaty, zadania w~tle). Jeden przebieg pokazuje wszystkie
problemy.

\paragraph{Konsekwencje.} Zestaw ruchu to 38~sesji dla 30~zaplanowanych scenariuszy.
Konfiguracja odtwarzania P5 wskazuje ten sam snapshot i~bazę w~resecie
\texttt{sqlServerSnapshot}. \path{*.skcap} zawiera wyłącznie dane sklepu przykładowego:
opublikowane hasło klienta przykładowego, hasło, z~którym rejestruje się R1, testowe
numery kart. Ciasteczka są domyślnie redagowane przez \texttt{sk capture}.

\zrodla{\path{letsgolegacy.secondkey-nopcommerce-bench/docs/P4-TRAFFIC-PLAN.md},
\path{letsgolegacy.secondkey-nopcommerce-bench/docs/adr/0005-nopcommerce-traffic-recording.md},
\path{letsgolegacy.secondkey-nopcommerce-bench/scripts/7-record-traffic.ps1},
\path{letsgolegacy.secondkey-nopcommerce-bench/scripts/traffic/NopCommerce.Scenarios.ps1}}

\subsubsection{Znaleziska z~nagrywania}\label{sec:bench-ruch-znaleziska}

Napisanie i~nagranie zestawu pokazało rzeczy, które P5 stwierdza w~kontrakcie, a~P7
porówna:
\begin{itemize}
  \item \textbf{T05 wymaga ustawienia katalogu.} Samo wycofanie publikacji zostawia stronę
    produktu otwartą dla wszystkich (\texttt{AllowViewUnpublishedProductPage = true}
    w~ustawieniach przykładowych).
  \item \textbf{Przywrócenie bazy przy działającym zadaniu harmonogramu kładzie sklep} --
    stąd wiersz zadań w~tle w~konfiguracji. Widziane raz w~pierwszym pełnym nagraniu:
    pierwsze żądania po przywróceniu dostały 500, potem aplikacja wystartowała ponownie.
  \item \textbf{Checkout zaczyna się od przycisku Checkout w~koszyku.} Sklep przykładowy
    ma wymagany atrybut zamówienia „Gift wrapping”, zapisywany przy wysłaniu formularza
    koszyka (\texttt{ShoppingCartController.StartCheckout}); zamówienie klienta, który
    poszedł wprost na \texttt{/onepagecheckout}, jest odrzucane z~„Please select Gift
    wrapping;”. Każdy checkout w~zestawie idzie: koszyk $\rightarrow$ Checkout
    $\rightarrow$ (gość: „Checkout as Guest”) $\rightarrow$ \texttt{/checkout}
    $\rightarrow$ checkout jednostronicowy, jak w~przeglądarce.
  \item \textbf{Szacunek wysyłki oferuje też odbiór osobisty „Pickup (\$1.99)”} --
    przykładowy punkt odbioru -- a~darmowa wysyłka nie znosi tej opłaty: koszyk MacBooków
    jest szacowany na \$0.00 dla każdej metody wysyłki i~\$1.99 dla odbioru (T25). R33
    i~R34 dotyczą metod wysyłki; opcja odbioru jest nagrana taka, jaka jest.
  \item \textbf{T18A: wymaganie roli klienta odmawia z~pustym komunikatem.} Reguła nie
    ustawia \texttt{UserError}, \texttt{ValidateDiscount} przekazuje ten \texttt{null}
    dalej, a~\texttt{ApplyDiscountCoupon} pokazuje go jako jedyny komunikat -- pole kuponu
    zawiera jeden pusty komunikat porażki, a~nie „The coupon code you entered couldn't be
    applied”. R24 jest spełniona; komunikat to zachowanie legacy, które migracja może
    „naprawić”.
  \item \textbf{T12B: \texttt{SAVE10} z~U+180E jest odrzucany} („The coupon code you
    entered couldn't be applied to your order”) na .NET Framework 4.8.
  \item \textbf{Kwoty ujemne są zapisywane w~nawiasach}: \texttt{(\$100.00)},
    \texttt{(\$34.00)}, \texttt{(\$75.00)} -- format walutowy \texttt{en-US} .NET Framework
    4.8 na Windows Server 2022 (pierwsza sonda NLS $\rightarrow$ ICU). Skrypty przy
    nagrywaniu przyjmują obie formy, \texttt{(\$x)} i~\texttt{-\$x}
    (\texttt{Get-NopNegative}); to kontrakt przypina formę, którą pisze system legacy.
  \item \textbf{Zerowa suma zamówienia pomija kroki płatności} -- decyduje o~tym
    \path{CheckoutController.OpcLoadStepAfterShippingMethod}. Dlatego T14 najpierw
    przechodzi checkout, a~\texttt{TAKE50} stosuje na końcu, zgodnie z~planem.
\end{itemize}

Dwa z~tych znalezisk zmieniły sposób nagrywania, a~nie tylko to, co stwierdza P5: krok~6
wyłącza zadania w~tle i~restartuje aplikację, a~krok~7 sprawdza pierwszą odpowiedź po
każdym przywróceniu; każdy checkout zaczyna się od przycisku Checkout w~koszyku.

\zrodla{\path{letsgolegacy.secondkey-nopcommerce-bench/docs/P4-TRAFFIC-PLAN.md},
\path{letsgolegacy.secondkey-nopcommerce-bench/results/p4/README.md},
\path{letsgolegacy.secondkey-nopcommerce-bench/scripts/traffic/NopCommerce.Scenarios.ps1}}

\subsubsection{Tryby awarii i~lista kontrolna planu ruchu}\label{sec:bench-ruch-awarie}

\begin{center}
\small
\begin{tabularx}{\linewidth}{@{}>{\raggedright\arraybackslash}X >{\raggedright\arraybackslash}X@{}}
\toprule
Objaw & Przyczyna \\
\midrule
Każda kwota podatku i~wysyłki w~nagraniu to \$0.00 & pominięto konfigurację (sekcja~2
  planu): dane przykładowe nie mają stawek \\
Kupon przechodzi na legacy, a~u~kandydata jest „not found” & wyszukiwanie rabatu
  kandydata przeszło z~SQL (kolacja, bez rozróżniania wielkości liter) do pamięci
  (porównanie porządkowe, ordinal) \\
Odtworzenia rozjeżdżają się od pierwszego scenariusza & odtworzenie nie wystartowało ze
  snapshotu S11 albo nagranie ominęło proxy \\
Reguła \texttt{never} przechodzi, niczego nie sprawdzając & scenariusz nie dotarł do tego,
  czego reguła pilnuje (np. T30 bez prawdziwego numeru zamówienia): każda reguła
  \texttt{never} potrzebuje pozytywnego bliźniaka w~tym samym scenariuszu -- R44 z~R43,
  R24 z~zarejestrowaną połową T18, R39 z~poprawną kartą w~T28 \\
Numery zamówień rozjeżdżają się między legacy a~kandydatem & scenariusze biegły
  z~różnych snapshotów albo dodatkowe żądanie utworzyło zamówienie \\
\bottomrule
\end{tabularx}
\end{center}

\begin{checklista}
  \item Konfiguracja z~sekcji~2 zastosowana przez panel administracyjny, potem zrobiony
    snapshot S11.
  \item Każde żądanie każdego scenariusza przeszło przez proxy nagrywające.
  \item Każda persona używała własnego słoika ciasteczek; tokeny anti-forgery czytane ze
    strony tam, gdzie są sprawdzane.
  \item Każda reguła z~sekcji~5 jest ćwiczona przez nagrany scenariusz, a~każda reguła
    \texttt{never} ma pozytywnego bliźniaka.
  \item Liczba scenariuszy i~skrót manifestu witryny legacy (\path{state\build.json}) są
    zapisane razem z~\path{*.skcap}.
  \item Sondy NLS $\rightarrow$ ICU z~sekcji~6 są w~nagraniu.
\end{checklista}

\zrodla{\path{letsgolegacy.secondkey-nopcommerce-bench/docs/P4-TRAFFIC-PLAN.md}}

\subsection{Kontrakt i~dowód A/A (ADR 0006, P5)}\label{sec:bench-kontrakt}

P5 zamienia sekcję~5 planu ruchu w~kontrakt Second Key (sekcja~\ref{sec:formaty}).
Zadanie jest zrobione, gdy kontrakt przechodzi walidację (\texttt{sk validate}) i~jest
spełniony na systemie legacy. Kandydata jeszcze nie ma (P6), więc „spełniony na legacy”
trzeba pokazać samym sklepem legacy, a~kontrakt musi być napisany tak, by P7 mogło
rozliczyć z~niego zmigrowany sklep. \path{contract/contract.yaml} opisuje, co sklep musi
robić i~czego nigdy nie może robić na zestawie ruchu P4, jako klauzule nad nagranymi
żądaniami; \path{contract/secondkey.yaml} opisuje, jak ruch jest odtwarzany.

\zrodla{\path{letsgolegacy.secondkey-nopcommerce-bench/docs/adr/0006-nopcommerce-contract-and-aa-proof.md},
\path{letsgolegacy.secondkey-nopcommerce-bench/contract/README.md}}

\subsubsection{Reguły planu i~sposób pisania klauzul}\label{sec:bench-kontrakt-reguly}

Sekcja~5 planu ruchu ma 46~reguł (R01--R46), każdą z~cytatem kodu 3.90, z~którego
pochodzi: 37 rodzaju \texttt{must} i~9 rodzaju \texttt{never} (R05, R19, R20, R23, R24,
R32, R39, R44, R45), czyli 19,6\%. Reguła \texttt{must} stwierdza, co zawiera odpowiedź
albo stan po niej; reguła \texttt{never} -- czego nie ma. „Pokazane” znaczy: w~HTML strony
wskazanej w~scenariuszu; fakty o~zamówieniach można też stwierdzać na delcie bazy (zadanie
rdzenia S4). Kontrakt pisze się według sześciu zasad:
\begin{itemize}
  \item \textbf{Z~reguł planu.} Każda klauzula w~polu \texttt{why} podaje reguły, które
    wyraża; tabela~\ref{tab:bench-mapa} mapuje reguły na klauzule.
  \item \textbf{Tytuł mówi dokładnie, co sprawdzają asercje} (ADR 0006, decyzja~4):
    „Billed to US / New York / 10001, the \$500.00 Lenovo with Ground shows tax 45.00
    (9\%)”, a~nie „tax is correct”. Werdykt i~pakiet drukują tytuły, więc czytelnik musi
    wiedzieć, co zostało spełnione, bez otwierania YAML.
  \item \textbf{Zakres przez sondę.} Klauzula wybiera żądanie scenariusza przez
    \texttt{when: \{ query: \{ probe: "\textasciicircum{}T11\$" \} \}}
    (sekcja~\ref{sec:bench-ruch-nagranie}).
  \item \textbf{HTML czytany przez ekstraktory.} Strona jest porównywana i~sprawdzana tylko
    przez wartości wyciągnięte ekstraktorami (\texttt{compare.html: extracts}): wiersze
    sum koszyka, komunikaty kuponów, tytuły list, sumy zamówień. Znaczniki nie są
    zachowaniem.
  \item \textbf{Każda klauzula \texttt{never} ma pozytywnego bliźniaka} w~tym samym
    scenariuszu, żeby nieobecność nie była spełniona dlatego, że nic się nie wydarzyło:
    odrzucony \texttt{MEMBERS15} gościa obok \texttt{(\$75.00)} klienta zarejestrowanego,
    odrzucona karta obok złożonego zamówienia, odrzucone strony zamówień obok własnych.
  \item \textbf{Zachowanie legacy przypięte takie, jakie jest} (ADR 0006, decyzja~2),
    łącznie z~tym, co migracja mogłaby chcieć „naprawić”: rabaty zapisywane jako
    \texttt{(\$100.00)}, pusty komunikat gościa dla \texttt{MEMBERS15}, opłata \$1.99 za
    odbiór, której darmowa wysyłka nie znosi, odmowa dla \texttt{SAVE10}+U+180E. Kandydat,
    który zmieni jedno z~nich, łamie klauzulę, a~człowiek rozstrzyga, czy to regresja,
    czy poprawka -- różnica nigdy nie jest wchłaniana przez regułę normalizacji.
\end{itemize}

{\footnotesize
\begin{xltabular}{\linewidth}{@{}>{\raggedright\arraybackslash}p{3.9cm} >{\raggedright\arraybackslash}X@{}}
\caption{Reguły planu i~klauzule kontraktu (za \texttt{contract/README.md})}\label{tab:bench-mapa}\\
\toprule
Reguły planu & Klauzule \\
\midrule
\endfirsthead
\toprule
Reguły planu & Klauzule \\
\midrule
\endhead
\bottomrule
\endlastfoot
R01 strony sklepu & \path{STOREFRONT-ANSWERS-WITH-FOOTER}, \path{COMPUTERS-LISTS-ITS-SUBCATEGORIES} \\
R02 strona główna & \path{HOME-LISTS-THE-FOUR-HOME-PAGE-PRODUCTS} \\
R03, R04 sortowanie, stronicowanie & \path{NOTEBOOKS-BY-PRICE-ASCENDING}, \path{NOTEBOOKS-BY-NAME}, \path{NOTEBOOKS-PAGE-TWO-OF-THREE-PER-PAGE} \\
R05 nieopublikowane (\texttt{never}) & \path{UNPUBLISHED-NEVER-LISTED}, \path{UNPUBLISHED-PAGE-NEVER-SERVED}; bliźniak \path{UNPUBLISHED-PAGE-NOT-FOUND} \\
R06 strony produktów & \path{LENOVO-PAGE-NAME-SKU-PRICE}, \path{MACBOOK-PAGE-NAME-SKU-PRICE}, \path{COMPUTER-PAGE-ATTRIBUTE-ADJUSTMENTS} \\
R07, R08 wyszukiwanie & \path{SEARCH-FINDS-MACBOOK-BY-NAME-CASE-AND-SKU}, \path{SEARCH-TERM-MINIMUM-LENGTH} \\
R09--R14 koszyk & \path{LISTING-ADD-TO-CART-SUCCEEDS}, \path{T07-CART-HOLDS-ONE-LENOVO}, \path{QUANTITY-ABOVE-STOCK-REFUSED}, \path{QUANTITY-BELOW-MINIMUM-REFUSED}, \path{REQUIRED-ATTRIBUTE-REFUSED}, \path{T09-ATTRIBUTES-PRICED-INTO-UNIT-PRICE}, \path{T10-*} \\
R15, R17, R18, R21, R22, R25--R27 rabaty & \path{SAVE10-*}, \path{COUPON-TRIMMED-AND-CASE-INSENSITIVE}, \path{*-COUPON-MESSAGE}, \path{PCT20MAX100-*}, \path{SHIPFREE-*}, \path{ONCE5-*}, \path{COUPON-123-*}, \path{TAKE50-*}, \path{MEMBERS15-FOR-A-REGISTERED-CUSTOMER}, \path{MEMBERS15-GUEST-GETS-AN-EMPTY-MESSAGE} \\
R19, R20, R23, R24 rabaty (\texttt{never}) & \path{REFUSED-COUPON-NEVER-*}, \path{CART-TOTAL-NEVER-NEGATIVE}, \path{GIFT-CARD-CART-NEVER-DISCOUNTED}, \path{MEMBERS15-NEVER-FOR-A-GUEST}, \path{MEMBERS15-GUEST-NEVER-TOLD-APPLIED} \\
R28--R30 podatki & \path{TAX-US-*}, \path{TAX-NO-RATE-GERMANY}, \path{TAX-ZIP-MATCH-IGNORES-CASE}, \path{TAX-FOLLOWS-BILLING-ADDRESS} \\
R32 zwolnienie z~podatku (\texttt{never}) & \path{TAX-EXEMPT-NEVER-TAXED}, \path{TAX-EXEMPT-CART-PRICED} \\
R33--R35 wysyłka & \path{ESTIMATE-OFFERS-FIXED-RATES-AND-PICKUP}, \path{CHECKOUT-OFFERS-THE-SAME-RATES}, \path{FREE-SHIPPING-*}, \path{NO-FREE-SHIPPING-AT-EXACTLY-1000} \\
R31, R36--R38, R40, R41 zamówienia & \path{GUEST-ORDER-*}, \path{REGISTERED-ORDER-TOTALS}, \path{ORDER-EMPTIES-THE-CART}, \path{ORDER-LOWERS-STOCK} \\
R39 cyfra kontrolna karty (\texttt{never}) & \path{BAD-CARD-NEVER-REACHES-CONFIRM}, \path{BAD-CARD-NUMBER-REFUSED} \\
R42, R43 własne zamówienia & \path{HISTORY-LISTS-OWN-ORDERS-ONLY}, \path{OWN-ORDER-DETAILS-SHOWN}, \path{OWN-ORDER-PDF-INVOICE} \\
R44, R45 cudze zamówienie (\texttt{never}) & \path{OTHERS-ORDER-NEVER-SHOWN}, \path{OTHERS-ORDER-NEVER-NAMES-ITS-OWNER}, \path{OTHERS-ORDER-SENDS-TO-LOGIN}, \path{GUEST-NEVER-GETS-ORDER-HISTORY}, \path{NOBODY-REORDERS-OTHERS-ORDER} \\
R46 format waluty & \path{CART-UNIT-PRICES-WRITTEN-IN-USD}; kwoty rabatów przypięte jako \texttt{(\$100.00)} w~klauzulach rabatów \\
§6 sondy NLS $\rightarrow$ ICU & \path{COUPON-WITH-U180E-REFUSED}; kwoty w~stylu \texttt{(\$100.00)} powyżej \\
-- & \path{SIGN-IN-RETURNS-HOME}, \path{REGISTRATION-COMPLETES}, \path{NEVER-A-SERVER-ERROR} \\
\end{xltabular}
}

Dwie reguły nie mają własnej klauzuli. \textbf{R16} (rabat procentowy liczony
w~pojedynczej precyzji i~zaokrąglany „half to even”): nagrane kwoty nie odróżniają tego
obliczenia od dokładnej arytmetyki dziesiętnej, więc klauzule rabatów przypinają kwoty.
\textbf{R37} (sumy złożonego zamówienia równe sumom z~kroku potwierdzenia): jedna klauzula
widzi jedno żądanie, więc klauzule szczegółów zamówienia stwierdzają sumy, które pokazał
krok potwierdzenia.

\zrodla{\path{letsgolegacy.secondkey-nopcommerce-bench/contract/README.md},
\path{letsgolegacy.secondkey-nopcommerce-bench/docs/P4-TRAFFIC-PLAN.md},
\path{letsgolegacy.secondkey-nopcommerce-bench/docs/adr/0006-nopcommerce-contract-and-aa-proof.md}}

\subsubsection{Struktura pliku i~liczby}\label{sec:bench-kontrakt-struktura}

Plik zaczyna się od \texttt{apiVersion: secondkey/v1} i~\texttt{kind: Contract}; blok
\texttt{metadata} podaje nazwę \path{nopcommerce-390}, tytuł, opis systemu („nopCommerce
release-3.90 (12d1f01), IIS, .NET Framework 4.8, SQL Server 2022 Express”), rewizję~1
i~właściciela \texttt{bench}. Sekcja \texttt{compare} porównuje
nagłówki \texttt{content-type} i~\texttt{location}, a~HTML -- przez ekstrakty
(\texttt{html: extracts}). Walidacja: \texttt{sk validate contract/contract.yaml}.
Kontrakt ma \textbf{76~klauzul -- 61 rodzaju \texttt{must} i~15 rodzaju \texttt{never}}
(udział nieobecności 19,7\%) -- nad 32~ekstraktorami i~z~trzema maskami normalizacji;
wszystkie klauzule są zaakceptowane. Każda ma proweniencję \texttt{origin: designed},
\texttt{status: accepted}, \texttt{author: bench}, \texttt{date: 2026-09-26}. Pełną listę
klauzul z~zakresami i~tytułami podaje dodatek~\ref{app:klauzule}.

Policzone z~pliku: klauzule mają łącznie 134~asercje. Operatory: \texttt{equals} (86),
\texttt{contains} (21), \texttt{matches} (13), \texttt{countEquals} (4),
\texttt{absent} (4), \texttt{exists} (3), \texttt{lt}, \texttt{gt} i~\texttt{gte} (po
1); kwantyfikator \texttt{any} występuje w~9~asercjach (wiersze sum zamówień). Selektory
to \path{response.status}, \path{response.body.json.<pole>},
\path{response.body.text}, \path{response.headers.<nazwa>[0]}
i~\path{extract.<nazwa>}; gdy operator wartości działa na ekstraktorze
z~\texttt{all: true}, selektor kończy się \texttt{[*]} (lekcja rozgrzewki,
sekcja~\ref{sec:bench-wyniki-p3}). Jedyna
klauzula bez \texttt{when}, \path{NEVER-A-SERVER-ERROR}, obejmuje każde żądanie zestawu
(\texttt{response.status gte 500}).

Ekstraktory (29 z~HTML, 3 z~JSON) grupują się według stron:

{\small
\begin{xltabular}{\linewidth}{@{}>{\raggedright\arraybackslash}p{3.4cm} >{\raggedright\arraybackslash}X@{}}
\toprule
Zakres \texttt{when} & Ekstraktory \\
\midrule
\endfirsthead
\toprule
Zakres \texttt{when} & Ekstraktory \\
\midrule
\endhead
\bottomrule
\endlastfoot
\texttt{GET} strony sklepu & \texttt{footer} (\texttt{.footer-powered-by} na dziewięciu
  stronach, które otwiera ruch); \texttt{homeProducts} (\path{^/$}),
  \texttt{listingProducts} (\path{^/(notebooks|search)$}), \texttt{subcategories}
  (\path{^/computers$}), \texttt{searchWarning} (\path{^/search$}); dla trzech stron
  produktów \texttt{productName}, \texttt{productSku}, \texttt{productPrice} (atrybut
  \texttt{content} elementu \texttt{span[itemprop=price]}, \texttt{as: decimal});
  \texttt{attributeOptions} (\path{^/build-your-own-computer$}) \\
\path{^/cart$}, każda metoda & \texttt{cartProducts}, \texttt{cartQuantities} (wartość
  pola \texttt{qty-input}, \texttt{integer}), \texttt{cartUnitPrices},
  \texttt{cartLineTotals}, \texttt{cartSubtotal}, \texttt{cartShipping},
  \texttt{cartTax}, \texttt{cartTotal} (\texttt{decimal}, \texttt{culture: en-US});
  \texttt{cartUnitPriceTexts}, \texttt{cartSubtotalDiscount},
  \texttt{cartTotalDiscount} (tekst, bez konwersji -- stąd przypięta forma
  \texttt{(\$100.00)}); \texttt{couponMessages} (\texttt{.message-success},
  \texttt{.message-failure} w~\texttt{.coupon-box}); \texttt{cartEmpty} \\
\texttt{POST} wysyłka & \texttt{shippingEstimates} (\path{^/cart/estimateshipping$},
  \texttt{.shipping-results .option-name}); \texttt{checkoutGround},
  \texttt{checkoutSecondDayAir}, \texttt{checkoutNextDayAir}
  (\path{^/checkout/OpcSaveBilling/$} z~sondą \texttt{T24}; \texttt{from: json},
  \texttt{selector: update\_section.html}, wyrażenie regularne na etykiecie metody) \\
\texttt{GET} zamówienia i~klienci & \texttt{placedOrderNumber}
  (\path{^/checkout/completed/$}, liczba z~\texttt{.order-number strong});
  \texttt{orderNumber}, \texttt{orderStatus}, \texttt{orderTotals} (wiersze
  \texttt{table.cart-total}) dla \path{^/orderdetails/[0-9]+$}; \texttt{historyOrders}
  (\path{^/order/history$}); \texttt{registrationResult} (\path{^/registerresult/1$}) \\
\end{xltabular}
}

\zrodla{\path{letsgolegacy.secondkey-nopcommerce-bench/contract/contract.yaml},
\path{letsgolegacy.secondkey-nopcommerce-bench/contract/README.md}}

\subsubsection{Normalizacja: trzy maski}\label{sec:bench-kontrakt-maski}

Każda maska dotyczy wartości, której sklep nie wyprowadza z~samego żądania i~snapshotu,
więc dwie odpowiedzi na to samo żądanie mogą się w~niej różnić (ADR 0006, decyzja~3):

\begin{lstlisting}
normalize:
  masks:
    patterns:
      - { name: pdf-bytes, regex: '^JVBERi0[A-Za-z0-9+/=]*$', replacement: '<pdf>' }
      - { name: anti-forgery-token, regex: '(?<=name="__RequestVerificationToken" type="hidden" value=")[^"]+', replacement: '<token>' }
      - { name: cart-line-picture, regex: '(?<=\b(?:alt="Picture of |title="Show details for ))[^"]*|(?<=\bsrc=")[^"]*/content/images/thumbs/[^"]*', replacement: '<cached picture>' }
\end{lstlisting}

\begin{itemize}
  \item \texttt{pdf-bytes} -- bajty faktury PDF. Każde renderowanie zawiera czas
    utworzenia i~losowy identyfikator dokumentu, więc dwa renderowania nigdy się nie
    zgadzają; maska obejmuje treści base64 zaczynające się od \texttt{\%PDF-}, a~kontrakt
    stwierdza status i~typ faktury.
  \item \texttt{anti-forgery-token} -- wartość tokenu anti-forgery, świeża w~każdym
    formularzu renderowanym przez MVC. Krok potwierdzenia checkoutu jednostronicowego
    odpowiada JSON-em, którego \texttt{update\_section.html} zawiera podsumowanie
    zamówienia -- formularz z~takim tokenem.
  \item \texttt{cart-line-picture} -- adres, tekst \texttt{alt} i~tytuł linku obrazka
    pozycji koszyka. Wersja 3.90 buforuje je przez trzy minuty według identyfikatora
    pozycji koszyka (\path{ShoppingCartModelFactory.PrepareCartItemPictureModel}); każdy
    scenariusz startuje z~tego samego snapshotu, więc jego pierwsza pozycja zawsze dostaje
    ten sam identyfikator i~przez te trzy minuty pokazuje obrazek produktu, który ostatnio
    miał ten identyfikator -- przebieg A/A zobaczył pozycję Lenovo z~obrazkiem karty
    podarunkowej \$25. Widać to w~minikoszyku odpowiedzi na dodanie do koszyka
    i~w~podsumowaniu zamówienia kroku potwierdzenia. Ten bufor to stan, którego
    przywrócenie bazy nie resetuje. Czyszczenie go przed każdym scenariuszem też by go
    zresetowało, ale tylko przez panel administracyjny albo restart aplikacji, przed
    każdym scenariuszem po obu stronach -- także u~kandydata P7 -- więc kontrakt maskuje
    te trzy wartości i~mówi to wprost.
\end{itemize}

JSON jest porównywany w~całości, łącznie z~HTML w~środku, więc maskowane są tylko te
wartości: reszta sekcji -- nazwy produktów, ceny, ilości, sumy -- nadal się liczy. Strony
są porównywane przez ekstrakty, a~żaden ekstraktor nie czyta tokenu ani obrazka. Nic
więcej nie jest normalizowane: z~jednego snapshotu i~tych samych żądań sklep legacy
odpowiada tak samo po obu stronach odtworzenia A/A, łącznie z~numerami zamówień. Różnica,
której nic nie wyjaśnia, jest regresją, więc zaliczony przebieg A/A pokazuje też, że te
trzy maski to wszystko, czego porównanie potrzebuje.

Plan ruchu (sekcja~7) wymieniał jako wartości do normalizacji przez komparator (zadanie
rdzenia S12) ciasteczko \texttt{Nop.customer} i~identyfikatory GUID klientów, tokeny
anti-forgery oraz ciasteczka \texttt{ASP.NET\_SessionId} i~\texttt{NOPCOMMERCE.AUTH},
identyfikatory GUID zamówień i~znaczniki czasu (daty zamówień, „created on”, nagłówek
\texttt{Date}) oraz nagłówki transportowe w~rodzaju \texttt{Content-Length}; numerów
zamówień nie normalizuje się (\texttt{CustomOrderNumberMask = "\{ID\}"}). Kontrakt nie ma
masek na te wartości poza tokenem: porównuje tylko nagłówki \texttt{content-type}
i~\texttt{location}, a~HTML -- wyłącznie przez ekstrakty.

\zrodla{\path{letsgolegacy.secondkey-nopcommerce-bench/contract/contract.yaml},
\path{letsgolegacy.secondkey-nopcommerce-bench/contract/README.md},
\path{letsgolegacy.secondkey-nopcommerce-bench/docs/adr/0006-nopcommerce-contract-and-aa-proof.md},
\path{letsgolegacy.secondkey-nopcommerce-bench/docs/P4-TRAFFIC-PLAN.md}}

\subsubsection{Konfiguracja odtwarzania i~dowód A/A}\label{sec:bench-kontrakt-aa}

\path{contract/secondkey.yaml} trzyma to, czego nie da się powiedzieć w~wierszu poleceń
(nagranie, przebieg i~oba adresy bazowe podaje \path{scripts/8-replay.ps1}); ścieżki są
względne wobec pliku:

\begin{lstlisting}
version: 1
replay:
  scenarioMode: session
  # Generous: the first request of a scenario follows a database restore.
  timeoutSeconds: 120
  legacy:
    reset:
      sqlServerSnapshot: { connectionStringEnv: NOPBENCH_SK_SQL, database: nopcommerce_legacy, snapshot: nopcommerce_legacy_snapshot }
  candidate:
    reset:
      sqlServerSnapshot: { connectionStringEnv: NOPBENCH_SK_SQL, database: nopcommerce_legacy, snapshot: nopcommerce_legacy_snapshot }
  # The cart, estimate-shipping and registration forms carry an anti-forgery token issued
  # with the page; the replay posts the one the page it just loaded issued.
  correlation:
    - name: csrf
      regex: 'name="__RequestVerificationToken" type="hidden" value="([^"]+)"'
      formField: __RequestVerificationToken
compare:
  contract: contract.yaml
evidence:
  contract: contract.yaml
\end{lstlisting}

Connection string resetu jest czytany ze zmiennej \texttt{NOPBENCH\_SK\_SQL}, którą
ustawia krok~8 (uwierzytelnianie Windows, bez sekretu).

\paragraph{Dowód: odtworzenie A/A na jednym sklepie i~jednej bazie (ADR 0006, 1).}
Krok~8 odtwarza nagranie P4 przez \texttt{sk replay}; legacy i~kandydat to oba
\url{http://localhost:8080/}, a~każda strona jest przywracana do snapshotu nagrania przed
każdym scenariuszem. \texttt{sk replay} kończy scenariusz po jednej stronie, zanim
zresetuje drugą, więc jedna baza obsługuje obie. Przebieg trafia do Linuksa, gdzie
krok~9 wymaga od \texttt{sk compare} wyniku \texttt{pass} z~każdą zaakceptowaną klauzulą
spełnioną i~żadną niewykonaną -- każda klauzula jest sprawdzana dwa razy, raz na stronę,
na odpowiedziach, których sklep legacy naprawdę udzielił. \finding{Kontrakt, który nie
przechodzi A/A, myli się co do systemu legacy albo odtworzenie nie jest powtarzalne;
jedno i~drugie odebrałoby sens werdyktowi P7 -- dlatego A/A jest bramką.} P7 odtworzy to
samo nagranie na kandydacie z~\texttt{-Candidate}, z~tym samym resetem i~kontraktem;
klauzula spełniona tutaj, a~złamana tam, jest regresją z~konstrukcji.

\paragraph{Co znalazł pierwszy przebieg A/A.} Pierwszy przebieg tej bramki nie przeszedł,
a~oba powody leżały po stronie kontraktu, nie sklepu:
\begin{itemize}
  \item \textbf{Strona główna wymienia cztery produkty, nie sześć.} Dane przykładowe
    oznaczają „show on home page” cztery produkty; dwa kolejne, których oczekiwał
    kontrakt, są zakomentowane w~instalatorze 3.90 (\texttt{CodeFirstInstallationService}).
    Klauzula była złamana po obu stronach (\texttt{violated-both}) -- kontrakt mylił się co
    do systemu legacy, a~po to jest bramka.
    \path{HOME-LISTS-THE-FOUR-HOME-PAGE-PRODUCTS} wymienia teraz te cztery.
  \item \textbf{17~odpowiedzi JSON różniło się między dwoma odtworzeniami tego samego
    żądania.} 14 to podsumowanie zamówienia kroku potwierdzenia ze świeżym tokenem
    anti-forgery; 3 to minikoszyk odpowiedzi na dodanie do koszyka, w~którym pozycja
    Lenovo miała po jednej stronie obrazek karty podarunkowej \$25. Obie wartości są
    teraz maskowane (same wartości), a~drugi przebieg przeszedł.
\end{itemize}

\zrodla{\path{letsgolegacy.secondkey-nopcommerce-bench/contract/secondkey.yaml},
\path{letsgolegacy.secondkey-nopcommerce-bench/docs/adr/0006-nopcommerce-contract-and-aa-proof.md},
\path{letsgolegacy.secondkey-nopcommerce-bench/results/p5/README.md}}

\subsection{Skrypty: kroki 1--9, biblioteki i~rozgrzewka}\label{sec:bench-skrypty}

Ponumerowane skrypty w~\path{scripts/} to strona legacy benchu w~kolejności uruchamiania;
CI uruchamia je bez zmian, więc porażkę z~CI odtwarza się na dowolnej maszynie Windows tym
samym krokiem. Konwencje wspólne: Windows PowerShell 5.1 i~PowerShell 7 (CI -- 5.1, joby
linuksowe -- \texttt{pwsh}); wyłącznie ASCII, bo 5.1 czyta skrypt bez BOM jako ANSI;
\texttt{Set-StrictMode -Version 3.0} i~\texttt{\$ErrorActionPreference = 'Stop'}; każdy
skrypt dot-source'uje \path{lib/NopBench.Common.ps1} (kroki łańcucha także
\path{lib/NopBench.Chain.ps1}) i~działa samodzielnie; sam sprawdza swoje warunki wstępne;
nic nie zapisuje w~repozytorium; kończy się jawnym \texttt{exit 0}, a~porażka rzuca
wyjątek, więc kod ostatniego polecenia natywnego nie decyduje o~wyniku. W~GitHub Actions
skrypty grupują log (\texttt{::group::}), dopisują tabele do podsumowania joba
i~maskują sekrety. Przy blokadzie polityki wykonywania: \texttt{powershell
-ExecutionPolicy Bypass -File}. Każda wartość ma domyślną
(tabela~\ref{tab:bench-parametry}; gwiazdka oznacza parametry, których nie wymienia README
skryptów).

{\footnotesize
\begin{xltabular}{\linewidth}{@{}>{\raggedright\arraybackslash}p{3.9cm} l >{\raggedright\arraybackslash}p{4.1cm} >{\raggedright\arraybackslash}X@{}}
\caption{Parametry i~zmienne środowiskowe kroków 1--9}\label{tab:bench-parametry}\\
\toprule
Nazwa & Krok & Domyślnie & Znaczenie \\
\midrule
\endfirsthead
\toprule
Nazwa & Krok & Domyślnie & Znaczenie \\
\midrule
\endhead
\bottomrule
\endlastfoot
\texttt{-WorkRoot}, \texttt{NOPBENCH\_WORK} & wszystkie & \path{C:\nopbench}
  (\path{$HOME/nopbench} poza Windows) & katalog roboczy \\
\texttt{-Destination}; \texttt{-Force} & 1 & \path{<work>\legacy-src}; wyłączony & miejsce
  checkoutu; ponowne pobranie mimo zweryfikowanego checkoutu \\
\texttt{-SourceDir} & 2 & z~\path{fetch.json}, inaczej \path{<work>\legacy-src} &
  czysty checkout przypiętego commita \\
\texttt{-SiteDir} & 2 & \path{<work>\legacy-site} & cel publikacji (zastępowany) \\
\texttt{-MSBuild}; \texttt{-NuGet} & 2 & z~\texttt{PATH}, inaczej \texttt{vswhere};
  z~\texttt{PATH}, inaczej przypięty \texttt{NuGet.CommandLine} & narzędzia buildu \\
\texttt{-SkipIis}, \texttt{-SkipSql} & 3 & wyłączone & zostaw tę połowę (naprawy) \\
\texttt{-SiteDir} & 4 & z~\path{build.json}, inaczej \path{<work>\legacy-site} &
  witryna; sprawdzana z~\path{legacy-site.sha256} \\
\texttt{-InstallDir} & 4 & \path{C:\inetpub\nopcommerce-legacy} & katalog IIS
  (zastępowany) \\
\texttt{-SiteName} & 4 & \texttt{nopcommerce-legacy} & witryna, pula i~login SQL
  \texttt{IIS APPPOOL\textbackslash{}<nazwa>} \\
\texttt{-Port} & 4 & \texttt{8080} & port HTTP; zajęty port jest odrzucany \\
\texttt{-SqlServer} & 4 & z~\path{platform.json}, inaczej \texttt{.\textbackslash{}SQLEXPRESS} & instancja SQL
  Server \\
\texttt{-DatabaseName} & 4 & \texttt{nopcommerce\_legacy} & usuwana i~tworzona przy
  każdym przebiegu \\
\texttt{-AdminEmail} & 4 & \texttt{admin@nopbench.invalid} & e-mail administratora \\
\texttt{-Install\-TimeoutSeconds}* & 4 & \texttt{1200} & limit klienta HTTP instalatora \\
\texttt{-Warmup\-TimeoutSeconds}* & 4; 5 & \texttt{900}; \texttt{600} & czas na pierwszą
  odpowiedź sklepu \\
\path{NOPBENCH_ADMIN_PASSWORD} & 4 & generowane & znane hasło administratora
  (poza plikami repozytorium) \\
\texttt{-BaseUrl} & 5, 6 & z~\path{deploy.json}, inaczej \url{http://localhost:8080/} &
  sklep do testu lub konfiguracji \\
\texttt{-Out}; \texttt{-ProxyUrl} & 7 & \path{<work>\traffic\nopcommerce.skcap};
  \url{http://127.0.0.1:8000/} & nagranie; nasłuch proxy \\
\texttt{-Only} & 7 & wszystkie & wybrane scenariusze; nagranie częściowe, nie zestaw ruchu \\
\texttt{-Capture}; \texttt{-Candidate} & 8 & nagranie z~\path{traffic.json}; sklep legacy
  (A/A) & co odtworzyć; adres drugiej strony (P7) \\
\texttt{-Out}* & 8 & \path{<work>\replay\run.skrun} & przebieg \\
\texttt{-Run}, \texttt{-OutDir} & 9 & \path{<work>\replay\run.skrun}, \path{<work>\replay} &
  przebieg do oceny; miejsce werdyktu i~pakietu \\
\texttt{-Pdf} & 9 & \texttt{auto} & \texttt{sk evidence -{}-pdf}: \texttt{auto},
  \texttt{required} (bez przeglądarki z~Chromium -- błąd), \texttt{off} \\
\texttt{NOPBENCH\_SK\_SQL} & ustawia 8 & z~\path{deploy.json} & connection string resetu
  dla \repopath{sk}; bez sekretu \\
\texttt{NOPBENCH\_SK}, \texttt{NOPBENCH\_PORTCULLIS} & łańcuch &
  \path{<work>\tools\chain\secondkey\SecondKey.Cli.dll}, \path{<work>\tools\chain\portcullis\Portcullis.Cli.dll}
  & narzędzia łańcucha \\
\end{xltabular}
}

\zrodla{\path{letsgolegacy.secondkey-nopcommerce-bench/scripts/README.md},
\path{letsgolegacy.secondkey-nopcommerce-bench/scripts/lib/NopBench.Common.ps1},
\path{letsgolegacy.secondkey-nopcommerce-bench/scripts/4-deploy-and-install.ps1},
\path{letsgolegacy.secondkey-nopcommerce-bench/scripts/5-smoke.ps1},
\path{letsgolegacy.secondkey-nopcommerce-bench/scripts/8-replay.ps1}}

\subsubsection{Kroki 1--3: pobranie, build, platforma}\label{sec:bench-krok1}

\paragraph{\texttt{1-fetch-nopcommerce.ps1}} (git; Windows lub Linux; tworzy
\path{<work>\legacy-src} i~\path{fetch.json}). \texttt{git ls-remote -{}-tags} dla tagu
i~jego rozwinięcia \texttt{\textasciicircum\{\}} (do 3~prób): commit tagu musi równać się
przypięciu, inaczej krok odmawia budowania przesuniętego tagu. Zweryfikowany checkout
jest używany ponownie (chyba że \texttt{-Force}); w~przeciwnym razie \texttt{git init},
\texttt{core.autocrlf=false} (końce linii upstream bajt w~bajt, więc pliki statyczne nie
zależą od ustawień git maszyny), \texttt{core.longpaths=true}, \texttt{fetch -{}-depth 1
-{}-no-tags} samego tagu i~\texttt{checkout -{}-detach}. Na końcu \texttt{HEAD} musi równać
się przypięciu, a~\texttt{git status -{}-porcelain -{}-untracked-files=no} być pusty (pliki
nieśledzone -- wyjście buildu -- są dozwolone).

\paragraph{\texttt{2-build-legacy.ps1}} (Windows, Visual Studio 2022 lub Build Tools
2022 z~obciążeniem web, git; bez uprawnień administratora; tworzy
\path{<work>\legacy-site}, \path{build.json}, \path{legacy-site.sha256} oraz
\path{logs\build.binlog} i~\path{publish.binlog}). Odmawia budowania drzewa innego niż
czysty checkout pinu. Znajduje MSBuild i~nuget.exe, rozpakowuje przypięte asemblacje
referencyjne (sprawdza \path{mscorlib.dll}, \path{System.Web.dll},
\path{RedistList\FrameworkList.xml}) i~zapisuje, czy maszyna ma własny pakiet 4.5.1.
Potem \texttt{nuget restore} (upstream commituje pakiety, więc to raczej weryfikacja),
build rozwiązania i~publikacja \texttt{Nop.Web}:

\begin{lstlisting}
common:  /p:Configuration=Release /p:TargetFrameworkRootPath=<package>\build
         /p:ImportDirectoryBuildProps=false /p:ImportDirectoryBuildTargets=false
build:   /m /nr:false /v:minimal /t:Build "/p:Platform=Any CPU" /bl:<work>\logs\build.binlog
publish: /p:Platform=AnyCPU /p:BuildProjectReferences=false /p:DeployOnBuild=true
         /p:DeployDefaultTarget=WebPublish /p:WebPublishMethod=FileSystem
         /p:PublishUrl=<site> /p:DeleteExistingFiles=true
         /p:AutoParameterizationWebConfigConnectionStrings=false /bl:<work>\logs\publish.binlog
\end{lstlisting}

Kontrola witryny: obecne \path{Global.asax}, \path{Web.config}, asemblacje
\texttt{Nop.Web}, \texttt{Nop.Admin}, \texttt{Nop.Core}, \texttt{Nop.Data},
\texttt{Nop.Services}, \texttt{Nop.Web.Framework}, widoki sklepu i~Administration,
motyw \texttt{DefaultClean}, \path{App_Data\Install\SqlServer.StoredProcedures.sql},
zasoby lokalizacji i~\path{Plugins\bin}; nieobecne \path{App_Data\Settings.txt}
i~\path{InstalledPlugins.txt} (oznaczyłyby witrynę jako zainstalowaną), pliki
\path{.cs} i~\path{.csproj} oraz \texttt{debug="true"}; każdy projekt wtyczki jako katalog
z~\path{Description.txt} i~\path{Nop.Plugin.*.dll}. Drzewo źródeł jest sprawdzane
ponownie po buildzie. Manifest to wiersze \texttt{<SHA-256>~~<ścieżka>} (separator
\texttt{/}) sortowane porządkowo, niezależnie od kultury maszyny, zgodne z~\texttt{sha256sum
-c}.

\paragraph{\texttt{3-install-iis-sql.ps1}}\label{sec:bench-krok3} (uprawnienia
administratora; idempotentny; tworzy \path{platform.json}). IIS: \texttt{dism.exe /Online
/Enable-Feature /All /NoRestart} z~funkcjami \path{IIS-DefaultDocument},
\path{IIS-StaticContent}, \path{IIS-HttpErrors}, \path{IIS-HttpLogging},
\path{IIS-RequestFiltering}, \path{IIS-NetFxExtensibility45},
\path{IIS-ISAPIExtensions}, \path{IIS-ISAPIFilter}, \path{IIS-ASPNET45} (kod 3010
jest akceptowany), start \texttt{WAS} i~\texttt{W3SVC}, kontrola handlera
\path{ExtensionlessUrlHandler-Integrated-4.0}. SQL Server (gdy nie ma usługi
\texttt{MSSQL\$SQLEXPRESS}): pobranie z~weryfikacją SHA-256, podpis Authenticode ze
statusem \texttt{Valid} i~\texttt{O=Microsoft Corporation}, rozpakowanie, test sektorów
(\texttt{fsutil fsinfo sectorinfo}; ponad 4096~B na dysku systemowym przenosi dane do
\path{<work>\sqldata} albo kończy krok błędem ze wskazówką
\texttt{ForcedPhysicalSectorSizeInBytes}) i~instalator (dwie próby; po nieudanej -- 60
ostatnich linii \texttt{ERRORLOG}, zdarzenia \texttt{MSSQL*}, deinstalacja i~sprzątanie):

\begin{lstlisting}
/Q /ACTION=Install /IACCEPTSQLSERVERLICENSETERMS /FEATURES=SQLENGINE
/INSTANCENAME=SQLEXPRESS /INSTANCEID=SQLEXPRESS /SQLSYSADMINACCOUNTS="BUILTIN\Administrators"
/SQLCOLLATION=SQL_Latin1_General_CP1_CI_AS /TCPENABLED=1 /NPENABLED=0 /UPDATEENABLED=False
/SKIPRULES=RebootRequiredCheck [/INSTALLSQLDATADIR="<work>\sqldata"]
\end{lstlisting}

Istniejąca instancja jest używana ponownie; inna kolacja serwera daje tylko ostrzeżenie,
bo krok~4 tworzy bazę z~przypiętą kolacją jawnie.

\zrodla{\path{letsgolegacy.secondkey-nopcommerce-bench/scripts/1-fetch-nopcommerce.ps1},
\path{letsgolegacy.secondkey-nopcommerce-bench/scripts/2-build-legacy.ps1},
\path{letsgolegacy.secondkey-nopcommerce-bench/scripts/3-install-iis-sql.ps1}}

\subsubsection{Kroki 4--5: wdrożenie, instalator, smoke test}\label{sec:bench-krok4}

\paragraph{\texttt{4-deploy-and-install.ps1}} (uprawnienia administratora, po krokach 2
i~3). Kolejno: kontrola każdego wiersza manifestu i~liczby plików (brak manifestu --
ostrzeżenie i~wdrożenie bez weryfikacji); usunięcie witryny, puli i~bazy
(\texttt{SINGLE\_USER WITH ROLLBACK IMMEDIATE}, \texttt{DROP DATABASE}) oraz do 30~s
oczekiwania na zwolnienie portu; pula .NET CLR v4.0 w~potoku zintegrowanym
z~\texttt{ApplicationPoolIdentity}, \texttt{loadUserProfile=true}, zerowym
\texttt{idleTimeout} i~\texttt{periodicRestart.time},
\texttt{enable32BitAppOnWin64=false}; \texttt{icacls} \texttt{(OI)(CI)M} dla tożsamości
puli \emph{przed} kopiowaniem (pliki dziedziczą), \texttt{robocopy /MIR} (kody poniżej 8
to sukces); witryna z~powiązaniem \texttt{http/*:<port>:}; login \texttt{CREATE LOGIN …
FROM WINDOWS} z~rolą \texttt{dbcreator}. Instalator: oczekiwanie na formularz
\texttt{GET /install} (z~polem \texttt{AdminEmail}) i~wysłanie pól \texttt{InstallModel}:

\begin{lstlisting}
AdminEmail=<-AdminEmail>  AdminPassword=<password>  ConfirmPassword=<password>
InstallSampleData=true  DataProvider=sqlserver  SqlConnectionInfo=sqlconnectioninfo_raw
DatabaseConnectionString=Data Source=<server>;Initial Catalog=<database>;Integrated Security=True;Persist Security Info=False
SqlAuthenticationType=windowsauthentication  SqlServerCreateDatabase=true
UseCustomCollation=true  Collation=SQL_Latin1_General_CP1_CI_AS
\end{lstlisting}

Sukcesem jest tylko 302 na \texttt{/}; inaczej krok drukuje błędy walidacji formularza
i~kończy się błędem. Potem strona główna musi zawierać „nopCommerce”, a~plik
\path{App_Data\Settings.txt} istnieć; krok odczytuje liczbę tabel, produktów
(\texttt{Deleted = 0}) i~kolację bazy do \path{deploy.json}, zapisuje plik danych
administratora (sekcja~\ref{sec:bench-stan}) i~podsumowanie.

\paragraph{\texttt{5-smoke.ps1}} (dowolna maszyna, która sięga sklepu; bez uprawnień).
Jeden gość, jeden słoik ciasteczek; najpierw oczekiwanie na stronę główną (pierwsze
żądanie kompiluje widoki i~ładuje wtyczki), potem pięć kontroli: \emph{home}
(\texttt{GET /}, „Welcome to our store”), \emph{category} (\texttt{GET /desktops},
\texttt{<h1>Desktops</h1>} i~Lenovo), \emph{product}
(\texttt{GET} \path{/lenovo-ideacentre-600-all-in-one-pc}: nazwa, SKU \texttt{LE\_IC\_600},
\texttt{\$500.00}; identyfikator produktu ze strony), \emph{add}
(\texttt{POST /addproducttocart/catalog/\{id\}/1/1}, \texttt{"success":true}) i~\emph{cart}
(\texttt{GET /cart}: \texttt{<h1>Shopping cart</h1>} i~produkt, bez „Your Shopping Cart is
empty!”). Każda strona HTML musi mieć stopkę „Powered by nopCommerce” -- dowód, że
odpowiedział nopCommerce, a~licencja wymaga jej na każdej stronie. Kontrole są
bezwarunkowe i~nigdy pomijane; wynik każdej trafia do \path{smoke.json} i~podsumowania,
a~krok zawodzi na końcu, jeśli zawiodła którakolwiek.

\zrodla{\path{letsgolegacy.secondkey-nopcommerce-bench/scripts/4-deploy-and-install.ps1},
\path{letsgolegacy.secondkey-nopcommerce-bench/scripts/5-smoke.ps1}}

\subsubsection{Krok 6: konfiguracja sklepu}\label{sec:bench-krok6}

\path{6-configure-store.ps1} wymaga pliku danych administratora i~\path{deploy.json}.
Sesja administratora idzie wprost do witryny z~nagłówkiem \texttt{x-bench-configuration}
(nie nagłówkiem scenariusza); każdy \texttt{POST} panelu niesie token anti-forgery ze
strony. Podatki: strona
\path{/Admin/Tax/ConfigureProvider?systemName=Tax.FixedOrByCountryStateZip}, adresy
\texttt{SaveMode}, \texttt{AddRateByCountryStateZip}
i~\texttt{RatesByCountryStateZipList} czytane ze skryptu strony, identyfikatory kategorii,
krajów i~stanów z~list wyboru; po dodaniu lista stawek musi mieć dokładnie pięć
oczekiwanych wierszy. Katalog: najpierw dowód, że adres HP Envy serwuje ten produkt
(inaczej późniejsze 404 niczego by nie dowodziło); \texttt{IsTaxExempt=true} dla MacBooka
przez formularz produktu, \texttt{Published=false} (HP Envy) i~\texttt{StockQuantity=3}
(Lenovo) przez edycję zbiorczą z~odczytem; \texttt{AllowViewUnpublishedProductPage=false},
po czym zwykły gość musi dostać 404. Wysyłka: stawki przez
\texttt{UpdateFixedShippingRate} i~ustawienia \texttt{FreeShippingOverXEnabled=true},
\texttt{FreeShippingOverXValue=1000}, \texttt{FreeShippingOverXIncludingTax=false}.
Rabaty: formularze \path{/Admin/Discount/Create} z~odczytem pól; wymaganie roli dla
\texttt{MEMBERS15} dodaje strona konfiguracji wtyczki
(\path{/Plugins/DiscountRulesCustomerRoles/Configure}). Zadania w~tle: \texttt{TaskUpdate}
z~\texttt{Enabled=false}, \texttt{POST} na \path{/Admin/Common/RestartApplication}
i~oczekiwanie do 300~s na nowy wpis „Application
started” w~tabeli \texttt{dbo.Log} -- sama odpowiedź nie dowodzi restartu. Porównanie
formularzy przed i~po zapisie pomija token i~pola \texttt{downloadurl<liczba>} generowane
przy każdym renderowaniu; w~ustawieniach wysyłki dopuszcza też zmianę
\texttt{ShippingOriginAddress.Id}, bo pierwszy zapis tworzy pusty rekord adresu nadania.
Nieudany obszar daje adnotację \texttt{::error::}, następne obszary i~tak działają, a~krok
kończy się błędem z~ich listą. Na sklepie już skonfigurowanym krok zawodzi, więc ponowna
konfiguracja wymaga kroku~4. Wynik: \path{store-config.json}.

\zrodla{\path{letsgolegacy.secondkey-nopcommerce-bench/scripts/6-configure-store.ps1},
\path{letsgolegacy.secondkey-nopcommerce-bench/scripts/README.md}}

\subsubsection{Kroki 7--9: nagranie, odtworzenie, werdykt}\label{sec:bench-krok7}

\paragraph{\texttt{7-record-traffic.ps1}} (host sklepu; wymaga \path{deploy.json},
\path{store-config.json} -- bez niego „każda kwota podatku i~wysyłki to \$0.00” --
i~\repopath{sk}). Robi snapshot (istniejący jest używany ponownie), startuje proxy i~czeka
do 60~s na wiersz \texttt{sk capture: recording} w~\path{logs\traffic-capture.log}:

\begin{lstlisting}
dotnet <SecondKey.Cli.dll> capture --listen http://127.0.0.1:8000 --target http://localhost:8080/ --out <work>\traffic\nopcommerce.skcap --session-key x-bench-scenario
\end{lstlisting}

Przed każdym scenariuszem przywraca snapshot (\texttt{SINGLE\_USER}, \texttt{RESTORE
DATABASE … FROM DATABASE\_SNAPSHOT}, \texttt{MULTI\_USER}; do 3~prób) i~odpytuje witrynę
wprost do pierwszego 200 (do 120~s; wcześniejsze odpowiedzi dają ostrzeżenie), potem
uruchamia scenariusz w~nowej sesji przez proxy. Porażka scenariusza daje
\texttt{::error::} i~pięć najnowszych wpisów \texttt{dbo.Log}; kolejne scenariusze
działają dalej. Proxy jest zatrzymywane po liczbie wymian, nagranie przechodzi
\texttt{sk validate}, a~\path{traffic.json} i~podsumowanie zawierają tabelę scenariuszy
z~identyfikatorami sesji.

\paragraph{\texttt{8-replay.ps1}} (host sklepu; \repopath{sk}; bez \texttt{-Capture}
wymaga \path{traffic.json}, które nie opisuje nagrania częściowego). Tryb „A/A”, gdy
kandydat równa się adresowi legacy, inaczej „legacy vs candidate”. Krok ustawia dla
procesu \repopath{sk} \texttt{NOPBENCH\_SK\_SQL} (\texttt{Data Source=<serwer>;Integrated
Security=True;Encrypt=False}), uruchamia z~korzenia repozytorium poniższe polecenie
i~usuwa zmienną:

\begin{lstlisting}
dotnet <SecondKey.Cli.dll> replay --config contract\secondkey.yaml --capture <skcap> --out <work>\replay\run.skrun --legacy <legacy> --candidate <candidate>
\end{lstlisting}

Następnie waliduje przebieg i~czyta z~niego wiersze \texttt{exchange.result} (wyniki,
błędy, statusy) i~\texttt{state.reset} (scenariusze, nieudane resety, czasy). Wynik bez
odpowiedzi, nieudany reset albo liczba scenariuszy inna niż w~\path{traffic.json} kończy
krok błędem („to nie jest uczciwe porównanie”). Wynik: \path{replay.json}.

\paragraph{\texttt{9-verdict.ps1}} (dowolna maszyna z~\repopath{sk}; w~CI Linux).
\texttt{sk validate} kontraktu i~przebiegu; \texttt{sk compare} do \path{verdict.json}
(akceptowane kody 0, 1, 2 -- ich znaczenie w~sekcji~\ref{sec:cli}) i~jego walidacja;
\texttt{sk evidence} z~\texttt{-{}-pdf} do pustego katalogu \path{replay\evidence}, po czym
każdy plik jest sprawdzany z~SHA-256 w~\path{manifest.json}, a~pakiet musi zawierać
\path{index.html}, \path{verdict.json}, \path{contract.yaml}, \path{run.skrun}
i~\path{statement.intoto.json}. Krok liczy, które maski wyjaśniły wymiany
\texttt{equal-under-contract} (\texttt{normalizedBy}), a~przy porażce loguje każdą
niespełnioną klauzulę z~maksymalnie trzema wymianami i~tekstem odpowiedzi legacy oraz do
dziesięciu wymian \texttt{regression}/\texttt{fix-candidate} z~różnicami.
\textbf{Bramka:} w~trybie A/A (także przy braku \path{replay.json}) krok zawodzi -- po
zapisaniu pakietu, który wyjaśnia porażkę -- gdy wynik nie jest \texttt{pass} albo
jakakolwiek zaakceptowana klauzula nie jest spełniona lub jest niewykonana. W~trybie
„legacy vs candidate” zapisuje werdykt i~pakiet bez tej bramki. Wynik:
\path{state\verdict.json}.

\paragraph{\texttt{collect-diagnostics.ps1}} nie jest krokiem: zbiera po porażce kroku
3, 4 albo 5 pliki przekazania, konfigurację IIS, logi IIS i~\texttt{HTTPERR}, zdarzenia
z~ostatnich 3~godzin, podsumowanie instalatora i~\texttt{ERRORLOG} SQL Server, pliki
\path{App_Data} i~50 ostatnich wpisów tabeli \texttt{Log} do \path{<work>\diagnostics}
(\texttt{-OutputDir}). Działa „w~miarę możliwości” (pominięte sekcje trafiają do
\path{skipped.txt}), sam nigdy nie zawodzi i~nigdy nie zbiera \path{<work>\secrets}.

\zrodla{\path{letsgolegacy.secondkey-nopcommerce-bench/scripts/7-record-traffic.ps1},
\path{letsgolegacy.secondkey-nopcommerce-bench/scripts/8-replay.ps1},
\path{letsgolegacy.secondkey-nopcommerce-bench/scripts/9-verdict.ps1},
\path{letsgolegacy.secondkey-nopcommerce-bench/scripts/collect-diagnostics.ps1}}

\subsubsection{Biblioteki}\label{sec:bench-biblioteki}

\begin{itemize}
  \item \path{lib/NopBench.Common.ps1}: log, grupy, maski i~podsumowanie Actions; katalog
    roboczy, \path{pins.json} i~pliki przekazania; \texttt{Invoke-BenchNative} (lista
    akceptowanych kodów wyjścia; łagodzi \texttt{ErrorActionPreference}, bo 5.1 zamienia
    przekierowany stderr w~błędy); \texttt{Invoke-BenchGit} z~ponawianiem;
    \texttt{Test-BenchPinnedCheckout}; \texttt{Get-BenchVerifiedDownload} (plik z~pasującym
    SHA-256 używany ponownie, pobranie do \path{.partial}, do 4~prób, niezgodny skrót
    usuwa plik -- nic niezweryfikowanego nie zostaje); pakiety NuGet; \texttt{appcmd};
    klient \texttt{System.Net.Http} z~ciasteczkami, bez automatycznych przekierowań
    i~z~User-Agent przeglądarki (zachowuje się tak samo w~5.1 i~7, w~przeciwieństwie do
    \texttt{Invoke-WebRequest}); SQL z~uwierzytelnianiem Windows; tworzenie
    i~przywracanie snapshotów (pliki \path{.ss} obok plików danych; snapshoty są w~każdym
    wydaniu SQL Server od 2016 SP1, także Express).
  \item \path{lib/NopBench.Chain.ps1}: narzędzia łańcucha (kody wyjścia \repopath{sk} są
    jego interfejsem, więc wywołujący wymienia akceptowane), proxy nagrywające, licznik
    wymian (wiersze \texttt{\{"type":"http.exchange"}), identyfikator sesji; sesja
    scenariusza z~licznikiem żądań liczonych także przy błędzie transportu (proxy mogło je
    nagrać); \texttt{Invoke-BenchStep} (status, treść obowiązkowa i~zakazana, typ treści,
    \texttt{Location}; niespełnienie rzuca wyjątek z~wyciągiem odpowiedzi); formularze
    jak w~przeglądarce (\texttt{Get-BenchForm}: pola w~kolejności dokumentu, zaznaczone
    pola wyboru, wybrane opcje, bez przycisków, plików i~pól wyłączonych).
  \item \path{lib/NopBench.Shop.ps1}: sklep 3.90 tak, jak używa go przeglądarka --
    dokładnie żądania stron i~skryptów sklepu (\path{public.ajaxcart.js},
    \path{public.onepagecheckout.js}): sondy, dodawanie do koszyka, formularz koszyka
    (wysyłany jako urlencoded, choć strona ma multipart -- akcje czytają go tak samo,
    a~korelacja odtwarzania sięga pola tokenu), kupony, szacowanie wysyłki, logowanie,
    rejestracja, wejście przyciskiem Checkout i~checkout jednostronicowy krok po kroku
    (domyślnie Ground i~\texttt{Payments.CheckMoneyOrder}; błędna karta wysyłana
    najpierw).
  \item \path{traffic/NopCommerce.Scenarios.ps1}: 38~scenariuszy, stałe danych
    przykładowych (Lenovo -- produkt~3, MacBook -- 4, komputer -- 1, książka -- 39, karta
    podarunkowa -- 43, HP Envy -- 8), \texttt{Assert-NopTotal} (wiersze bloku sum jako
    wyrażenia regularne) i~\texttt{Get-NopNegative} (przy nagrywaniu obie formy kwoty
    ujemnej).
\end{itemize}

\zrodla{\path{letsgolegacy.secondkey-nopcommerce-bench/scripts/lib/NopBench.Common.ps1},
\path{letsgolegacy.secondkey-nopcommerce-bench/scripts/lib/NopBench.Chain.ps1},
\path{letsgolegacy.secondkey-nopcommerce-bench/scripts/lib/NopBench.Shop.ps1},
\path{letsgolegacy.secondkey-nopcommerce-bench/scripts/traffic/NopCommerce.Scenarios.ps1}}

\subsubsection{Skrypty rozgrzewki (\texttt{scripts/eshop})}\label{sec:bench-skrypty-eshop}

Te same konwencje; CI uruchamia je w~\path{warmup-eshop.yml}. IIS dostarcza
\path{scripts/3-install-iis-sql.ps1} z~\texttt{-SkipSql}.

{\footnotesize
\begin{xltabular}{\linewidth}{@{}>{\raggedright\arraybackslash}p{3.3cm} >{\raggedright\arraybackslash}p{3.9cm} >{\raggedright\arraybackslash}X@{}}
\toprule
Skrypt & Parametry (domyślne) & Co robi i~co tworzy \\
\midrule
\endfirsthead
\toprule
Skrypt & Parametry (domyślne) & Co robi i~co tworzy \\
\midrule
\endhead
\bottomrule
\endlastfoot
\path{1-build-eshop.ps1} & \texttt{-SiteDir} (\path{<work>\eshop-site}),
  \texttt{-MSBuild} & commit pobierany po identyfikatorze (pin bez tagu), czysty checkout,
  asemblacje 4.6.1 i~4.7.2 pod jednym korzeniem, \texttt{/t:Restore} i~publikacja,
  kontrola plików; \texttt{UseMockData=true} w~opublikowanym \path{Web.config};
  \path{eshop-build.json} \\
\path{2-deploy-eshop.ps1} & \texttt{-SiteDir}, \texttt{-SiteName} (\texttt{eshop-legacy}),
  \texttt{-Port} (8081), \texttt{-WarmupTimeoutSeconds} (600) & czyste wdrożenie jak
  krok~4 (Modify, bo aplikacja pisze log log4net); czeka na katalog 12~pozycji;
  \path{eshop-deploy.json} \\
\path{3-record-eshop.ps1} & \texttt{-Out} (\path{<work>\eshop\traffic.skcap}),
  \texttt{-ProxyUrl} (\url{http://127.0.0.1:8001/}) & reset przed każdym scenariuszem,
  E01--E07 przez proxy, zatrzymanie po liczbie wymian, \texttt{sk validate}; porażka
  scenariusza przerywa nagranie; \path{eshop-traffic.json} \\
\path{4-replay-eshop.ps1} & \texttt{-Capture}, \texttt{-Out} (\path{<work>\eshop\run.skrun})
  & \texttt{sk replay} A/A z~\path{warmup\eshop\secondkey.yaml}; błąd przy wyniku bez
  odpowiedzi lub nieudanym resecie; \path{eshop-replay.json} \\
\path{5-scan-eshop.ps1} & \texttt{-Out} (\path{<work>\eshop\portcullis.sarif}) &
  \texttt{Portcullis.Cli.dll scan} ze \texttt{-{}-provenance-repo} i~\texttt{-{}-sarif};
  raport JSON zapisany obok; kod 1 (bramka blokuje) to wynik, nie porażka; SARIF musi mieć
  wersję 2.1.0; \path{eshop-scan.json} \\
\path{6-verdict-eshop.ps1} & \texttt{-Run}, \texttt{-Sarif}, \texttt{-OutDir}
  (\path{<work>\eshop}), \texttt{-Pdf} (\texttt{auto}) & \texttt{sk validate},
  \texttt{sk compare} (kody 0--2), \texttt{sk gate -{}-sarif} (kody 0--1), \texttt{sk
  evidence} z~SARIF, kontrola manifestu; błąd, jeśli A/A nie przechodzi w~pełni;
  \path{eshop-verdict.json} \\
\path{Reset-EShop.ps1} & \texttt{-TimeoutSeconds} (300) & zatrzymanie i~start puli (po
  60~s na stan), potem oczekiwanie na „Showing 10 of 12 products” \\
\end{xltabular}
}

Narzędzia łańcucha da się zbudować ręcznie (jedynym wymaganiem jest .NET 10 SDK) albo
wskazać zbudowane gdzie indziej przez \texttt{NOPBENCH\_SK} i~\texttt{NOPBENCH\_PORTCULLIS}:

\begin{lstlisting}
git clone https://github.com/konradcinkusz/letsgolegacy.secondkey.git secondkey
git -C secondkey checkout 967c49c0005c4412180914f43437d75b43f148dc   # chain.secondKey.commit
dotnet build secondkey/src/SecondKey.Cli/SecondKey.Cli.csproj -c Release -o <work>/tools/chain/secondkey
git clone https://github.com/konradcinkusz/letsgolegacy.portcullis.git portcullis
git -C portcullis checkout 84e1925fe0d85cf23415f0d33c98590f86128722  # chain.portcullis.commit
dotnet build portcullis/src/Portcullis.Cli/Portcullis.Cli.csproj -c Release -o <work>/tools/chain/portcullis
\end{lstlisting}

\zrodla{\path{letsgolegacy.secondkey-nopcommerce-bench/scripts/eshop/1-build-eshop.ps1},
\path{letsgolegacy.secondkey-nopcommerce-bench/scripts/eshop/2-deploy-eshop.ps1},
\path{letsgolegacy.secondkey-nopcommerce-bench/scripts/eshop/3-record-eshop.ps1},
\path{letsgolegacy.secondkey-nopcommerce-bench/scripts/eshop/4-replay-eshop.ps1},
\path{letsgolegacy.secondkey-nopcommerce-bench/scripts/eshop/5-scan-eshop.ps1},
\path{letsgolegacy.secondkey-nopcommerce-bench/scripts/eshop/6-verdict-eshop.ps1},
\path{letsgolegacy.secondkey-nopcommerce-bench/scripts/eshop/Reset-EShop.ps1},
\path{letsgolegacy.secondkey-nopcommerce-bench/scripts/README.md}}

\subsection{Workflowy CI}\label{sec:bench-ci}

Repozytorium ma cztery workflowy. Wszystkie mają \texttt{permissions: contents: read}
i~przypinają akcje przez SHA commita z~wersją w~komentarzu (Dependabot aktualizuje SHA
i~komentarz razem, jednym cotygodniowym pull requestem). Checkouty w~workflowach łańcucha
używają \texttt{persist-credentials: false}, a~checkout w~\path{secret-scan.yml} pobiera
całą historię (\texttt{fetch-depth: 0}); \path{legacy-build.yml}
i~\path{warmup-eshop.yml} mają \texttt{concurrency} na workflow i~numer PR (albo ref) z~\texttt{cancel-in-progress}.
Joby, które uruchamiają skrypty, ustawiają przez \texttt{GITHUB\_ENV} zmienną
\texttt{NOPBENCH\_WORK} na \path{$env:RUNNER_TEMP\nopbench} (Windows) albo
\path{$env:RUNNER_TEMP/nopbench} (Linux).
Rozkład jobów łańcucha (TOPOLOGY):

\begin{lstlisting}
legacy-build.yml
  lint  (ubuntu-24.04)   PSScriptAnalyzer on scripts/
  build (windows-2022)   1-fetch -> 2-build          -> artifact legacy-site (site + manifest)
  tools (ubuntu-24.04)   chain-tools.yml: sk and Portcullis at their pinned commits -> chain-tools
  iis   (windows-2022)   needs build, tools: 3-install-iis-sql -> 4-deploy-and-install -> 5-smoke
                         -> 6-configure-store -> 7-record-traffic (sk capture) -> nopcommerce-traffic
                         -> 8-replay (sk replay, A/A)                        -> nopcommerce-replay
  verdict (ubuntu-24.04) needs tools, iis: 9-verdict (sk compare, sk evidence) -> nopcommerce-evidence-pack

warmup-eshop.yml (P3)
  tools   (ubuntu-24.04)  chain-tools.yml: sk and Portcullis at their pinned commits -> chain-tools
  legacy  (windows-2022)  eshop 1-build -> IIS -> 2-deploy -> 3-record (sk capture)
                          -> 4-replay (sk replay, A/A) -> 5-scan (Portcullis) -> eshop-recordings
  verdict (ubuntu-24.04)  6-verdict: sk compare, sk gate, sk evidence       -> eshop-evidence-pack
\end{lstlisting}

\paragraph{\texttt{legacy-build.yml}.} Triggery: \texttt{pull\_request}, \texttt{push} na
\texttt{main} i~\texttt{workflow\_dispatch} z~wejściem \texttt{runner}
(\texttt{windows-2022} -- domyślny i~wspierany, \texttt{windows-2025} -- do sprawdzenia
przenośności). Domyślna powłoka to \texttt{powershell}, czyli Windows PowerShell 5.1:
skrypty muszą w~niej działać, więc CI to udowadnia. Joby: \texttt{lint} (ubuntu-24.04,
10~min) -- PSScriptAnalyzer i~kontrola ASCII (sekcja~\ref{sec:bench-higiena});
\texttt{build} (wybrany runner, 45~min) -- \texttt{microsoft/setup-msbuild},
\texttt{NuGet/setup-nuget} z~nuget.exe 7.9.0, bufor \path{downloads} z~kluczem
\texttt{legacy-build-downloads-<hash(pins.json)>} (każdy plik ponownie weryfikowany), kroki
1 i~2; \texttt{tools} -- \path{chain-tools.yml}; \texttt{iis} (needs: build, tools; 60~min)
-- pobranie artefaktów \texttt{legacy-site} i~\texttt{chain-tools}, .NET 10 przez
\texttt{actions/setup-dotnet}, bufor nośnika SQL Server
(\texttt{legacy-iis-downloads-<hash>}, ponownie weryfikowany skrótem i~podpisem), kroki
3--8, przy porażce \path{collect-diagnostics.ps1}; \texttt{verdict} (needs: tools, iis;
ubuntu-24.04, 15~min, \texttt{pwsh}) -- pobranie narzędzi i~przebiegu, krok~9
z~\texttt{-Pdf required}, bo obraz hostowany ma Google Chrome, więc pakiet musi mieć PDF.

\paragraph{\texttt{chain-tools.yml}.} Tylko \texttt{workflow\_call}. Job \texttt{build}
(ubuntu-24.04, 15~min, \texttt{DOTNET\_NOLOGO}, \texttt{DOTNET\_CLI\_TELEMETRY\_OPTOUT}):
odczyt pinów przez \texttt{jq}, checkout \repopath{sk} i~Portcullis na przypiętych
commitach, SDK z~\path{global.json} repozytorium \repopath{sk}, asercja, że \texttt{HEAD}
obu checkoutów równa się pinom, \texttt{dotnet build} obu projektów CLI w~Release do
\path{chain-tools/secondkey} i~\path{chain-tools/portcullis}, próba
\texttt{SecondKey.Cli.dll -{}-help}, plik \path{COMMITS} i~tabela w~podsumowaniu.

\paragraph{\texttt{warmup-eshop.yml}.} Triggery jak w~\path{legacy-build.yml} (bez
wejścia). Joby: \texttt{tools}; \texttt{legacy} (windows-2022, 45~min) -- narzędzia, .NET
10, MSBuild, bufor \texttt{warmup-eshop-downloads-<hash>}, krok~1 rozgrzewki, IIS przez
\path{3-install-iis-sql.ps1} z~\texttt{-SkipSql}, kroki 2--5; \texttt{verdict}
(ubuntu-24.04, 15~min) -- krok~6 z~\texttt{-Pdf required}.

\paragraph{\texttt{secret-scan.yml}.} Triggery \texttt{pull\_request} i~\texttt{push} na
\texttt{main}; job \texttt{gitleaks} (ubuntu-latest) pobiera całą historię
(\texttt{fetch-depth: 0} -- poświadczenie dodane i~usunięte później i~tak zostało
wypchnięte), instaluje gitleaks 8.28.0 zweryfikowany skrótem SHA-256 z~pliku sum wydania
(akcje są przypięte przez SHA, więc skaner -- przez skrót) i~uruchamia
\texttt{gitleaks detect -{}-redact -{}-config .gitleaks.toml -{}-verbose}.

{\footnotesize
\begin{xltabular}{\linewidth}{@{}>{\raggedright\arraybackslash}p{3.9cm} >{\raggedright\arraybackslash}p{2.6cm} >{\raggedright\arraybackslash}X l@{}}
\caption{Artefakty CI}\label{tab:bench-artefakty}\\
\toprule
Artefakt & Job (kiedy) & Zawartość (względem \texttt{<work>}) & Dni \\
\midrule
\endfirsthead
\toprule
Artefakt & Job (kiedy) & Zawartość (względem \texttt{<work>}) & Dni \\
\midrule
\endhead
\bottomrule
\endlastfoot
\path{legacy-site} & build (brak plików = błąd) & \path{legacy-site},
  \path{state\build.json}, \path{state\legacy-site.sha256} & 7 \\
\path{legacy-build-logs} & build (porażka) & \path{logs} (binlogi MSBuild) & 7 \\
\path{chain-tools} & chain-tools (brak = błąd) & \path{secondkey}, \path{portcullis},
  \path{COMMITS} & 7 \\
\path{nopcommerce-traffic} & iis (zawsze) & \path{traffic}, \path{state\traffic.json},
  \path{state\store-config.json}, \path{logs\traffic-capture.log*} & 30 \\
\path{nopcommerce-replay} & iis (zawsze) & \path{replay}, \path{state\replay.json},
  \path{state\traffic.json} & 30 \\
\path{legacy-iis-diagnostics} & iis (porażka) & \path{diagnostics} & 7 \\
\path{nopcommerce-evidence-pack} & verdict (zawsze) & \path{replay/evidence} & 30 \\
\path{eshop-recordings} & legacy (zawsze) & \path{eshop}, \path{state\eshop-*.json},
  \path{logs\eshop-capture.log*} & 14 \\
\path{eshop-build-logs} & legacy (porażka) & \path{logs\*.binlog} & 7 \\
\path{eshop-evidence-pack} & verdict (zawsze; brak = błąd) & \path{eshop/evidence} & 30 \\
\end{xltabular}
}

Artefakty są przesyłane względem katalogu roboczego, więc pobranie do innego katalogu
roboczego odtwarza \path{legacy-site}, \path{traffic}, \path{replay} i~\path{state} tak,
jak były. \textbf{Co jest bramką:} pull request przechodzi, gdy przejdą: lint (żadnego
znaleziska PSScriptAnalyzer, same znaki ASCII); build (pin, czyste drzewo, kontrole
witryny); kroki 3--6 (smoke test, konfiguracja z~odczytem); nagranie (każdy scenariusz
dotarł do celu, liczba wymian się zgadza, \texttt{sk validate}); odtworzenie (bez wyników
bez odpowiedzi i~nieudanych resetów, pełna liczba scenariuszy); werdykt A/A kontraktu
i~werdykt A/A rozgrzewki (\texttt{pass}, każda zaakceptowana klauzula spełniona, żadna
niewykonana); skan gitleaks. Wynik \texttt{sk gate} na drzewie eShopLegacyMVC jest
zapisywany, a~nie egzekwowany. Workflowy CI wszystkich repozytoriów frameworku opisuje
dodatek~\ref{app:ci}.

\zrodla{\path{letsgolegacy.secondkey-nopcommerce-bench/.github/workflows/legacy-build.yml},
\path{letsgolegacy.secondkey-nopcommerce-bench/.github/workflows/chain-tools.yml},
\path{letsgolegacy.secondkey-nopcommerce-bench/.github/workflows/warmup-eshop.yml},
\path{letsgolegacy.secondkey-nopcommerce-bench/.github/workflows/secret-scan.yml},
\path{letsgolegacy.secondkey-nopcommerce-bench/.github/dependabot.yml},
\path{letsgolegacy.secondkey-nopcommerce-bench/docs/TOPOLOGY.md}}

\subsection{Wyniki P3, P4 i~P5}\label{sec:bench-wyniki}

Liczby pochodzą z~przebiegów CI, które zamknęły zadania; pakiety i~nagrania z~tych
przebiegów są artefaktami przechowywanymi przez 30~dni, a~każdy kolejny przebieg
workflowu tworzy je od nowa.

\zrodla{\path{letsgolegacy.secondkey-nopcommerce-bench/results/p3/README.md},
\path{letsgolegacy.secondkey-nopcommerce-bench/results/p4/README.md},
\path{letsgolegacy.secondkey-nopcommerce-bench/results/p5/README.md}}

\subsubsection{P3: pierwszy pakiet rozgrzewki}\label{sec:bench-wyniki-p3}

Zrobione, gdy powstaje pakiet: przebieg \texttt{warmup-eshop} 36240136388 (2026-09-26,
commit \texttt{3510b98}), artefakt \texttt{eshop-evidence-pack}.

\begin{center}
\small
\begin{tabularx}{\linewidth}{@{}l X@{}}
\toprule
Aplikacja & eShopLegacyMVC z~\texttt{dotnet-architecture/eShopModernizing@63bc9ec}, dane
  mock, IIS na \texttt{windows-2022} \\
Narzędzia & \repopath{sk} na \texttt{967c49c}, Portcullis na \texttt{84e1925} \\
Nagrane & \textbf{7~scenariuszy, 27~wymian} przez \texttt{sk capture} (E01--E07) \\
Odtworzone & A/A: 54~wyniki, 14~resetów (restart puli przed każdym scenariuszem po każdej
  stronie), 0~bez odpowiedzi, 50,5~s \\
Werdykt & \textbf{pass} -- 27 \texttt{equal}, 0 \texttt{equal-under-contract},
  0 \texttt{regression}, 0 \texttt{fix-candidate} \\
Kontrakt & 13~klauzul, wszystkie zaakceptowane: \textbf{13 spełnionych}, 0~niewykonanych;
  udział nieobecności 23\% (3 \texttt{never}) \\
Bramka & Portcullis na kodzie legacy, całe drzewo: 27~znalezisk --
  \texttt{PORTCULLIS\_MIG\_SYSTEM\_WEB} 20 ostrzeżeń,
  \texttt{PORTCULLIS\_MIG\_HTTPCONTEXT\_CURRENT} 4 błędy,
  \texttt{PORTCULLIS\_MIG\_CONFIGURATION\_MANAGER} 3 błędy. \texttt{sk gate}:
  \texttt{fail}, 7 blokujących -- zapisane, nie egzekwowane: to idiomy, które migracja
  musi usunąć, a~żadna zmiana nie była bramkowana \\
Pakiet & \path{contract.yaml}, \path{index.html}, \path{report.pdf} (Google Chrome na
  runnerze linuksowym), \path{run.skrun}, \path{sarif/portcullis.sarif},
  \path{statement.intoto.json}, \path{verdict.json}, \path{manifest.json} -- każdy
  SHA-256 sprawdzony ponownie z~manifestem \\
\bottomrule
\end{tabularx}
\end{center}

Czasy w~tym przebiegu: build i~publikacja 52~s; IIS 22~s; wdrożenie 21~s (pierwsza
odpowiedź po 14~s); nagranie 28~s (około 3~s na restart); odtworzenie 50~s; skan 3~s;
werdykt i~pakiet 13~s; cały workflow 4,5~minuty. Czego rozgrzewka nauczyła przebiegi
nopCommerce: łańcuch działa od początku do końca na hoście Windows i~jobie linuksowym,
między którymi przechodzą tylko pliki; scenariusz jako sesja nazwana nagłówkiem i~proxy
zatrzymywane po liczbie wymian przeszły do P4 bez zmian; \textbf{operator wartości na
ekstraktorze z~\texttt{all: true} wymaga \texttt{[*]}} -- \texttt{gt} na
\texttt{extract.catalogPrices} porównuje samą tablicę i~nigdy nie jest spełniony, więc
klauzula wybiera \texttt{extract.catalogPrices[*]} (wyłapane przed CI, na lokalnym
zastępniku aplikacji; kontrakt P5 stosuje tę samą regułę).

\zrodla{\path{letsgolegacy.secondkey-nopcommerce-bench/results/p3/README.md}}

\subsubsection{P4: zestaw ruchu nopCommerce}\label{sec:bench-wyniki-p4}

Zrobione, gdy liczba scenariuszy jest zapisana: \finding{38~scenariuszy i~290~wymian},
nagranych przez \texttt{sk capture} przed sklepem z~P2 w~przebiegu \texttt{legacy-build}
36253613672 (2026-09-26, commit \texttt{45fb761}); nagranie to artefakt
\texttt{nopcommerce-traffic} tego przebiegu. Każdy przebieg \texttt{legacy-build}
nagrywa zestaw od nowa i~zapisuje liczby scenariusz po scenariuszu w~podsumowaniu joba.

\begin{center}
\small
\begin{tabularx}{\linewidth}{@{}l X@{}}
\toprule
Aplikacja & nopCommerce 3.90 (\texttt{12d1f01}) z~danymi przykładowymi, IIS na
  \texttt{windows-2022}, SQL Server 2022 Express \\
Sklep & skonfigurowany według sekcji~2 planu przez panel (krok~6), każda zmiana odczytana:
  stawki podatku według kraju, stanu i~kodu, stałe stawki wysyłki i~darmowa wysyłka powyżej
  \$1,000.00, siedem kuponów, zmiany katalogu; cztery zadania harmonogramu wyłączone,
  aplikacja zrestartowana -- z~powrotem po 14~s, z~nowym wpisem „Application started” \\
Narzędzie & \texttt{sk capture} na \texttt{967c49c} na \texttt{127.0.0.1:8000}; jedna
  sesja na scenariusz, nazwana nagłówkiem \texttt{x-bench-scenario}; \texttt{sk validate}
  akceptuje \path{.skcap} \\
Nagrane & \textbf{38~scenariuszy, 290~wymian} w~166~s -- 30~zaplanowanych scenariuszy
  (T01--T30), niektóre jako kilka sesji \\
Stan startowy & każdy scenariusz ze snapshotu \texttt{nopcommerce\_legacy\_snapshot}
  skonfigurowanego sklepu: 38~przywróceń po 3,3--3,4~s; pierwsze żądanie po każdym dostało
  od razu 200 \\
\bottomrule
\end{tabularx}
\end{center}

\begin{center}
\small
\begin{tabularx}{\linewidth}{@{}l r r X@{}}
\toprule
Obszar & Scenariusze & Wymiany & \\
\midrule
Katalog & 5 & 10 & T01--T05 \\
Wyszukiwanie & 1 & 5 & T06 \\
Koszyk & 4 & 19 & T07--T10 \\
Rabaty & 12 & 80 & T11, T12, T12B, T13--T17, T18A, T18B, T19, T20 \\
Podatki & 8 & 98 & T21, T22A--T22F, T23 -- każdy to checkout na adres rozliczeniowy \\
Wysyłka & 3 & 27 & T24--T26 \\
Checkout & 2 & 36 & T27 (gość), T28 (zarejestrowany: odrzucona karta, potem zamówienie) \\
Zamówienia & 3 & 15 & T29, T30A, T30B \\
\midrule
\textbf{Razem} & \textbf{38} & \textbf{290} & \\
\bottomrule
\end{tabularx}
\end{center}

\zrodla{\path{letsgolegacy.secondkey-nopcommerce-bench/results/p4/README.md}}

\subsubsection{P5: werdykt A/A kontraktu}\label{sec:bench-wyniki-p5}

Zrobione, gdy kontrakt przechodzi walidację i~jest spełniony na systemie legacy.
\texttt{sk validate} akceptuje \path{contract.yaml} i~\path{secondkey.yaml}, a~odtworzenie
A/A w~przebiegu \texttt{legacy-build} 36255313866 (2026-09-26, commit \texttt{c6732ab})
daje \texttt{sk compare} wynik \textbf{pass: 76 z~76 zaakceptowanych klauzul spełnionych,
0~niewykonanych}. Pakiet to artefakt \texttt{nopcommerce-evidence-pack}. Każdy przebieg
\texttt{legacy-build} nagrywa, odtwarza, porównuje i~pakuje od nowa i~zawodzi, jeśli werdykt
A/A nie jest taki sam.

\begin{center}
\small
\begin{tabularx}{\linewidth}{@{}l X@{}}
\toprule
Nagranie & zestaw P4 nagrany w~tym samym przebiegu: 38~scenariuszy, 290~wymian \\
Narzędzia & \repopath{sk} na \texttt{967c49c}; odtworzenie na hoście Windows obok IIS,
  porównanie i~pakiet w~Linuksie \\
Odtworzone & A/A, legacy i~kandydat to \url{http://localhost:8080}: 580~wyników
  (200~×~500, 302~×~78, 404~×~2), 0~bez odpowiedzi; 76~resetów
  \texttt{sqlServerSnapshot}, 0~nieudanych, średnio 3,3~s; łącznie 300~s \\
Werdykt & \textbf{pass} -- 290~wymian w~38~scenariuszach: 272 \texttt{equal},
  18 \texttt{equal-under-contract}, 0 \texttt{regression}, 0 \texttt{fix-candidate} \\
Kontrakt & 76~klauzul, wszystkie zaakceptowane: \textbf{76 spełnionych}, 0~niespełnionych,
  0~niewykonanych; 15 \texttt{never} -- udział nieobecności 19,7\%; 32~ekstraktory \\
Normalizacja & 18 wymian \texttt{equal-under-contract} według masek, które usunęły ich
  różnice (wymiana może mieć dwie): \texttt{anti-forgery-token} 14 (podsumowanie zamówienia
  kroku potwierdzenia), \texttt{cart-line-picture} 17 (te same 14 i~3 minikoszyki),
  \texttt{pdf-bytes} 1 (faktura) \\
Pakiet & \path{contract.yaml}, \path{index.html}, \path{report.pdf} (Google Chrome na
  runnerze linuksowym), \path{run.skrun}, \path{statement.intoto.json},
  \path{verdict.json}, \path{manifest.json} -- każdy SHA-256 sprawdzony ponownie \\
\bottomrule
\end{tabularx}
\end{center}

\zrodla{\path{letsgolegacy.secondkey-nopcommerce-bench/results/p5/README.md}}

\subsection{Runbook P6: przebieg modernize-dotnet (praca człowieka)}\label{sec:bench-p6}

P6 wykonuje człowiek: osoba z~subskrypcją GitHub Copilot uruchamia agenta migracyjnego
Microsoftu (GitHub Copilot app modernization for .NET, „modernize-dotnet”) na przypiętym
nopCommerce i~przekazuje wynik benchowi. Runbook \repopath{docs/P6-RUNBOOK.md} czyni to
mechanicznym: gdzie trafia kod, jakie standardy dostaje agent, co wpisać, co zapisać i~czego
potrzebują P7 i~P8. Kryterium ukończenia (WORKPLAN): kandydat się buduje, a~istniejące
testy przechodzą. Wejście: nopCommerce \texttt{release-3.90} na
\texttt{12d1f0139149a9d6e84964d325d62bae6ac748f6} oraz skill standardów z~zadania R4
w~\path{letsgolegacy.secondkey-standards} (sekcja~\ref{sec:standardy}). Wyjście: pull request
w~forku nopCommerce, którego commit czołowy jest kandydatem, zapis przebiegu i~godziny
w~dzienniku P10. Runbook napisano przed pierwszym przebiegiem P6; zapis tego przebiegu ma się
stać przykładem.

\begin{enumerate}
  \item \textbf{Przed startem.} Potrzebne: GitHub Copilot z~trybem agenta (Pro, Business
    albo Enterprise); jeden klient z~agentem modernizacji -- Visual Studio 2026 (albo 2022
    17.14+) z~\emph{GitHub Copilot app modernization}, VS Code z~rozszerzeniem o~tej nazwie
    albo GitHub Copilot CLI (polecenia runbooka są te same, różni się tylko krok~5.2);
    Windows z~MSBuild Visual Studio dla .NET Framework, .NET 10 SDK, git i~GitHub CLI
    (\texttt{gh}) -- bazowy build legacy idzie przez MSBuild, kandydat przez
    \texttt{dotnet}; prawo zapisu do repozytorium benchu (dla zapisu przebiegu);
    \textbf{zrobione R4}: otagowane wydanie repozytorium standardów z~\path{generated/skills/}
    -- bez niego P6 się nie zaczyna. Uruchom stoper i~zapisz czas startu: wszystko do
    otwarcia pull requesta liczy się jako godziny P6; czas CI benchu -- nie.
  \item \textbf{Fork i~gałęzie.} Kandydat to zmodyfikowany nopCommerce, więc żyje tam, gdzie
    obowiązuje już licencja nopCommerce -- w~forku upstream -- i~nigdy w~repozytorium
    benchu.
\begin{lstlisting}
# 2.1 Fork upstream once, into the account that owns the bench (public, as upstream is).
gh repo fork nopSolutions/nopCommerce --clone=false

# 2.2 Take the pinned commit from upstream, shallow, with upstream's exact bytes.
git init nop-p6 && cd nop-p6
git config core.autocrlf false
git remote add upstream https://github.com/nopSolutions/nopCommerce.git
git remote add fork https://github.com/<your-account>/nopCommerce.git
git fetch --depth 1 upstream tag release-3.90
git rev-parse "release-3.90^{commit}"          # must print 12d1f0139149a9d6e84964d325d62bae6ac748f6
git checkout -b secondkey/p6-base 12d1f0139149a9d6e84964d325d62bae6ac748f6
\end{lstlisting}
    Jeśli \texttt{rev-parse} wypisze cokolwiek innego, zatrzymaj się: tag przesunął się
    w~upstream (krok~1 benchu też by go odrzucił). Gałąź \texttt{secondkey/p6-base} to
    przypięty commit plus jeden commit ze skillem standardów i~nigdy się nie zmienia;
    \texttt{secondkey/p6-candidate} powstaje z~niej i~zawiera wszystko, co zmienią agent
    i~człowiek, tylko w~trakcie P6. Pull request idzie z~kandydata do bazy, więc jego diff
    to dokładnie migracja.
  \item \textbf{Standardy dla agenta.}
\begin{lstlisting}
# 3.1 Get the standards at the R4 release tag (write the tag into the record).
git clone --depth 1 --branch <standards-tag> https://github.com/konradcinkusz/letsgolegacy.secondkey-standards.git ../standards

# 3.2 Copy every generated skill into the target's .github/skills/.
mkdir -p .github/skills
cp -R ../standards/generated/skills/. .github/skills/
git -C ../standards rev-parse HEAD             # the standards commit, for the record

# 3.3 Commit it alone, on the base branch.
git add .github/skills
git commit -m "Add Second Key standards skills (<standards-tag>, <standards-commit>)"
git push -u fork secondkey/p6-base
git checkout -b secondkey/p6-candidate
git push -u fork secondkey/p6-candidate
\end{lstlisting}
    Każdy skill to katalog z~\path{SKILL.md}; Copilot odkrywa skille w~\path{.github/skills/}
    otwartego repozytorium. Nie edytuj ich tutaj -- zmiana standardu to zmiana R4.
  \item \textbf{Bazowa linia testów legacy.} „Istniejące testy przechodzą” wymaga linii
    bazowej: które z~własnych testów nopCommerce przechodzą na niezmienionym kodzie. 3.90
    ma pięć projektów NUnit~3 w~\path{src/Tests} (\texttt{Nop.Core.Tests},
    \texttt{Nop.Data.Tests}, \texttt{Nop.Services.Tests}, \texttt{Nop.Web.MVC.Tests},
    \texttt{Nop.Tests}). Drzewo legacy to to samo, które buduje bench, więc z~klonu benchu:
\begin{lstlisting}
.\scripts\1-fetch-nopcommerce.ps1     # C:\nopbench\legacy-src, verified at 12d1f01
.\scripts\2-build-legacy.ps1          # builds every project, the test projects included
nuget install NUnit.ConsoleRunner -Version 3.22.0 -OutputDirectory C:\nopbench\tools
$runner = 'C:\nopbench\tools\NUnit.ConsoleRunner.3.22.0\tools\nunit3-console.exe'
& $runner (Get-ChildItem C:\nopbench\legacy-src\src\Tests\*\bin\Release\*.Tests.dll).FullName --result=C:\nopbench\p6-legacy-tests.xml
\end{lstlisting}
    Zapisz sumę, zaliczone, niezaliczone i~pominięte dla każdej asemblacji. Test, który
    zawodzi na legacy, nie jest regresją, gdy zawodzi u~kandydata; test zaliczony na legacy
    i~niezaliczony u~kandydata -- jest.
  \item \textbf{Uruchomienie agenta.} (5.1) Otwórz \texttt{nop-p6} na gałęzi
    \texttt{secondkey/p6-candidate} w~wybranym kliencie, zalogowany do GitHub Copilot,
    i~przed czymkolwiek innym zapytaj: „List the skills you have loaded from this
    repository.” -- odpowiedź wklej do zapisu. Jeśli brakuje skilli Second Key, zatrzymaj
    się i~popraw krok~3 (akceptacja R4 to właśnie to: agent podejmuje skill). (5.2) Uruchom
    modernizację: w~Visual Studio otwórz \path{src\NopCommerce.sln} i~uruchom
    \emph{GitHub Copilot app modernization} (menu kontekstowe rozwiązania albo Copilot Chat
    w~trybie agenta z~wybranym agentem modernizacji); w~VS Code otwórz katalog, Copilot
    Chat w~trybie \textbf{Agent} z~agentem app modernization; w~Copilot CLI --
    \texttt{copilot} w~korzeniu repozytorium z~włączonym agentem. Nazwy menu zmieniają się
    między wydaniami: \emph{gdzie} kliknąć mówi dokumentacja narzędzia, \emph{co} mu dać --
    runbook. Wyślij dosłownie:
\begin{lstlisting}
Modernize the solution in src/NopCommerce.sln from .NET Framework 4.5.1 to .NET 10.
Migrate the ASP.NET MVC 5 web application (Nop.Web, the Administration area and the
plugins) to ASP.NET Core on .NET 10. Keep the SQL Server database schema, the public
URLs and the observable behaviour of the storefront unchanged. Keep the existing unit
test projects and make them build and run on .NET 10. Follow the standards in
.github/skills. Commit your work on the current branch in logical steps.
\end{lstlisting}
    Każdy kolejny prompt, każde pytanie agenta i~każdy wybór trafiają do zapisu dosłownie,
    po kolei. Agent commituje sam; nie rób squasha. (5.3) Jeśli musisz zmienić kod sam
    (agent utknął albo prosi), zrób osobny commit z~komunikatem zaczynającym się od
    \texttt{human:} i~uzasadnieniem -- pakiet dowodowy odróżnia zmiany agenta od zmian
    człowieka po tym prefiksie.
  \item \textbf{Build i~testy kandydata.}
\begin{lstlisting}
dotnet --version                                   # record it
dotnet build src\NopCommerce.sln -c Release        # must succeed
dotnet test  src\NopCommerce.sln -c Release --logger "trx;LogFileName=p6-candidate-tests.trx"
\end{lstlisting}
    P6 jest zrobione, gdy build się udaje, a~każdy test zaliczony w~linii bazowej
    przechodzi. Zapisz liczby obok bazowych. Test zaliczony na legacy, który agent usunął,
    pominął albo opróżnił, nie liczy się jako zaliczony -- napisz o~tym w~zapisie.
  \item \textbf{Pull request kandydata.}
\begin{lstlisting}
git push fork secondkey/p6-candidate
gh pr create --repo <your-account>/nopCommerce --base secondkey/p6-base --head secondkey/p6-candidate \
  --title "P6: modernize-dotnet candidate - nopCommerce 3.90 to .NET 10" \
  --body-file p6-record.md
git rev-parse HEAD                                  # the candidate commit
\end{lstlisting}
    Nie scalaj go. Jego commit czołowy jest kandydatem; późniejsze commity tworzą nowego
    kandydata.
  \item \textbf{Zapis przebiegu.} Napisz \path{p6-record.md} (treść pull requesta), dodaj
    ten sam plik do benchu jako \path{results/p6/run-record.md} w~pull requeście i~wpisz
    godziny do dziennika P10. Szablon:
\begin{lstlisting}
# P6 run record

| | |
|---|---|
| Operator | <name> |
| Started / finished (UTC) | <yyyy-mm-dd hh:mm> / <yyyy-mm-dd hh:mm> |
| Hours (hands-on / waiting on the agent) | <h> / <h> |
| Client and version | <Visual Studio 2026 18.x / VS Code x.y + extension a.b / Copilot CLI x.y> |
| Modernization agent version | <as shown by the client> |
| Model | <model selected in Copilot> |
| .NET SDK | <dotnet --version> |
| Legacy commit | 12d1f0139149a9d6e84964d325d62bae6ac748f6 |
| Standards | <standards-tag> at <standards-commit> |
| Skills the agent listed | <paste> |
| Candidate PR / head commit | <url> / <sha> |
| Legacy tests (total / passed / failed / skipped) | <n / n / n / n> |
| Candidate tests (total / passed / failed / skipped) | <n / n / n / n> |
| Tests passing on legacy but not on the candidate | <none, or list> |
| Human commits | <none, or list with reasons> |

## Prompts and answers

1. <prompt, verbatim>
   - Agent: <summary of the answer, or the question it asked>
   - Choice: <what you answered>

## Notes

<anything surprising: files the agent deleted, packages it replaced, warnings it ignored>
\end{lstlisting}
  \item \textbf{Przekazanie do P7 i~P8.} P7 bierze commit czołowy PR kandydata, przypina go
    w~benchu, buduje, wdraża na hoście Windows obok legacy (\url{http://localhost:8090/},
    własna baza przywrócona z~tego samego snapshotu), odtwarza na obu nagranie P4
    i~porównuje: \path{verdict.json}. P8 uruchamia reguły migracji Portcullis na liniach
    zmienionych przez PR: SARIF dołączony do pakietu. P9 publikuje pakiet, który cytuje PR,
    jego commit czołowy, wersję standardów i~ten zapis. Nowy push do
    \texttt{secondkey/p6-candidate} po starcie P7 to nowy kandydat: P7 i~P8 biegną ponownie
    na nowym commicie, a~zapis to odnotowuje.
\end{enumerate}

\begin{center}
\small
\begin{tabularx}{\linewidth}{@{}>{\raggedright\arraybackslash}X >{\raggedright\arraybackslash}X@{}}
\toprule
Objaw & Przyczyna \\
\midrule
Agent nie wspomina standardów Second Key & \path{.github/skills/} nie ma w~otwartym
  katalogu albo wersja klienta nie ładuje skilli repozytorium: krok~5.1 wyłapuje to przed
  jakąkolwiek zmianą kodu \\
Diff PR ma tysiące niezwiązanych linii & gałęzie nie powstały z~przypiętego commita albo
  przekonwertowano końce linii (\texttt{core.autocrlf}) \\
„Testy przechodzą”, ale uruchomiono ich mniej & agent wykluczył, pominął albo usunął
  projekty testów; porównaj liczby z~linią bazową \\
P7 nie odtwarza buildu & kandydat zależy od czegoś, co jest tylko na maszynie operatora
  (lokalne źródło NuGet, SDK w~wersji preview); zapisz SDK i~każde źródło \\
Kandydat zmienił się w~trakcie P7 & ktoś wypchnął zmiany na gałąź kandydata: P7 przypina
  commit, nigdy gałąź \\
\texttt{git push} odmawia z~„shallow update not allowed” & fork utworzono z~„copy the
  default branch only”, więc brakuje historii przypiętego commita: pobierz tag bez
  \texttt{-{}-depth} albo utwórz fork ponownie z~ustawieniami domyślnymi \texttt{gh repo
  fork} \\
\bottomrule
\end{tabularx}
\end{center}

\begin{checklista}
  \item Fork utworzony; \texttt{secondkey/p6-base} to przypięty commit plus commit skilli,
    nic więcej.
  \item Skille skopiowane z~otagowanego wydania standardów; tag i~commit są w~zapisie.
  \item Linia bazowa testów legacy zapisana (liczby dla każdej asemblacji).
  \item Agent wymienił skille Second Key, zanim cokolwiek zmienił.
  \item Każdy prompt, pytanie i~wybór zapisane dosłownie; zmiany człowieka to osobne
    commity \texttt{human:}.
  \item \texttt{dotnet build} się udaje; każdy test zaliczony na legacy przechodzi
    u~kandydata.
  \item Pull request otwarty do \texttt{secondkey/p6-base}, nie scalony; jego commit
    czołowy zapisany.
  \item Zapis zacommitowany do benchu jako \path{results/p6/run-record.md}; godziny dodane
    do dziennika P10.
\end{checklista}

\zrodla{\path{letsgolegacy.secondkey-nopcommerce-bench/docs/P6-RUNBOOK.md},
\path{letsgolegacy.secondkey-nopcommerce-bench/docs/WORKPLAN.md}}

\subsection{Status i~plany P7--P10}\label{sec:bench-plany}

\begin{center}
\small
\begin{tabularx}{\linewidth}{@{}l >{\raggedright\arraybackslash}X >{\raggedright\arraybackslash}X l@{}}
\toprule
Id & Produkt & Zrobione, gdy & Status \\
\midrule
P1 & repozytorium, przypięty nopCommerce 3.x, build na runnerze Windows & CI go buduje &
  zrobione (\#1) \\
P2 & legacy pod IIS (topologia demo) z~SQL Server i~danymi & sklep odpowiada pod URL-em &
  zrobione (\#2) \\
P3 & rozgrzewka łańcucha na eShopModernizing & powstaje pakiet; po S15 rdzenia &
  zrobione (\#4) \\
P4 & zestaw ruchu nopCommerce nagrany do \path{*.skcap} & liczba scenariuszy zapisana &
  zrobione (\#5) \\
P5 & kontrakt nopCommerce: reguły „must” i~„never”, około 20\% asercji nieobecności &
  przechodzi walidację i~jest spełniony na legacy; po S8 rdzenia & zrobione (\#6) \\
P6 & przebieg modernize-dotnet ze skillem standardów z~R4 & kandydat się buduje,
  istniejące testy przechodzą -- \textbf{praca człowieka, wymaga GitHub Copilot} &
  planowane \\
P7 & odtworzenie i~porównanie legacy z~kandydatem; przegląd różnic & istnieje
  \path{verdict.json}, a~różnica NLS $\rightarrow$ ICU jest znaleziona albo jej brak
  udokumentowany & planowane \\
P8 & reguły migracji Portcullis na PR kandydata (R2, R3) & SARIF jest dołączony do
  pakietu & planowane \\
P9 & publikacja pakietu, kontraktu i~liczb (scenariusze, różnice, regresje, zabite
  mutanty) & publiczne & planowane \\
P10 & dziennik godzin wdrożenia dla P4--P7 & godziny są w~opublikowanym raporcie &
  planowane \\
\bottomrule
\end{tabularx}
\end{center}

P9 tego repozytorium jest bramką fazy 01 (sekcja~\ref{sec:status}). Plan P4, runbook P6
i~dziennik P10 przygotowano przed zadaniami, które ich używają (\#3).

\paragraph{P7: legacy kontra kandydat.} Mechanizm jest gotowy: krok~8 z~\texttt{-Candidate}
odtwarza to samo nagranie na kandydacie z~tym samym resetem i~kontraktem, a~krok~9
w~trybie „legacy vs candidate” zapisuje werdykt i~pakiet bez bramki A/A. Kandydat ma
działać na \url{http://localhost:8090/} z~własną kopią bazy; jego tożsamość dostaje własny
login SQL (ADR 0003). Sondy NLS $\rightarrow$ ICU (sekcja~6 planu) to hipotezy do
potwierdzenia albo obalenia, nie znaleziska:

\begin{center}
\small
\begin{tabularx}{\linewidth}{@{}>{\raggedright\arraybackslash}p{2.8cm} >{\raggedright\arraybackslash}X >{\raggedright\arraybackslash}p{1.9cm}@{}}
\toprule
Sonda & Ścieżka w~kodzie i~dlaczego może się różnić & Gdzie \\
\midrule
\textbf{Ujemne kwoty walutowe} w~sumach koszyka & \texttt{ShoppingCartModelFactory}
  formatuje \texttt{-discount} przez \texttt{PriceFormatter.GetCurrencyString}
  (\texttt{ToString("C", new CultureInfo("en-US"))}); NLS zapisuje ujemną kwotę
  \texttt{en-US} jako \texttt{(\$100.00)}, dane ICU -- jako \texttt{-\$100.00}.
  Najmocniejszy kandydat & T11, T14, T15 -- R27 \\
Wielkość liter kuponu & \texttt{CustomerExtensions.ApplyDiscountCouponCode} zapisuje
  \texttt{couponCode.Trim().ToLower()} w~\emph{bieżącej} kulturze, a~\texttt{ValidateDiscount}
  porównuje z~\texttt{InvariantCultureIgnoreCase}; wielkość liter i~porównanie znaków spoza
  ASCII może się różnić (P4 może dodać kupony z~\texttt{ß}, \texttt{İ}) & T12 -- R17 \\
Przycinanie kuponu & \texttt{String.Trim()} idzie za danymi Unicode środowiska: U+180E był
  białym znakiem w~starszych wersjach Unicode, w~nowszych nie jest & T12 \\
Porównanie kodu pocztowego & \texttt{GetTaxRate} z~\texttt{InvariantCultureIgnoreCase};
  różnić mogą się tylko litery, więc sondą są kody kanadyjskie & T23 -- R29 \\
Sortowanie po nazwie & \texttt{ProductLoadAllPaged} sortuje w~SQL; decyduje kolacja bazy,
  więc różnica oznacza, że kandydat przeniósł sortowanie do .NET & T02 -- R03 \\
Daty i~godziny & daty zamówień na stronach szczegółów i~historii; nowsze dane ICU wstawiają
  wąską twardą spację przed AM/PM w~\texttt{en-US} & T29 \\
\bottomrule
\end{tabularx}
\end{center}

\paragraph{P8, P9.} P8 uruchomi reguły migracji Portcullis (zadania R2, R3) na liniach,
które zmienia pull request kandydata, i~dołączy SARIF do pakietu
(sekcja~\ref{sec:portcullis}). P9 opublikuje pakiet, kontrakt i~liczby: scenariusze,
różnice, regresje i~zabite mutanty; pakiet cytuje PR, jego commit czołowy, wersję
standardów i~zapis P6.

\paragraph{P10: dziennik godzin.} Liczba godzin potrzebnych do objęcia prawdziwego systemu
legacy przez Second Key jest jedną z~liczb raportu fazy 01. Zakres: P4--P7 to wdrożenie
we właściwym sensie; P1 i~P2 są zapisywane osobno jako jednorazowy nakład pracy na uruchomienie strony
legacy. Zasady: jeden wiersz na sesję pracy na jednym zadaniu i~wykonawcę; godziny zegarowe
faktycznie spędzone, dziesiętnie (0,25 = 15~minut), a~czekanie na CI albo agenta tylko wtedy,
gdy nic innego się nie działo -- jako osobny wiersz „waiting”; wykonawca to osoba albo
agent (\texttt{agent: <model> (<tool>)}), nigdy oba, a~godziny agenta i~człowieka nie są
sumowane; każdy wiersz ma dowód (PR, przebieg, zapis); szacunki są oznaczone. Stan
dziennika w~migawce: siedem wierszy agenta z~jednej sesji około dwóch godzin (P1 0,40 i~0,35;
P2 0,30 i~0,45; przygotowanie P4 0,40; przygotowanie P6 0,15; P10 0,05; podział między
wiersze to szacunek z~czasów commitów i~przebiegów), zero godzin człowieka. Suma: P1--P2
-- 1,50~h agenta, przygotowania P4 i~P6 -- 0,55~h, P10 -- 0,05~h, wdrożenie P4--P7 --
0. Minuty runnerów są prowadzone osobno: \#1 -- 3~przebiegi (jeden nieudany), 4,5~min
Windows i~1,2~min Linux; \#2 -- 3~przebiegi (jeden nieudany), 28,3 i~1,5~min; \#3 -- około
7 i~0,5~min na przebieg. Zielony przebieg całego workflowu legacy trwał wtedy około
7~minut (build 1,5, IIS z~SQL Server, instalatorem i~smoke testem 5--5,5, lint 0,5
równolegle); nieudany przebieg P2 trwał 13,4~minuty, z~czego 11 w~instalatorze SQL Server.

\zrodla{\path{letsgolegacy.secondkey-nopcommerce-bench/docs/WORKPLAN.md},
\path{letsgolegacy.secondkey-nopcommerce-bench/docs/P4-TRAFFIC-PLAN.md},
\path{letsgolegacy.secondkey-nopcommerce-bench/docs/P6-RUNBOOK.md},
\path{letsgolegacy.secondkey-nopcommerce-bench/docs/P10-ONBOARDING-HOURS.md},
\path{letsgolegacy.secondkey-nopcommerce-bench/docs/TOPOLOGY.md},
\path{letsgolegacy.secondkey-nopcommerce-bench/docs/adr/0006-nopcommerce-contract-and-aa-proof.md}}

\subsection{Higiena repozytorium: lint, skan sekretów, polityka nagrań}\label{sec:bench-higiena}

\paragraph{Lint.} \path{PSScriptAnalyzerSettings.psd1} ustawia \texttt{Severity = Error,
Warning}, wyłącza \texttt{PSAvoidUsingWriteHost} (skrypty są interaktywnymi runbookami,
których wyjściem jest log dla człowieka, a~nie dane potoku) i~włącza
\texttt{PSUseCompatibleSyntax} dla wersji 5.1 i~7.0 oraz \texttt{PSUseCompatibleCommands}
z~profilami Windows PowerShell 5.1 i~PowerShell 7.0. Job \texttt{lint} uruchamia
\texttt{Invoke-ScriptAnalyzer -Path scripts -Recurse -Settings
./PSScriptAnalyzerSettings.psd1} i~zawodzi przy jakimkolwiek znalezisku, a~osobny krok
zawodzi przy znaku spoza ASCII w~\path{scripts/}. Wyjątki są w~kodzie, jako
\texttt{SuppressMessageAttribute} z~uzasadnieniem (np.
\texttt{PSUseShouldProcessForStateChangingFunctions} dla funkcji tworzących obiekty
w~pamięci, \texttt{PSAvoidUsingPlainTextForPassword} dla logowania klientów sklepu
przykładowego).

\paragraph{Skan sekretów.} Każdy klon włącza raz, przed pierwszym commitem, hook
\texttt{pre-commit} (\texttt{git config core.hooksPath .githooks}; Git for Windows
uruchamia go w~swoim bashu). Hook wymaga gitleaks i~bez niego odmawia commitu -- CI i~tak
skanuje -- a~uruchamia \texttt{gitleaks protect -{}-staged -{}-redact -{}-config
.gitleaks.toml}. Po trafieniu: jeśli to poświadczenie, zdejmij je ze stage i~czytaj ze
środowiska; jeśli nie (dane przykładowe, które muszą wyglądać jak sekret), dodaj ścieżkę
albo kształt do allowlisty \path{.gitleaks.toml} w~tym samym commicie, żeby wyjątek przeszedł
review. Ten sam skan biegnie w~CI po całej historii (\path{secret-scan.yml}).
\path{.gitleaks.toml} rozszerza reguły domyślne (\texttt{useDefault = true}) o~regułę
\texttt{sqlserver-connection-string-with-password}: connection string SQL Server
z~hasłem w~treści (sekretem jest samo hasło, grupa~1). Bench łączy się przez
\texttt{Integrated Security} (ADR 0003), więc hasło w~treści nigdy nie jest tu uprawnione;
allowlista przepuszcza tylko wartość zaczynającą się od \texttt{<} -- symbol zastępczy
mówiący, skąd pochodzi hasło. Sekrety trafiają do skryptów wyłącznie przez zmienne
środowiskowe.

\paragraph{Polityka nagrań.} Nagranie (\path{*.skcap}, \path{*.skrun}) czegokolwiek poza
własnym sklepem przykładowym benchu nigdy nie jest commitowane. \path{.gitignore} ignoruje
\path{*.skcap} i~\path{*.skrun} (poza \path{samples/**}, \path{schemas/samples/**}
i~\path{tests/**/fixtures/**}), a~także \path{.secondkey/}, \path{evidence-out/}
i~lokalne pliki \path{.env}. Nagrania opuszczają CI wyłącznie jako artefakty
(tabela~\ref{tab:bench-artefakty}). Żądania konfiguracyjne administratora nigdy nie
przechodzą przez proxy nagrywające, więc \path{*.skcap} zawiera tylko dane sklepu
przykładowego (opublikowane hasło klienta przykładowego, hasło R1, testowe numery kart),
a~ciasteczka są domyślnie redagowane przez \texttt{sk capture}. Zasady bezpieczeństwa
całego frameworku opisuje sekcja~\ref{sec:bezpieczenstwo}.

\paragraph{Własność i~zależności.} \path{.github/CODEOWNERS} przypisuje każdą ścieżkę
właścicielowi repozytorium i~jawnie wymienia te, które decydują o~tym, co bench dowodzi:
\path{/.github/}, \path{/.githooks/}, \path{/.gitleaks.toml}, \path{/contract/}
i~\path{/pins.json}. \path{.github/dependabot.yml} aktualizuje akcje GitHub raz
w~tygodniu, jedną grupą.

\zrodla{\path{letsgolegacy.secondkey-nopcommerce-bench/PSScriptAnalyzerSettings.psd1},
\path{letsgolegacy.secondkey-nopcommerce-bench/.github/workflows/legacy-build.yml},
\path{letsgolegacy.secondkey-nopcommerce-bench/.githooks/pre-commit},
\path{letsgolegacy.secondkey-nopcommerce-bench/.gitleaks.toml},
\path{letsgolegacy.secondkey-nopcommerce-bench/.github/workflows/secret-scan.yml},
\path{letsgolegacy.secondkey-nopcommerce-bench/.gitignore},
\path{letsgolegacy.secondkey-nopcommerce-bench/README.md},
\path{letsgolegacy.secondkey-nopcommerce-bench/.github/CODEOWNERS},
\path{letsgolegacy.secondkey-nopcommerce-bench/.github/dependabot.yml},
\path{letsgolegacy.secondkey-nopcommerce-bench/scripts/lib/NopBench.Shop.ps1}}
